\documentclass[11pt,a4paper]{article}
\usepackage[margin=1in]{geometry}
\usepackage{amsmath,amssymb}
\usepackage{enumitem}
\usepackage{multicol}
\usepackage{array}
\usepackage{booktabs}
\usepackage{titlesec}
\usepackage{fancyhdr}

\pagestyle{fancy}
\fancyhf{}
\fancyhead[C]{The Hope Science Academy -- Statistics 11th (Combined)}
\fancyfoot[C]{\thepage}
\setlength{\headheight}{14pt}

\titleformat{\section}{\large\bfseries}{}{0em}{}

\newlist{mcqlist}{enumerate}{1}
\setlist[mcqlist,1]{label=\textbf{\arabic*.}, leftmargin=2.2em, itemsep=8pt, topsep=4pt}
\newlist{shortlist}{enumerate}{1}
\setlist[shortlist,1]{label=\textbf{\arabic*.}, leftmargin=2.5em, itemsep=8pt, topsep=4pt}
\newlist{longlist}{enumerate}{1}
\setlist[longlist,1]{label=\textbf{\arabic*.}, leftmargin=*, itemsep=10pt, topsep=6pt}

\title{\textbf{The Hope Science Academy}\\[2pt]\large Statistics -- Class 11th\\[2pt]\normalsize Combined Question Bank (Full Book)}
\author{}
\date{}

\begin{document}
\maketitle
\tableofcontents
\newpage

\section{Multiple Choice Questions (1--440)}

\section*{Choose the correct answer.}

\begin{mcqlist}

\item In the plural sense, statistics mean:
\begin{quote}\begin{multicols}{2}
\noindent (A) Methods \\
(B) Numerical Data \\
(C) Sample values \\
(D) Population values
\end{multicols}\end{quote}

\item In the singular sense, statistics means:
\begin{quote}\begin{multicols}{2}
\noindent (A) Methodology \\
(B) Numerical data \\
(C) Sample values \\
(D) Count data
\end{multicols}\end{quote}

\item Major area of statistics today is concerned with drawing conclusions from:
\begin{quote}\begin{multicols}{2}
\noindent (A) Samples \\
(B) Populations \\
(C) Complete study \\
(D) Complete universe
\end{multicols}\end{quote}

\item A collection of all the elements in a group is called:
\begin{quote}\begin{multicols}{2}
\noindent (A) Population \\
(B) Sample \\
(C) Data \\
(D) Registration
\end{multicols}\end{quote}

\item A collection of some of the elements from a group is called:
\begin{quote}\begin{multicols}{2}
\noindent (A) Census \\
(B) Population \\
(C) Sample \\
(D) Registration
\end{multicols}\end{quote}

\item Statistic is a characteristic calculated from:
\begin{quote}\begin{multicols}{2}
\noindent (A) Population data \\
(B) Sample data \\
(C) Arrayed data \\
(D) Fictitious data
\end{multicols}\end{quote}

\item The word statistics is at present used in:
\begin{quote}\begin{multicols}{2}
\noindent (A) Two senses \\
(B) Three senses \\
(C) Four senses \\
(D) None of these
\end{multicols}\end{quote}

\item If $k$ is a constant then $\displaystyle\sum_{i=1}^{5} k$:
\begin{quote}\begin{multicols}{2}
\noindent (A) $5k$ \\
(B) $5+k$ \\
(C) $nk$ \\
(D) $k_1+k_2+k_3+k_4+k_5$
\end{multicols}\end{quote}

\item Statistics came from the German word:
\begin{quote}\begin{multicols}{2}
\noindent (A) Status \\
(B) Statista \\
(C) Statistik \\
(D) Statistique
\end{multicols}\end{quote}

\item The mid point of group $5.5$---$7.5$ is:
\begin{quote}\begin{multicols}{2}
\noindent (A) 6 \\
(B) 7 \\
(C) 7.5 \\
(D) 6.5
\end{multicols}\end{quote}

\item In $60$---$70$, the lower limit is:
\begin{quote}\begin{multicols}{2}
\noindent (A) 50 \\
(B) 55 \\
(C) 60 \\
(D) 70
\end{multicols}\end{quote}

\item A portion of population selected for study is called:
\begin{quote}\begin{multicols}{2}
\noindent (A) Parameter \\
(B) Statistic \\
(C) Population \\
(D) Sample
\end{multicols}\end{quote}

\item If ``$a$'' is a constant, then $\displaystyle\sum_{i=1}^{5} a$ is equal to:
\begin{quote}\begin{multicols}{2}
\noindent (A) $a_1+a_2+a_3+a_4+a_5$ \\
(B) $a$ \\
(C) $5a$ \\
(D) None
\end{multicols}\end{quote}

\item Statistical laws are true:
\begin{quote}\begin{multicols}{2}
\noindent (A) On the average \\
(B) Always \\
(C) In the long run \\
(D) A and C
\end{multicols}\end{quote}

\item A numerical characteristic of a sample is called:
\begin{quote}\begin{multicols}{2}
\noindent (A) Statistic \\
(B) Parameter \\
(C) Variable \\
(D) Sample
\end{multicols}\end{quote}

\item The relative frequency multiplied by 100 is called:
\begin{quote}\begin{multicols}{2}
\noindent (A) Percentage frequency \\
(B) Cumulative frequency \\
(C) Sample frequency \\
(D) Bivariate frequency
\end{multicols}\end{quote}

\item A quantity calculated from a population is called:
\begin{quote}\begin{multicols}{2}
\noindent (A) Frequency \\
(B) Statistics \\
(C) Parameter \\
(D) Sample
\end{multicols}\end{quote}

\item Hourly temperature recorded by the Weather Bureau represents:
\begin{quote}\begin{multicols}{2}
\noindent (A) Discrete data \\
(B) Continuous data \\
(C) Qualitative data \\
(D) Secondary data
\end{multicols}\end{quote}

\item Weight of the whole earth is:
\begin{quote}\begin{multicols}{2}
\noindent (A) Discrete variable \\
(B) Qualitative variable \\
(C) Continuous variable \\
(D) Constant
\end{multicols}\end{quote}

\item Measurements usually provide:
\begin{quote}\begin{multicols}{2}
\noindent (A) Discrete data \\
(B) Continuous data \\
(C) Qualitative data \\
(D) Primary data
\end{multicols}\end{quote}

\item Counting or enumeration usually provides:
\begin{quote}\begin{multicols}{2}
\noindent (A) Continuous data \\
(B) Qualitative data \\
(C) Discrete data \\
(D) Grouped data
\end{multicols}\end{quote}

\item Number of accidents recorded yesterday in Lahore is a:
\begin{quote}\begin{multicols}{2}
\noindent (A) Discrete variable \\
(B) Continuous variable \\
(C) Qualitative variable \\
(D) Constant
\end{multicols}\end{quote}

\item The amount of milk given by a cow is a:
\begin{quote}\begin{multicols}{2}
\noindent (A) Qualitative variable \\
(B) Discrete variable \\
(C) Continuous variable \\
(D) Constant
\end{multicols}\end{quote}

\item Colour of hair is a:
\begin{quote}\begin{multicols}{2}
\noindent (A) Continuous variable \\
(B) Discrete variable \\
(C) Qualitative variable \\
(D) Quantitative variable
\end{multicols}\end{quote}

\item Smoking habits of residents of a city are:
\begin{quote}\begin{multicols}{2}
\noindent (A) Qualitative data \\
(B) Quantitative data \\
(C) Discrete data \\
(D) Continuous data
\end{multicols}\end{quote}

\item Continuous data are differentiated from discrete data in that:
\begin{quote}\begin{multicols}{2}
\noindent (A) Discrete data classes are represented by fractions \\
(B) Continuous data classes may be represented by fractions \\
(C) Continuous data take on only whole numbers \\
(D) Discrete data can take on any real number
\end{multicols}\end{quote}

\item Life of a T.V. tube is a:
\begin{quote}\begin{multicols}{2}
\noindent (A) Discrete variable \\
(B) Continuous variable \\
(C) Qualitative variable \\
(D) Constant
\end{multicols}\end{quote}

\item A constant can assume:
\begin{quote}\begin{multicols}{2}
\noindent (A) Only one value \\
(B) More than one value \\
(C) Different values \\
(D) No value at all
\end{multicols}\end{quote}

\item A statistic which is not measurable is called:
\begin{quote}\begin{multicols}{2}
\noindent (A) A constant \\
(B) An attribute \\
(C) A variable \\
(D) A parameter
\end{multicols}\end{quote}

\item The number $4.50001$ rounded off to the nearest unit (whole number) is:
\begin{quote}\begin{multicols}{2}
\noindent (A) 4 \\
(B) 5 \\
(C) 4.5 \\
(D) 4.01
\end{multicols}\end{quote}

\item The number $5.56500$ rounded off to the nearest hundredth (two decimal places) is:
\begin{quote}\begin{multicols}{2}
\noindent (A) 5.57 \\
(B) 5.56 \\
(C) 6.00 \\
(D) 5.00
\end{multicols}\end{quote}

\item The number $143.9500$ rounded off to the nearest tenth (one decimal place) is:
\begin{quote}\begin{multicols}{2}
\noindent (A) 143.9 \\
(B) 144.0 \\
(C) 143.0 \\
(D) 144
\end{multicols}\end{quote}

\item The number $5.0435001$ rounded off to the nearest thousandth (three decimal places) is:
\begin{quote}\begin{multicols}{2}
\noindent (A) 5.043 \\
(B) 5.044 \\
(C) 5.050 \\
(D) 5.000
\end{multicols}\end{quote}

\item The number $136.500$ rounded off to the nearest unit (whole number) is:
\begin{quote}\begin{multicols}{2}
\noindent (A) 136 \\
(B) 137 \\
(C) 136.5 \\
(D) 136.0
\end{multicols}\end{quote}

\item Continuous variable can be measured at:
\begin{quote}\begin{multicols}{2}
\noindent (A) Specific points \\
(B) All possible points \\
(C) No points \\
(D) Integer points
\end{multicols}\end{quote}

\item A quantitative variable whose values are countable is called:
\begin{quote}\begin{multicols}{2}
\noindent (A) Categorical variable \\
(B) Continuous variable \\
(C) Discrete variable \\
(D) None of these
\end{multicols}\end{quote}

\item Weight of students in a class makes:
\begin{quote}\begin{multicols}{2}
\noindent (A) Discrete data \\
(B) Continuous data \\
(C) Qualitative data \\
(D) Constant data
\end{multicols}\end{quote}

\item The lifetime of fans, data is:
\begin{quote}\begin{multicols}{2}
\noindent (A) Discrete \\
(B) Continuous \\
(C) Unchanged \\
(D) Qualitative
\end{multicols}\end{quote}

\item Statistics must be:
\begin{quote}\begin{multicols}{2}
\noindent (A) Comparable \\
(B) Not comparable \\
(C) Quantitative only \\
(D) None of these
\end{multicols}\end{quote}

\item The phase of statistics that is concerned with the description and analysis of sample or population data is called:
\begin{quote}\begin{multicols}{2}
\noindent (A) Inferential statistics \\
(B) Descriptive statistics \\
(C) Inductive statistics \\
(D) Sample statistics
\end{multicols}\end{quote}

\item The phase of statistics that is concerned with the procedures and methodology of obtaining valid conclusions is called:
\begin{quote}\begin{multicols}{2}
\noindent (A) Descriptive statistics \\
(B) Deductive statistics \\
(C) Inferential statistics \\
(D) Sample statistics
\end{multicols}\end{quote}

\item Data used by an agency which originally collected them are:
\begin{quote}\begin{multicols}{2}
\noindent (A) Primary data \\
(B) Raw data \\
(C) Secondary data \\
(D) Grouped data
\end{multicols}\end{quote}

\item Registration is the source of:
\begin{quote}\begin{multicols}{2}
\noindent (A) Primary data \\
(B) Secondary data \\
(C) Ogive \\
(D) Histogram
\end{multicols}\end{quote}

\item Registration in the Population Census Reports are:
\begin{quote}\begin{multicols}{2}
\noindent (A) Ungrouped data \\
(B) Secondary data \\
(C) Primary data \\
(D) Arrayed data
\end{multicols}\end{quote}

\item Information recorded in its original form, whether counts or measurements, is referred to as:
\begin{quote}\begin{multicols}{2}
\noindent (A) Continuous data \\
(B) Raw data \\
(C) Discrete data \\
(D) Arrayed data
\end{multicols}\end{quote}

\item Issuing a national identity card is an example of:
\begin{quote}\begin{multicols}{2}
\noindent (A) Census \\
(B) Registration \\
(C) Sampling \\
(D) Investigation through enumerators
\end{multicols}\end{quote}

\item The data which have not undergone any statistical treatment are:
\begin{quote}\begin{multicols}{2}
\noindent (A) Primary data \\
(B) Secondary data \\
(C) Discrete data \\
(D) Qualitative data
\end{multicols}\end{quote}

\item The data which have already been collected by someone else are called:
\begin{quote}\begin{multicols}{2}
\noindent (A) Secondary data \\
(B) Primary data \\
(C) Arrayed data \\
(D) None of these
\end{multicols}\end{quote}

\item How many sources are used for the collection of data:
\begin{quote}\begin{multicols}{2}
\noindent (A) 1 \\
(B) 2 \\
(C) 3 \\
(D) 4
\end{multicols}\end{quote}

\item How many methods are used for the collection of data:
\begin{quote}\begin{multicols}{2}
\noindent (A) 1 \\
(B) 2 \\
(C) 3 \\
(D) 4
\end{multicols}\end{quote}

\item Classification is the process of arranging data according to:
\begin{quote}\begin{multicols}{2}
\noindent (A) One characteristic \\
(B) Two or more characteristics \\
(C) Similar characteristics \\
(D) Different characteristics
\end{multicols}\end{quote}

\item Important bases of classification are:
\begin{quote}\begin{multicols}{2}
\noindent (A) Two \\
(B) Three \\
(C) Four \\
(D) More than four
\end{multicols}\end{quote}

\item Data classified by attributes are called:
\begin{quote}\begin{multicols}{2}
\noindent (A) Continuous data \\
(B) Quantitative data \\
(C) Qualitative data \\
(D) Grouped data
\end{multicols}\end{quote}

\item Data which have been arranged in ascending or descending order are called:
\begin{quote}\begin{multicols}{2}
\noindent (A) Raw data \\
(B) Grouped data \\
(C) Arrayed data \\
(D) Ungrouped data
\end{multicols}\end{quote}

\item As a general rule, when arranging data, statisticians tend to use:
\begin{quote}\begin{multicols}{2}
\noindent (A) Less than six classes \\
(B) Between six and fifteen classes \\
(C) More than twenty classes \\
(D) Only fifteen classes
\end{multicols}\end{quote}

\item Classification of data by their time of occurrence is called:
\begin{quote}\begin{multicols}{2}
\noindent (A) Temporal or chronological classification \\
(B) Geographical classification \\
(C) Quantitative classification \\
(D) Qualitative classification
\end{multicols}\end{quote}

\item Classification of data by attributes is called:
\begin{quote}\begin{multicols}{2}
\noindent (A) Temporal classification \\
(B) Geographical classification \\
(C) Quantitative classification \\
(D) Qualitative classification
\end{multicols}\end{quote}

\item Classification of data by quantitative characteristics is called:
\begin{quote}\begin{multicols}{2}
\noindent (A) Qualitative classification \\
(B) Quantitative classification \\
(C) Geographical classification \\
(D) Temporal classification
\end{multicols}\end{quote}

\item Classification of data according to locations or areas is called:
\begin{quote}\begin{multicols}{2}
\noindent (A) Temporal classification \\
(B) Geographical classification \\
(C) Quantitative classification \\
(D) Qualitative classification
\end{multicols}\end{quote}

\item The title of a table should be in:
\begin{quote}\begin{multicols}{2}
\noindent (A) Lower case letters \\
(B) Capital letters \\
(C) Italic and lower case letters \\
(D) Twenty letters
\end{multicols}\end{quote}

\item The process of arranging data into rows and columns is called:
\begin{quote}\begin{multicols}{2}
\noindent (A) Frequency distribution \\
(B) Classification \\
(C) Tabulation \\
(D) Array
\end{multicols}\end{quote}

\item A statistical table has at least:
\begin{quote}\begin{multicols}{2}
\noindent (A) Five parts \\
(B) Four parts \\
(C) Three parts \\
(D) Two parts
\end{multicols}\end{quote}

\item The section of a table that contains the column captions is called:
\begin{quote}\begin{multicols}{2}
\noindent (A) Box head \\
(B) Box plot \\
(C) Stub \\
(D) Body
\end{multicols}\end{quote}

\item Two methods of data arrangement are:
\begin{quote}\begin{multicols}{2}
\noindent (A) Array and frequency distribution \\
(B) Frequency distribution and histogram \\
(C) Array and frequency polygon \\
(D) Histogram and array
\end{multicols}\end{quote}

\item Which of the following is an example of compressed data:
\begin{quote}\begin{multicols}{2}
\noindent (A) Array \\
(B) Frequency distribution \\
(C) Histogram \\
(D) Ogive
\end{multicols}\end{quote}

\item Given classes $1-6, 6-10,\ldots$, the class interval is:
\begin{quote}\begin{multicols}{2}
\noindent (A) 5.5 \\
(B) 3 \\
(C) 4 \\
(D) 5
\end{multicols}\end{quote}

\item The grouped data is:
\begin{quote}\begin{multicols}{2}
\noindent (A) Primary \\
(B) Secondary \\
(C) Raw data \\
(D) None of these
\end{multicols}\end{quote}

\item An arrangement of data to show the frequency of occurrence is called:
\begin{quote}\begin{multicols}{2}
\noindent (A) Frequency distribution \\
(B) Probability distribution \\
(C) Data array \\
(D) Cumulative distribution
\end{multicols}\end{quote}

\item The average of the lower and upper class limits is called:
\begin{quote}\begin{multicols}{2}
\noindent (A) Class boundary \\
(B) Class frequency \\
(C) Class mark \\
(D) Class limit
\end{multicols}\end{quote}

\item Frequency is denoted by:
\begin{quote}\begin{multicols}{2}
\noindent (A) $c$ \\
(B) $q$ \\
(C) $P$ \\
(D) $f$
\end{multicols}\end{quote}

\item Mid point of the group $5.5$---$7.5$ is:
\begin{quote}\begin{multicols}{2}
\noindent (A) 6 \\
(B) 6.5 \\
(C) 7 \\
(D) 7.5
\end{multicols}\end{quote}

\item When constructing a frequency distribution, the first step is:
\begin{quote}\begin{multicols}{2}
\noindent (A) Divide the data into at least five classes \\
(B) Arrange the data into an array \\
(C) Decide on the type and number of classes for dividing the data \\
(D) None of these
\end{multicols}\end{quote}

\item Which of the following statements is true of cumulative frequency polygons or ogives for a particular set of data:
\begin{quote}\begin{multicols}{2}
\noindent (A) Both `less than' and `or more' curves have the same shape \\
(B) `Or more' curves slope up and to the right \\
(C) `Less than' curves slope down and to the right \\
(D) `Less than' curves slope up and to the right
\end{multicols}\end{quote}

\item In constructing a frequency distribution for a sample, the number of classes depends on:
\begin{quote}\begin{multicols}{2}
\noindent (A) The number of data points \\
(B) The range of the data collected \\
(C) The size of the population \\
(D) Both
\end{multicols}\end{quote}

\item A relative frequency distribution presents frequencies in terms of:
\begin{quote}\begin{multicols}{2}
\noindent (A) Fractions \\
(B) Whole numbers \\
(C) Percentages \\
(D) Both
\end{multicols}\end{quote}

\item A graph of a cumulative frequency distribution is called:
\begin{quote}\begin{multicols}{2}
\noindent (A) Histogram \\
(B) Frequency polygon \\
(C) Ogive \\
(D) None of these
\end{multicols}\end{quote}

\item An ogive is a:
\begin{quote}\begin{multicols}{2}
\noindent (A) Frequency curve \\
(B) Frequency polygon \\
(C) Cumulative frequency polygon \\
(D) Frequency histogram
\end{multicols}\end{quote}

\item If a curve can be divided into two parts that are mirror images, it is called a:
\begin{quote}\begin{multicols}{2}
\noindent (A) Skewed curve \\
(B) Symmetrical curve \\
(C) J-Shaped curve \\
(D) Frequency curve
\end{multicols}\end{quote}

\item Mid points of the top of the rectangles of a histogram are joined to get:
\begin{quote}\begin{multicols}{2}
\noindent (A) Polygon \\
(B) Frequency curve \\
(C) Ogive \\
(D) Histogram
\end{multicols}\end{quote}

\item As the number of observations and classes increase, the shape of a frequency polygon:
\begin{quote}\begin{multicols}{2}
\noindent (A) Tends to become increasingly smooth \\
(B) Tends to become jagged \\
(C) Stays the same \\
(D) Varies if data become more reliable
\end{multicols}\end{quote}

\item Graphs of frequency distributions are used because:
\begin{quote}\begin{multicols}{2}
\noindent (A) They have a long history in practical applications \\
(B) They attract attention to data patterns \\
(C) They account for biased or incomplete data \\
(D) None of these
\end{multicols}\end{quote}

\item Continuous data are differentiated from discrete data in that:
\begin{quote}\begin{multicols}{2}
\noindent (A) Discrete data classes are represented by fractions \\
(B) Continuous data classes may be represented by fractions \\
(C) Continuous data take on only whole numbers \\
(D) Discrete data can take on any real number
\end{multicols}\end{quote}

\item If a curve has a longer tail to the right, it is called a:
\begin{quote}\begin{multicols}{2}
\noindent (A) Positively skewed curve \\
(B) Negatively skewed curve \\
(C) J-shaped curve \\
(D) Symmetrical curve
\end{multicols}\end{quote}

\item If a curve has a longer tail to the left, it is called a:
\begin{quote}\begin{multicols}{2}
\noindent (A) Symmetrical curve \\
(B) Positively skewed curve \\
(C) Negatively skewed curve \\
(D) None of these
\end{multicols}\end{quote}

\item A chart in which total magnitude and its components are compared is called a:
\begin{quote}\begin{multicols}{2}
\noindent (A) Component bar chart \\
(B) Pie chart \\
(C) Percentage component bar chart \\
(D) All of these
\end{multicols}\end{quote}

\item Total angle of a pie-chart is:
\begin{quote}\begin{multicols}{2}
\noindent (A) $270^\circ$ \\
(B) $300^\circ$ \\
(C) $320^\circ$ \\
(D) $360^\circ$
\end{multicols}\end{quote}

\item Histogram is a graph of:
\begin{quote}\begin{multicols}{2}
\noindent (A) Frequency distribution \\
(B) Time series \\
(C) Qualitative data \\
(D) Ogive
\end{multicols}\end{quote}

\item The number of classes in a frequency distribution is obtained by dividing the range of the variable by the:
\begin{quote}\begin{multicols}{2}
\noindent (A) Total frequency \\
(B) Class interval \\
(C) Mid-points \\
(D) Relative frequency
\end{multicols}\end{quote}

\item A frequency polygon is a closed figure which is:
\begin{quote}\begin{multicols}{2}
\noindent (A) One sided \\
(B) Two sided \\
(C) Three sided \\
(D) Many sided
\end{multicols}\end{quote}

\item The frequency of a class divided by the total frequency is called:
\begin{quote}\begin{multicols}{2}
\noindent (A) Class frequency \\
(B) Cumulative \\
(C) Relative frequency \\
(D) Total frequency
\end{multicols}\end{quote}

\item Brand of a soap is \underline{\hspace{1.2cm}} variable:
\begin{quote}\begin{multicols}{2}
\noindent (A) Quantitative \\
(B) Qualitative \\
(C) Imaginary \\
(D) Continuous
\end{multicols}\end{quote}

\item The sum of rows or sum of columns of a bivariate frequency distribution are equal:
\begin{quote}\begin{multicols}{2}
\noindent (A) $\Sigma f$ \\
(B) $\Sigma X$ \\
(C) $\Sigma fX$ \\
(D) 100
\end{multicols}\end{quote}

\item What is the major assumption we make when computing a mean from grouped data:
\begin{quote}\begin{multicols}{2}
\noindent (A) All values are discrete \\
(B) Every value in a class is equal to the midpoint \\
(C) No value occurs more than once \\
(D) Each class contains exactly the same number of values
\end{multicols}\end{quote}

\item When calculating the average rate of debt growth for a company, the correct mean to use is:
\begin{quote}\begin{multicols}{2}
\noindent (A) Arithmetic mean \\
(B) Weighted arithmetic mean \\
(C) Geometric mean \\
(D) None of these
\end{multicols}\end{quote}

\item Which of the following is the first step in calculating the median of a data set:
\begin{quote}\begin{multicols}{2}
\noindent (A) Average the middle two values of the data set \\
(B) Array the data \\
(C) Determine the relative weights of the data values in terms of importance \\
(D) None of these
\end{multicols}\end{quote}

\item Departure from symmetry is called:
\begin{quote}\begin{multicols}{2}
\noindent (A) Kurtosis \\
(B) Skewness \\
(C) Dispersion \\
(D) None of these
\end{multicols}\end{quote}

\item When a distribution is symmetrical and has one mode, the highest point on the curve is called the:
\begin{quote}\begin{multicols}{2}
\noindent (A) Mode \\
(B) Median \\
(C) Mean \\
(D) All of these
\end{multicols}\end{quote}

\item When referring to a curve that tails off to the left end, you would call it:
\begin{quote}\begin{multicols}{2}
\noindent (A) Symmetrical \\
(B) Skewed to the right \\
(C) Positively skewed \\
(D) None of these
\end{multicols}\end{quote}

\item In which of these cases would the mode be most useful as an indicator of central tendency:
\begin{quote}\begin{multicols}{2}
\noindent (A) Every value in a data set occurs exactly once \\
(B) All but three values in a data set occur once; three values occur 5 times each \\
(C) All values in a data set occur 10 times each \\
(D) Every observation in a data set has the same value
\end{multicols}\end{quote}

\item It is the reciprocal of the simple average of the reciprocal of all the values:
\begin{quote}\begin{multicols}{2}
\noindent (A) A.M. \\
(B) G.M. \\
(C) H.M. \\
(D) Mode
\end{multicols}\end{quote}

\item When referring to a curve whose longer tail is to the right, you would call it:
\begin{quote}\begin{multicols}{2}
\noindent (A) Symmetrical \\
(B) Positively skewed \\
(C) Negatively skewed \\
(D) None of these
\end{multicols}\end{quote}

\item If the mean is less than the mode, the distribution is:
\begin{quote}\begin{multicols}{2}
\noindent (A) Positively skewed \\
(B) Negatively skewed \\
(C) Symmetrical \\
(D) None of these
\end{multicols}\end{quote}

\item The mean is affected by:
\begin{quote}\begin{multicols}{2}
\noindent (A) Change of origin \\
(B) Scale of measurement \\
(C) Both \\
(D) None of these
\end{multicols}\end{quote}

\item Sum of squares of deviations of the values is least when deviations are taken from:
\begin{quote}\begin{multicols}{2}
\noindent (A) Median \\
(B) Mode \\
(C) Mean \\
(D) Harmonic mean
\end{multicols}\end{quote}

\item When all the values in a series occur the same number of times, then it is not possible to compute the:
\begin{quote}\begin{multicols}{2}
\noindent (A) Mean \\
(B) Median \\
(C) Mode \\
(D) Weighted mean
\end{multicols}\end{quote}

\item The most central value of an arrayed data is:
\begin{quote}\begin{multicols}{2}
\noindent (A) Mode \\
(B) Median \\
(C) Mean \\
(D) Harmonic mean
\end{multicols}\end{quote}

\item The geometric mean of $X_1$ and $X_2$ is:
\begin{quote}\begin{multicols}{2}
\noindent (A) $\dfrac{2}{X_1+X_2}$ \\
(B) $\dfrac{2X_1X_2}{X_1+X_2}$ \\
(C) $\dfrac{X_1+X_2}{2}$ \\
(D) $\sqrt{X_1X_2}$
\end{multicols}\end{quote}

\item Mode of $2, 10$ and $7$ is:
\begin{quote}\begin{multicols}{2}
\noindent (A) 10 \\
(B) 2 \\
(C) 7 \\
(D) None of these
\end{multicols}\end{quote}

\item The mean of the first $n$ natural numbers is:
\begin{quote}\begin{multicols}{2}
\noindent (A) $n(n+1)/2$ \\
(B) $(n+1)/2$ \\
(C) $(n-1)/2$ \\
(D) $n/2$
\end{multicols}\end{quote}

\item The suitable average for the qualitative data is:
\begin{quote}\begin{multicols}{2}
\noindent (A) Mean \\
(B) Mode \\
(C) Geometric mean \\
(D) Weighted mean
\end{multicols}\end{quote}

\item If $\bar{X}=10$ and $Y=2X+5$, then $\bar{Y}=$:
\begin{quote}\begin{multicols}{2}
\noindent (A) 20 \\
(B) 25 \\
(C) 30 \\
(D) None of these
\end{multicols}\end{quote}

\item If $\displaystyle\sum_{i=1}^{5}(X_i-20)=0$, then $\bar{X}=$:
\begin{quote}\begin{multicols}{2}
\noindent (A) 0 \\
(B) 20 \\
(C) 5 \\
(D) 25
\end{multicols}\end{quote}

\item Mode of the series $2,2,2,3,3,3,2,3,3,4$ is:
\begin{quote}\begin{multicols}{2}
\noindent (A) 3 \\
(B) 2 and 3 \\
(C) 4 \\
(D) None of these
\end{multicols}\end{quote}

\item Coding method is used for calculation of the:
\begin{quote}\begin{multicols}{2}
\noindent (A) Median \\
(B) Mode \\
(C) Mean \\
(D) Weighted mean
\end{multicols}\end{quote}

\item In a symmetrical distribution:
\begin{quote}\begin{multicols}{2}
\noindent (A) $Q_1=Q_3$ \\
(B) $P_{25}=P_{50}=P_{75}$ \\
(C) A.M. = G.M. = H.M. \\
(D) A.M. = Med = Mode
\end{multicols}\end{quote}

\item The mean of 10 observations is 10. All observations are increased by 10\%. The mean of the increased observations shall be:
\begin{quote}\begin{multicols}{2}
\noindent (A) 20 \\
(B) 11 \\
(C) 10 \\
(D) 100
\end{multicols}\end{quote}

\item Which is appropriate average for finding the average speed of a journey:
\begin{quote}\begin{multicols}{2}
\noindent (A) Mean \\
(B) Geometric mean \\
(C) Harmonic mean \\
(D) Weighted mean
\end{multicols}\end{quote}

\item Which is the suitable average for calculating average percent increace in population:
\begin{quote}\begin{multicols}{2}
\noindent (A) Median \\
(B) Geometric mean \\
(C) Harmonic mean \\
(D) Mean
\end{multicols}\end{quote}

\item For a positively skewed distribution:
\begin{quote}\begin{multicols}{2}
\noindent (A) Mean > Mode \\
(B) Mode > Mean \\
(C) Median > mean \\
(D) None of these
\end{multicols}\end{quote}

\item If any value in a series is zero, then we cannot calculate the:
\begin{quote}\begin{multicols}{2}
\noindent (A) Mean \\
(B) Harmonic mean \\
(C) Mode \\
(D) Median
\end{multicols}\end{quote}

\item If the values in a series are not of equal importance, we compute the:
\begin{quote}\begin{multicols}{2}
\noindent (A) Median \\
(B) Harmonic mean \\
(C) Mean \\
(D) Weighted mean
\end{multicols}\end{quote}

\item Which is the suitable average for calculating the average price at which articles are sold:
\begin{quote}\begin{multicols}{2}
\noindent (A) Geometric mean \\
(B) Arithmetic mean \\
(C) Harmonic median \\
(D) Mode
\end{multicols}\end{quote}

\item For a negatively skewed distribution:
\begin{quote}\begin{multicols}{2}
\noindent (A) Mean > Mode \\
(B) Mode > Mean \\
(C) Mean > Median \\
(D) None of these
\end{multicols}\end{quote}

\item Which of the following measures cannot be calculated for the numbers 5, 8, 12, 6, 9, 13, 10:
\begin{quote}\begin{multicols}{2}
\noindent (A) Median \\
(B) Mode \\
(C) Mean \\
(D) None of these
\end{multicols}\end{quote}

\item The suitable average for shoe or collar size is:
\begin{quote}\begin{multicols}{2}
\noindent (A) Geometric mean \\
(B) Arightmetic mean \\
(C) Mode \\
(D) Median
\end{multicols}\end{quote}

\item Which pair of measures cannot be calculated when one of numbers in the series is zero:
\begin{quote}\begin{multicols}{2}
\noindent (A) G.M. and A.M. \\
(B) G.M and H.M. \\
(C) H.M. and A.M. \\
(D) None of these
\end{multicols}\end{quote}

\item The sum of deviations from arithmetic mean is $\Sigma (X - \bar{X}) = $:
\begin{quote}\begin{multicols}{2}
\noindent (A) 0 \\
(B) 1 \\
(C) 2 \\
(D) 3
\end{multicols}\end{quote}

\item $\Sigma (X - 20) = 25$, $\Sigma (X - 18) = 0$ then arithmetic mean will be:
\begin{quote}\begin{multicols}{2}
\noindent (A) 25 \\
(B) 18 \\
(C) 38 \\
(D) 20
\end{multicols}\end{quote}

\item If Mean = 40, Mode = 42, then distribution is:
\begin{quote}\begin{multicols}{2}
\noindent (A) Positively skewed \\
(B) Negatively skewed \\
(C) Symmetrical \\
(D) All of these
\end{multicols}\end{quote}

\item Which average cannot be computed if any value is less than zero:
\begin{quote}\begin{multicols}{2}
\noindent (A) G.M. \\
(B) Median \\
(C) Mode \\
(D) A.M.
\end{multicols}\end{quote}

\item Geometric mean of numbers 2, 4, 8 and 64 is:
\begin{quote}\begin{multicols}{2}
\noindent (A) 2 \\
(B) 4 \\
(C) 8 \\
(D) 64
\end{multicols}\end{quote}

\item The Mean is based on:
\begin{quote}\begin{multicols}{2}
\noindent (A) All the values \\
(B) Small values \\
(C) Extreme values \\
(D) Large values
\end{multicols}\end{quote}

\item The arithmetic mean of 'n' values $y_1, y_2, ...., y_n$ is:
\begin{quote}\begin{multicols}{2}
\noindent (A) $\dfrac{\Sigma fy}{\Sigma f}$ \\
(B) $y_1+y_2+y_3+...+y_n$ \\
(C) $\dfrac{\Sigma y}{\Sigma yf}$ \\
(D) $\dfrac{y_1+y_2+y_3+...+y_n}{n}$
\end{multicols}\end{quote}

\item Median divides the data into:
\begin{quote}\begin{multicols}{2}
\noindent (A) 2 parts \\
(B) 3 parts \\
(C) 4 parts \\
(D) 10 parts
\end{multicols}\end{quote}

\item If a distribution has two modes, then it is called:
\begin{quote}\begin{multicols}{2}
\noindent (A) Uni-model \\
(B) Bi-model \\
(C) Tri-model \\
(D) Multi model
\end{multicols}\end{quote}

\item The most frequent value of the data if it exists is:
\begin{quote}\begin{multicols}{2}
\noindent (A) A.M. \\
(B) G.M. \\
(C) Mode \\
(D) Median
\end{multicols}\end{quote}

\item Mean is affected by the change of:
\begin{quote}\begin{multicols}{2}
\noindent (A) Origin \\
(B) Scale \\
(C) Both a and b \\
(D) None
\end{multicols}\end{quote}

\item If n=7 and $\Sigma (1/y) = 1.929$, then H.M. is:
\begin{quote}\begin{multicols}{2}
\noindent (A) 5.77 \\
(B) 3.63 \\
(C) 5.89 \\
(D) 5.97
\end{multicols}\end{quote}

\item If X=0, 2, 2, 4, 8, 10, then G.M. is:
\begin{quote}\begin{multicols}{2}
\noindent (A) 4 \\
(B) 8 \\
(C) Zero \\
(D) 10
\end{multicols}\end{quote}

\item If $\Sigma (X-20)^2=0$, $\Sigma (X-18)^2=10$ then $\bar{X}$ could be:
\begin{quote}\begin{multicols}{2}
\noindent (A) 20 \\
(B) 18 \\
(C) 28 \\
(D) 25
\end{multicols}\end{quote}

\item The median of -3, 0, -5, is:
\begin{quote}\begin{multicols}{2}
\noindent (A) -3 \\
(B) 0 \\
(C) -5 \\
(D) Does not exist
\end{multicols}\end{quote}

\item Which of the following average cannot be calculated from the observation 2, 2, 4, 4, 6, 6, 8, 8, 10, 10:
\begin{quote}\begin{multicols}{2}
\noindent (A) Mean \\
(B) Median \\
(C) Mode \\
(D) All of these
\end{multicols}\end{quote}

\item In case of positively skewed distribution the extreme values lie in the:
\begin{quote}\begin{multicols}{2}
\noindent (A) Middle \\
(B) Left tail \\
(C) Right tail \\
(D) Anywhere
\end{multicols}\end{quote}

\item Which of the following averages is affected by extreme values:
\begin{quote}\begin{multicols}{2}
\noindent (A) Arithmetic mean \\
(B) Median \\
(C) Mode \\
(D) All of these
\end{multicols}\end{quote}

\item $\Sigma (Y - \bar{Y}) = $:
\begin{quote}\begin{multicols}{2}
\noindent (A) 0 \\
(B) 1 \\
(C) Least \\
(D) > 0
\end{multicols}\end{quote}

\item Which of the following is a relative measure of dispersion:
\begin{quote}\begin{multicols}{2}
\noindent (A) Standard deviation \\
(B) Variance \\
(C) Coefficient of variation \\
(D) All of these
\end{multicols}\end{quote}

\item If you divide the standard deviation of the values by the mean of the same values and multiply by 100, you will get:
\begin{quote}\begin{multicols}{2}
\noindent (A) Variance \\
(B) Coefficient of variation \\
(C) Mean deviation \\
(D) Mean coefficient of dispersion
\end{multicols}\end{quote}

\item The position square root of the variance of a distribution is the:
\begin{quote}\begin{multicols}{2}
\noindent (A) Standard deviation \\
(B) Mean deviation \\
(C) Absolute deviation \\
(D) None of these
\end{multicols}\end{quote}

\item The main advantage of using the range as a measure of dispersion is that:
\begin{quote}\begin{multicols}{2}
\noindent (A) It is easy to calculate \\
(B) It is heavily influenced by extreme values \\
(C) It can change drastically from one data set to the next \\
(D) It is determined by only two points in the data set
\end{multicols}\end{quote}

\item Why is it necessary to square the differences from the mean when computing the standard deviation:
\begin{quote}\begin{multicols}{2}
\noindent (A) So that the extreme values will not affect the calculation \\
(B) Some of the differences will be positive and some will be negative \\
(C) The number of values involved is very small \\
(D) None of these
\end{multicols}\end{quote}

\item The sum of absolute deviations is a minimum if these deviations are taken from the:
\begin{quote}\begin{multicols}{2}
\noindent (A) Mean \\
(B) Median \\
(C) Mode \\
(D) All of these
\end{multicols}\end{quote}

\item Range can be calculated in open-end classes:
\begin{quote}\begin{multicols}{2}
\noindent (A) Never \\
(B) Always \\
(C) Often \\
(D) Seldom
\end{multicols}\end{quote}

\item If s.d. (X) = 5, then s.d. $\left(\dfrac{2X+5}{2}\right)$ is equal to:
\begin{quote}\begin{multicols}{2}
\noindent (A) 5 \\
(B) 10 \\
(C) 15 \\
(D) 7.5
\end{multicols}\end{quote}

\item The standard deviation is independent of:
\begin{quote}\begin{multicols}{2}
\noindent (A) Change of origin \\
(B) Change of scale of measurement \\
(C) Change of origin and scale of measurement \\
(D) None of these
\end{multicols}\end{quote}

\item If $\bar{X}=5$, which of the following expressions is a minimum:
\begin{quote}\begin{multicols}{2}
\noindent (A) $\Sigma(X-25)^2$ \\
(B) $\Sigma(X-5)^2$ \\
(C) $\Sigma|X-5|$ \\
(D) $\Sigma|X-25|$
\end{multicols}\end{quote}

\item For moderately skewed distributions, approximately 68\% of the cases are included between:
\begin{quote}\begin{multicols}{2}
\noindent (A) $\bar{X}-S$ and $\bar{X}+S$ \\
(B) $\bar{X}-2S$ and $\bar{X}+2S$ \\
(C) $\bar{X}-3S$ and $\bar{X}+3S$ \\
(D) None of these
\end{multicols}\end{quote}

\item For moderately skewed distributions, the approximate percentage of cases included between $\bar{X}-2S$ and $\bar{X}+2S$ is:
\begin{quote}\begin{multicols}{2}
\noindent (A) 99.7\\% \\
(B) 68\\% \\
(C) 95\\% \\
(D) 50\\%
\end{multicols}\end{quote}

\item Which of the following measures of dispersion is independent of the units employed:
\begin{quote}\begin{multicols}{2}
\noindent (A) Standard deviation \\
(B) Quartile deviation \\
(C) Coefficient of variation \\
(D) Variance
\end{multicols}\end{quote}

\item If X and Y are independent, then Var(X - Y) is equal to:
\begin{quote}\begin{multicols}{2}
\noindent (A) Var\\{X\\} + Var\\{Y\\} \\
(B) Var\\{X\\} - Var\\{Y\\} \\
(C) Var\\{X + Y\\} \\
(D) Zero
\end{multicols}\end{quote}

\item If X and Y are independent, then s.d. (X - Y) is equal to:
\begin{quote}\begin{multicols}{2}
\noindent (A) $\sqrt{Var(X) - Var(Y)}$ \\
(B) $\sqrt{Var(X) + Var(Y)}$ \\
(C) S.D.\\{X\\} + S.D.\\{Y\\} \\
(D) S.D.\\{X\\} - S.D.\\{Y\\}
\end{multicols}\end{quote}

\item For a moderately skewed distribution, which of the following empirical formula holds:
\begin{quote}\begin{multicols}{2}
\noindent (A) $M.D. = \dfrac{4}{5}(S.D)$ \\
(B) $Q.D. = \dfrac{2}{3}(S.D)$ \\
(C) $Q.D. = \dfrac{5}{6}(M.D)$ \\
(D) All of these
\end{multicols}\end{quote}

\item The mean deviation is least if deviations are taken from:
\begin{quote}\begin{multicols}{2}
\noindent (A) Median \\
(B) A.M. \\
(C) Mode \\
(D) G.M.
\end{multicols}\end{quote}

\item If the third moment about mean is zero ($m_3=0$), then the distribution is:
\begin{quote}\begin{multicols}{2}
\noindent (A) Mesokurtic \\
(B) Positively skewed \\
(C) Symmetrical \\
(D) Negatively skewed
\end{multicols}\end{quote}

\item If the moment coefficient of kurtosis $b_2=3$, then the distribution is:
\begin{quote}\begin{multicols}{2}
\noindent (A) Platykurtic \\
(B) Positively skewed \\
(C) Symmetrical \\
(D) Mesokurtic
\end{multicols}\end{quote}

\item Which measure of dispersion is considered as the best general purpose measure of dispersion:
\begin{quote}\begin{multicols}{2}
\noindent (A) Range \\
(B) Semi-interquartile range \\
(C) Standard deviation \\
(D) Mean deviation
\end{multicols}\end{quote}

\item To compare the variation of two or more than two series, we use:
\begin{quote}\begin{multicols}{2}
\noindent (A) Mean absolute deviation \\
(B) Variance \\
(C) Coefficient of variation \\
(D) Corrected standard deviation
\end{multicols}\end{quote}

\item The distribution is symmetrical if the moment coefficient of skewness $\sqrt{b_1}$ is:
\begin{quote}\begin{multicols}{2}
\noindent (A) Negative \\
(B) Positive \\
(C) 0 \\
(D) 3
\end{multicols}\end{quote}

\item The distribution is mesokurtic if the moment coefficient of kurtosis $b_2$ is:
\begin{quote}\begin{multicols}{2}
\noindent (A) Equal to 0 \\
(B) Equal to 3 \\
(C) Less than 3 \\
(D) Greater than zero
\end{multicols}\end{quote}

\item The distribution is positively skewed if:
\begin{quote}\begin{multicols}{2}
\noindent (A) Mean < Mode \\
(B) Mean > Mode \\
(C) Mean > Median \\
(D) Both
\end{multicols}\end{quote}

\item Moment ratios $b_1$ and $b_2$ are:
\begin{quote}\begin{multicols}{2}
\noindent (A) Expressed in original unit of the data \\
(B) Dimensionless quantities \\
(C) Independent of origin and scale of measurement \\
(D) Both
\end{multicols}\end{quote}

\item For moderately skewed distributions, approximately 80\% of the cases are included between:
\begin{quote}\begin{multicols}{2}
\noindent (A) $\bar{X}-S$ and $\bar{X}+S$ \\
(B) $\bar{X}-2S$ and $\bar{X}+2S$ \\
(C) $\bar{X}-3S$ and $\bar{X}+3S$ \\
(D) None of these
\end{multicols}\end{quote}

\item The sum of squares of deviations is a minimum if these deviations are taken from the:
\begin{quote}\begin{multicols}{2}
\noindent (A) Median \\
(B) Mean \\
(C) Mode \\
(D) All of these
\end{multicols}\end{quote}

\item Mean deviation is always:
\begin{quote}\begin{multicols}{2}
\noindent (A) More than S.D. \\
(B) Equal to S.D. \\
(C) Less than S.D. \\
(D) None of these
\end{multicols}\end{quote}

\item If X and Y are independent, then s.d. (X + Y) is equal to:
\begin{quote}\begin{multicols}{2}
\noindent (A) $\sqrt{Var(X)+Var(Y)}$ \\
(B) s.d.(X) + s.d.(Y) \\
(C) $\sqrt{Var(X) \cdot Var(Y)}$ \\
(D) None of these
\end{multicols}\end{quote}

\item Which of the following statements is correct:
\begin{quote}\begin{multicols}{2}
\noindent (A) Every symmetrical curve is mesokurtic \\
(B) Standard deviation is the mean squared deviations from the mean \\
(C) The standard deviation of a constant is constant \\
(D) The second moment about zero equals variance
\end{multicols}\end{quote}

\item The types of dispersion are:
\begin{quote}\begin{multicols}{2}
\noindent (A) 2 \\
(B) 3 \\
(C) 4 \\
(D) 5
\end{multicols}\end{quote}

\item For a symmetrical distribution:
\begin{quote}\begin{multicols}{2}
\noindent (A) $\beta_1=0$ \\
(B) $\beta_1=3$ \\
(C) $\beta_2=3$ \\
(D) $\beta_3=3$
\end{multicols}\end{quote}

\item Second moment about mean is:
\begin{quote}\begin{multicols}{2}
\noindent (A) 0 \\
(B) 1 \\
(C) Variance \\
(D) Standard deviation
\end{multicols}\end{quote}

\item The variance of 5, 5, 5, 5, 5 is:
\begin{quote}\begin{multicols}{2}
\noindent (A) 0 \\
(B) 25 \\
(C) 5 \\
(D) 125
\end{multicols}\end{quote}

\item In symmetrical distribution if $Q_1=4$, $Q_3=12$ then median is:
\begin{quote}\begin{multicols}{2}
\noindent (A) 4 \\
(B) 6 \\
(C) 8 \\
(D) Zero
\end{multicols}\end{quote}

\item A disadvantage of range is that it is based on:
\begin{quote}\begin{multicols}{2}
\noindent (A) Absolute deviation \\
(B) Square deviation \\
(C) Two extreme observation \\
(D) Upper and quartile
\end{multicols}\end{quote}

\item Relative depression is of types:
\begin{quote}\begin{multicols}{2}
\noindent (A) 1 \\
(B) 4 \\
(C) 3 \\
(D) 2
\end{multicols}\end{quote}

\item If Y = X + A, the range of Y = \_\_\_\_\_\_\_\_:
\begin{quote}\begin{multicols}{2}
\noindent (A) Range \\
(B) Range \\
(C) Zero \\
(D) A
\end{multicols}\end{quote}

\item For symmetrical distribution $\mu_3$ is:
\begin{quote}\begin{multicols}{2}
\noindent (A) Zero \\
(B) 1 \\
(C) 3 \\
(D) 4
\end{multicols}\end{quote}

\item If the values of a variable are -2, -3, -5, -10 then range is:
\begin{quote}\begin{multicols}{2}
\noindent (A) -12 \\
(B) 8 \\
(C) -8 \\
(D) 0
\end{multicols}\end{quote}

\item The measures of dispersion are changed by the change of:
\begin{quote}\begin{multicols}{2}
\noindent (A) Origin \\
(B) Scale \\
(C) Both \\
(D) None of these
\end{multicols}\end{quote}

\item M.D. of the values 4, 4, 4, 4 is:
\begin{quote}\begin{multicols}{2}
\noindent (A) 0 \\
(B) 4 \\
(C) 8 \\
(D) 12
\end{multicols}\end{quote}

\item In Sheppard's correction $\mu_2$ is equal to:
\begin{quote}\begin{multicols}{2}
\noindent (A) $\mu^2 - \dfrac{h^2}{12}$ \\
(B) $\mu^2 + \dfrac{h^2}{12}$ \\
(C) $\mu^2 - \dfrac{h}{2}$ \\
(D) $\mu + \dfrac{h}{12}$
\end{multicols}\end{quote}

\item The mean deviation of dispersion can be negative:
\begin{quote}\begin{multicols}{2}
\noindent (A) Never \\
(B) Always \\
(C) Sometimes \\
(D) Often
\end{multicols}\end{quote}

\item For Leptokurtic distribution:
\begin{quote}\begin{multicols}{2}
\noindent (A) $b_2 > 3$ \\
(B) $b_2 < 3$ \\
(C) $b_2 = 3$ \\
(D) $b_1 > 3$
\end{multicols}\end{quote}

\item First central moment is always:
\begin{quote}\begin{multicols}{2}
\noindent (A) Zero \\
(B) 1 \\
(C) -1 \\
(D) 2
\end{multicols}\end{quote}

\item The S.D. of 8, 8, 8, 8, 8 is:
\begin{quote}\begin{multicols}{2}
\noindent (A) 8 \\
(B) $8^2$ \\
(C) $\sqrt{8}$ \\
(D) Zero
\end{multicols}\end{quote}

\item Which set has the minimum variation:
\begin{quote}\begin{multicols}{2}
\noindent (A) 46, 48, 50 \\
(B) 30, 40, 50 \\
(C) 40, 50, 60 \\
(D) 48, 48, 49
\end{multicols}\end{quote}

\item In a symmetrical distribution, the coefficient of skewness will always be:
\begin{quote}\begin{multicols}{2}
\noindent (A) Negative \\
(B) Zero \\
(C) 1 \\
(D) -1
\end{multicols}\end{quote}

\item If $Q_3=20$ and $Q_1=10$ the coefficient of quartile deviation is:
\begin{quote}\begin{multicols}{2}
\noindent (A) 3 \\
(B) $\dfrac{1}{3}$ \\
(C) $\dfrac{2}{3}$ \\
(D) 1
\end{multicols}\end{quote}

\item The value of standard deviation changes by change of:
\begin{quote}\begin{multicols}{2}
\noindent (A) Origin \\
(B) Scale \\
(C) Algebraic sign \\
(D) None
\end{multicols}\end{quote}

\item The range of the scores 19, 3, 140, 25, 95 is:
\begin{quote}\begin{multicols}{2}
\noindent (A) 140 \\
(B) 3 \\
(C) 137 \\
(D) 143
\end{multicols}\end{quote}

\item The range of the values -2, -4, -6 and -8 is:
\begin{quote}\begin{multicols}{2}
\noindent (A) -6 \\
(B) 6 \\
(C) -10 \\
(D) -4
\end{multicols}\end{quote}

\item Var (bX) = \_\_\_\_\_\_\_:
\begin{quote}\begin{multicols}{2}
\noindent (A) b Var \\
(B) $b^2$Var \\
(C) Var \\
(D) Var
\end{multicols}\end{quote}

\item The variance expresses the variability of data in \_\_\_\_\_\_ as unit of data:
\begin{quote}\begin{multicols}{2}
\noindent (A) Square of unit \\
(B) Square root of unit \\
(C) Same unit \\
(D) All of these
\end{multicols}\end{quote}

\item If an index number calculation over 8 years with a base value of 100 gave an index for 1993 of 110, what would be the percentage relative for 1993:
\begin{quote}\begin{multicols}{2}
\noindent (A) 110 \\
(B) 90.0 \\
(C) 13.75 \\
(D) 880
\end{multicols}\end{quote}

\item If all the values considered in calculating an index are of equal importance, the index is:
\begin{quote}\begin{multicols}{2}
\noindent (A) Weighted \\
(B) Simple \\
(C) Unweighted \\
(D) None of these
\end{multicols}\end{quote}

\item When the base year values are used as weights, the weighted average of relatives price index is the same as:
\begin{quote}\begin{multicols}{2}
\noindent (A) The Paasche's index \\
(B) The Laspeyre's index \\
(C) The unweighted average of relatives price index \\
(D) None of these
\end{multicols}\end{quote}

\item Simple index number involves commodities:
\begin{quote}\begin{multicols}{2}
\noindent (A) One \\
(B) Two \\
(C) Four \\
(D) None of these
\end{multicols}\end{quote}

\item If an unweighted average of relatives index $\dfrac{p_n}{p_0}\times 100$ is calculated for each product in the composite, what is then done with these values to finish the calculation:
\begin{quote}\begin{multicols}{2}
\noindent (A) The values are multiplied together \\
(B) The largest value is found \\
(C) The values are averaged \\
(D) None of these
\end{multicols}\end{quote}

\item The weights used in a quantity index are:
\begin{quote}\begin{multicols}{2}
\noindent (A) Percentages of total quantity \\
(B) Prices \\
(C) Average of quantities \\
(D) None of these
\end{multicols}\end{quote}

\item The weights used in a price index are:
\begin{quote}\begin{multicols}{2}
\noindent (A) Percentage of total price \\
(B) Quantities \\
(C) Average of prices \\
(D) None of these
\end{multicols}\end{quote}

\item The base period can be described as a 'norma' period if:
\begin{quote}\begin{multicols}{2}
\noindent (A) It is neither the peak nor the trough of a fluctuation \\
(B) It is the most recent period for which we have data \\
(C) It is the average of several consecutive periods \\
(D) None of these
\end{multicols}\end{quote}

\item Link Relative is equal to:
\begin{quote}\begin{multicols}{2}
\noindent (A) $\dfrac{P_n}{P_0}\times 100$ \\
(B) $\dfrac{P_n}{P_{n-1}}\times 100$ \\
(C) $\dfrac{P_o}{P_n}\times 100$ \\
(D) $\dfrac{P_n}{P_n}\times 100$
\end{multicols}\end{quote}

\item Which of the following methods uses quantities consumed in the current period when computing a weighted index:
\begin{quote}\begin{multicols}{2}
\noindent (A) Laspeyre's method \\
(B) Paasche's method \\
(C) Marshall - Edgeworth's method \\
(D) Fisher's method
\end{multicols}\end{quote}

\item Which of the following methods uses quantities consumed in the base period when computing a weighted index:
\begin{quote}\begin{multicols}{2}
\noindent (A) Laspeyre's method \\
(B) Paasche's method \\
(C) Fisher's method \\
(D) None of these
\end{multicols}\end{quote}

\item Which of the following indices satisfies both the time reversal and factor reversal tests:
\begin{quote}\begin{multicols}{2}
\noindent (A) Laspeyre's index \\
(B) Paasche's index \\
(C) Fisher's index \\
(D) Marshall - Edgeworth
\end{multicols}\end{quote}

\item Circular test is satisfied by:
\begin{quote}\begin{multicols}{2}
\noindent (A) Laspeyre's index \\
(B) Paasche's index \\
(C) Fisher's index \\
(D) None of these
\end{multicols}\end{quote}

\item To measure how much the cost of some variable changes over time, you would use:
\begin{quote}\begin{multicols}{2}
\noindent (A) A value index \\
(B) An inflation index \\
(C) A quantity index \\
(D) None of these
\end{multicols}\end{quote}

\item Commodities subject to considerable price variation should best be measured by:
\begin{quote}\begin{multicols}{2}
\noindent (A) Quantity index \\
(B) Price index \\
(C) Value index \\
(D) None of these
\end{multicols}\end{quote}

\item Theoretically best average used in the construction of composite index is:
\begin{quote}\begin{multicols}{2}
\noindent (A) The arithmetic mean \\
(B) The geometric mean \\
(C) The median \\
(D) The harmonic mean
\end{multicols}\end{quote}

\item Which of the following indices has an upward bias:
\begin{quote}\begin{multicols}{2}
\noindent (A) Laspeyre's index \\
(B) Paasche's index \\
(C) Fisher's index \\
(D) None of these
\end{multicols}\end{quote}

\item Paasche's Index Number:
\begin{quote}\begin{multicols}{2}
\noindent (A) $\dfrac{\Sigma P_nq_o}{\Sigma P_nq_o}\times 100$ \\
(B) $\dfrac{\Sigma P_nq_o}{\Sigma P_oq_n}\times 100$ \\
(C) $\dfrac{\Sigma P_nq_o}{\Sigma P_oq_o}$ \\
(D) $\dfrac{\Sigma P_nq_n}{\Sigma P_oq_n}\times 100$
\end{multicols}\end{quote}

\item Which is the most suitable average in chain base method:
\begin{quote}\begin{multicols}{2}
\noindent (A) Arithmetic mean \\
(B) Median \\
(C) Geometric mean \\
(D) Weighted arithmetic mean
\end{multicols}\end{quote}

\item The prices used in the construction of consumer price index numbers are:
\begin{quote}\begin{multicols}{2}
\noindent (A) The retail prices \\
(B) The wholesale prices \\
(C) The fixed prices \\
(D) None of these
\end{multicols}\end{quote}

\item The index given by $\dfrac{\Sigma p_nq_n}{\Sigma p_oq_o}\times 100$ is:
\begin{quote}\begin{multicols}{2}
\noindent (A) The Laspeyre's index \\
(B) The Paasche's index \\
(C) The value index \\
(D) None of these
\end{multicols}\end{quote}

\item The consumer price index number is also called:
\begin{quote}\begin{multicols}{2}
\noindent (A) The cost of living index number \\
(B) The retail price index number \\
(C) The wholesale price index number \\
(D) Both
\end{multicols}\end{quote}

\item Which method of construction of consumer price index number is the Laspeyre's index number:
\begin{quote}\begin{multicols}{2}
\noindent (A) Aggregate expenditure method \\
(B) Family budget method \\
(C) Both \\
(D) None of these
\end{multicols}\end{quote}

\item Which of the following price indices are prepared by the Federal Bureau of Statistics:
\begin{quote}\begin{multicols}{2}
\noindent (A) Wholesale price index \\
(B) Consumer price index \\
(C) Sensitive price indicator \\
(D) All of these
\end{multicols}\end{quote}

\item Which index number has a wide scope:
\begin{quote}\begin{multicols}{2}
\noindent (A) Special \\
(B) General \\
(C) Price \\
(D) Quantity
\end{multicols}\end{quote}

\item In chain base method, the base period is:
\begin{quote}\begin{multicols}{2}
\noindent (A) Fixed \\
(B) Changed \\
(C) Constant \\
(D) None of these
\end{multicols}\end{quote}

\item Base year weighted index numbers are:
\begin{quote}\begin{multicols}{2}
\noindent (A) Laspeyre's \\
(B) Paasche's \\
(C) Fisher's \\
(D) C.P.I.
\end{multicols}\end{quote}

\item Index for base period is always taken as:
\begin{quote}\begin{multicols}{2}
\noindent (A) 100 \\
(B) One \\
(C) 200 \\
(D) Zero
\end{multicols}\end{quote}

\item Which formula is used in chain indices:
\begin{quote}\begin{multicols}{2}
\noindent (A) $\dfrac{\Sigma P_n}{\Sigma P_o}\times 100$ \\
(B) $\dfrac{P_n}{\Sigma P_{n-1}}\times 100$ \\
(C) $\dfrac{P_n}{P_{n-1}}\times 100$ \\
(D) $\dfrac{\Sigma P_nq_n}{\Sigma P_oq_o}$
\end{multicols}\end{quote}

\item Link relatives can be obtained by dividing $P_n$ by:
\begin{quote}\begin{multicols}{2}
\noindent (A) $P_o$ \\
(B) $P_{n-1}$ \\
(C) $q_o$ \\
(D) $q_{n-1}$
\end{multicols}\end{quote}

\item If Laspeyre's index number is 200, Paasche's index number is 200, the Fisher's index number is:
\begin{quote}\begin{multicols}{2}
\noindent (A) 100 \\
(B) 200 \\
(C) Zero \\
(D) 1000
\end{multicols}\end{quote}

\item Price relative is a ratio of current year price and:
\begin{quote}\begin{multicols}{2}
\noindent (A) Base year quantity \\
(B) Previous year quantity \\
(C) Base year price \\
(D) Current year quantity
\end{multicols}\end{quote}

\item An index that measures the change for a fixed time period is called:
\begin{quote}\begin{multicols}{2}
\noindent (A) Chain base method \\
(B) Fixed base method \\
(C) Simple aggregative method \\
(D) Cost of living method
\end{multicols}\end{quote}

\item How many basic types of index numbers are there:
\begin{quote}\begin{multicols}{2}
\noindent (A) 2 \\
(B) 3 \\
(C) 4 \\
(D) 5
\end{multicols}\end{quote}

\item In fixed base method, the base period should be:
\begin{quote}\begin{multicols}{2}
\noindent (A) Far away \\
(B) Normal \\
(C) Unreliable \\
(D) Abnormal
\end{multicols}\end{quote}

\item In a fixed base method, which period is always taken as 100:
\begin{quote}\begin{multicols}{2}
\noindent (A) Preceding \\
(B) Following \\
(C) Base \\
(D) Current
\end{multicols}\end{quote}

\item When all the commodities are not of equal importance, the index numbers are called:
\begin{quote}\begin{multicols}{2}
\noindent (A) Simple \\
(B) Un-weighted \\
(C) Weighted \\
(D) Value
\end{multicols}\end{quote}

\item $\dfrac{\Sigma p_n}{\Sigma p_o}\times 100$ is called:
\begin{quote}\begin{multicols}{2}
\noindent (A) Simple index \\
(B) Weighted index \\
(C) Simple aggregative index \\
(D) Price relative
\end{multicols}\end{quote}

\item CPI falls in the category of:
\begin{quote}\begin{multicols}{2}
\noindent (A) A simple index \\
(B) An aggregative index \\
(C) An inflationary index \\
(D) Wholesale price index
\end{multicols}\end{quote}

\item Which of the following formulas is Laspeyre's price index number:
\begin{quote}\begin{multicols}{2}
\noindent (A) $\dfrac{\Sigma p_nq_0}{\Sigma p_0q_0}\times 100$ \\
(B) $\dfrac{\Sigma p_0q_n}{\Sigma p_0q_0}\times 100$ \\
(C) $\dfrac{\Sigma p_nq_n}{\Sigma p_0q_0}\times 100$ \\
(D) All of these
\end{multicols}\end{quote}

\item For computing chain index, we compute:
\begin{quote}\begin{multicols}{2}
\noindent (A) Price relative \\
(B) Link relative \\
(C) Weighted indices \\
(D) None of these
\end{multicols}\end{quote}

\item Consumer price index is obtained by:
\begin{quote}\begin{multicols}{2}
\noindent (A) Paasche's formula \\
(B) Fisher's ideal formula \\
(C) Marshall Edgeworth formula \\
(D) Family Budget Method formula
\end{multicols}\end{quote}

\item Composite index number involves commodities:
\begin{quote}\begin{multicols}{2}
\noindent (A) One \\
(B) Two \\
(C) Three \\
(D) More than one
\end{multicols}\end{quote}

\item An arrangement of all or some of a set of objects in a definite order is called:
\begin{quote}\begin{multicols}{2}
\noindent (A) Factorial \\
(B) Combinations \\
(C) Permutations \\
(D) None of these
\end{multicols}\end{quote}

\item A set is usually denoted by a:
\begin{quote}\begin{multicols}{2}
\noindent (A) Small letter \\
(B) Capital letter \\
(C) Both a and b \\
(D) None of these
\end{multicols}\end{quote}

\item Two sets $A$ and $B$ are called overlapping sets if they have some:
\begin{quote}\begin{multicols}{2}
\noindent (A) Different elements \\
(B) Common elements \\
(C) Negative element \\
(D) None of these
\end{multicols}\end{quote}

\item The product set of $A$ and $B$ consists of all ordered pairs denoted by:
\begin{quote}\begin{multicols}{2}
\noindent (A) $A\cdot B$ \\
(B) $B\cdot A$ \\
(C) $A\times B$ \\
(D) None of these
\end{multicols}\end{quote}

\item The complement of set $A$ is:
\begin{quote}\begin{multicols}{2}
\noindent (A) $U-A'$ \\
(B) $U-A$ \\
(C) $A'-U$ \\
(D) None of these
\end{multicols}\end{quote}

\item How many possible permutations can be formed from the word ``PAKISTAN'':
\begin{quote}\begin{multicols}{2}
\noindent (A) $8!$ \\
(B) $20160$ \\
(C) $6!$ \\
(D) None of these
\end{multicols}\end{quote}

\item An arrangement of all objects without caring for the order is called:
\begin{quote}\begin{multicols}{2}
\noindent (A) Permutation \\
(B) Binomial Expansion \\
(C) Combination \\
(D) None of these
\end{multicols}\end{quote}

\item ${}^{8}C_3=$\underline{\hspace{1cm}}:
\begin{quote}\begin{multicols}{2}
\noindent (A) 20 \\
(B) 56 \\
(C) 6720 \\
(D) None of these
\end{multicols}\end{quote}

\item In the expansion of a binomial, the coefficient of the first and the last term is always:
\begin{quote}\begin{multicols}{2}
\noindent (A) 2 \\
(B) 3 \\
(C) 1 \\
(D) None of these
\end{multicols}\end{quote}

\item The number of terms in the expansion of the binomial $(p+q)^n$ is:
\begin{quote}\begin{multicols}{2}
\noindent (A) $n$ \\
(B) $n+1$ \\
(C) $n+2$ \\
(D) None of these
\end{multicols}\end{quote}

\item The difference between $A$ and $B$ is shown by:
\begin{quote}\begin{multicols}{2}
\noindent (A) $A-B$ \\
(B) $B-A$ \\
(C) $A/B$ \\
(D) None of these
\end{multicols}\end{quote}

\item ${}^{n}C_r=$\underline{\hspace{1cm}}:
\begin{quote}\begin{multicols}{2}
\noindent (A) ${}^{n-1}C_r$ \\
(B) ${}^{n}C_{n-r}$ \\
(C) ${}^{n}C_{n-r-1}$ \\
(D) None of these
\end{multicols}\end{quote}

\item The product of $n$ natural numbers is read as:
\begin{quote}\begin{multicols}{2}
\noindent (A) $n^n$ \\
(B) $n!$ \\
(C) $\Sigma n$ \\
(D) None of these
\end{multicols}\end{quote}

\item $\dfrac{12!}{9!}=$\underline{\hspace{1cm}}:
\begin{quote}\begin{multicols}{2}
\noindent (A) 336 \\
(B) 1320 \\
(C) 3360 \\
(D) None of these
\end{multicols}\end{quote}

\item For $r>n$, ${}^{n}C_r=$\underline{\hspace{1cm}}:
\begin{quote}\begin{multicols}{2}
\noindent (A) $n$ \\
(B) 1 \\
(C) 0 \\
(D) Not possible
\end{multicols}\end{quote}

\item A set can be specified in two ways:
\begin{quote}\begin{multicols}{2}
\noindent (A) Tabular and empty set form \\
(B) Tabular and set builder form \\
(C) Union and set builder form \\
(D) None of these
\end{multicols}\end{quote}

\item A 300-year-old person is an example of a:
\begin{quote}\begin{multicols}{2}
\noindent (A) Difference set \\
(B) Null set \\
(C) Subset \\
(D) None of these
\end{multicols}\end{quote}

\item $B$ is a proper subset of $A$, if:
\begin{quote}\begin{multicols}{2}
\noindent (A) $A=B$ \\
(B) $A\subset B$ \\
(C) $A\ne B$ \\
(D) None of these
\end{multicols}\end{quote}

\item In how many ways can 4 seats on a sofa set be occupied by 7 persons:
\begin{quote}\begin{multicols}{2}
\noindent (A) 60 \\
(B) 840 \\
(C) 28 \\
(D) None of these
\end{multicols}\end{quote}

\item The set of all subsets of a set is called:
\begin{quote}\begin{multicols}{2}
\noindent (A) The power set \\
(B) Null set \\
(C) Product set \\
(D) None of these
\end{multicols}\end{quote}

\item If $A\subset B$ and $B\subset C$, then $A$ and $B$ are:
\begin{quote}\begin{multicols}{2}
\noindent (A) Difference \\
(B) Equal \\
(C) Not equal \\
(D) None of these
\end{multicols}\end{quote}

\item ${}^{n}C_r=$\underline{\hspace{1cm}}:
\begin{quote}\begin{multicols}{2}
\noindent (A) $r! \times {}^{n}P_r$ \\
(B) $\dfrac{1}{r!}\times {}^{n}P_r$ \\
(C) ${}^{n}P_{r-1}$ \\
(D) None of these
\end{multicols}\end{quote}

\item The total number of subsets of a set containing $n$ elements is:
\begin{quote}\begin{multicols}{2}
\noindent (A) $2n$ \\
(B) $2^n$ \\
(C) $n^2$ \\
(D) None of these
\end{multicols}\end{quote}

\item A universal set is usually denoted by:
\begin{quote}\begin{multicols}{2}
\noindent (A) $A$ \\
(B) $S$ \\
(C) $U$ \\
(D) None of these
\end{multicols}\end{quote}

\item ${}^{8}P_6=$\underline{\hspace{1cm}}:
\begin{quote}\begin{multicols}{2}
\noindent (A) $\dfrac{8!}{(8-6)!}$ \\
(B) $\dfrac{8!}{(8-6)!\,6!}$ \\
(C) $\dfrac{8!}{6!}$ \\
(D) None of these
\end{multicols}\end{quote}

\item ${}^{n}C_{n-1}=$\underline{\hspace{1cm}}:
\begin{quote}\begin{multicols}{2}
\noindent (A) $n$ \\
(B) $r$ \\
(C) $n$ \\
(D) None of these
\end{multicols}\end{quote}

\item The permutation of the word ``ARRANGEMENT'' is:
\begin{quote}\begin{multicols}{2}
\noindent (A) 2449800 \\
(B) 248900 \\
(C) 2494800 \\
(D) None of these
\end{multicols}\end{quote}

\item The permutation of the word ``PAKPATTAN'' is:
\begin{quote}\begin{multicols}{2}
\noindent (A) 11520 \\
(B) 15120 \\
(C) 15210 \\
(D) None of these
\end{multicols}\end{quote}

\item The permutation of the word ``STATISTICS'' is:
\begin{quote}\begin{multicols}{2}
\noindent (A) 500000 \\
(B) 50400 \\
(C) 54000 \\
(D) None of these
\end{multicols}\end{quote}

\item Four books can be arranged on a shelf in:
\begin{quote}\begin{multicols}{2}
\noindent (A) 8 ways \\
(B) 12 ways \\
(C) 16 ways \\
(D) 24 ways
\end{multicols}\end{quote}

\item If one event is unaffected by the outcome of another event, the two events are said to be:
\begin{quote}\begin{multicols}{2}
\noindent (A) Dependent \\
(B) Independent \\
(C) Mutually exclusive \\
(D) Both
\end{multicols}\end{quote}

\item Symbolically, a marginal probability is:
\begin{quote}\begin{multicols}{2}
\noindent (A) $P\{AB\}$ \\
(B) $P\{A\}$ \\
(C) $P(A\mid B)$ \\
(D) Both a and b
\end{multicols}\end{quote}

\item The simple probability of occurrence of an event is called the:
\begin{quote}\begin{multicols}{2}
\noindent (A) Joint probability \\
(B) Conditional probability \\
(C) Marginal probability \\
(D) Subjective probability
\end{multicols}\end{quote}

\item Why are the outcomes of a coin toss mutually exclusive:
\begin{quote}\begin{multicols}{2}
\noindent (A) The outcome of any toss is not affected by the outcome of those preceding it \\
(B) Both a head and a tail cannot turn up on any one toss \\
(C) The probability of getting a head and the probability of getting a tail is the same \\
(D) All of these
\end{multicols}\end{quote}

\item What is the probability that a value chosen at random from a particular population is larger than the median of the population:
\begin{quote}\begin{multicols}{2}
\noindent (A) 0.25 \\
(B) 0.5 \\
(C) 1.0 \\
(D) 0.67
\end{multicols}\end{quote}

\item What is the probability that a ball drawn at random from the bag is (as specified):
\begin{quote}\begin{multicols}{2}
\noindent (A) 0.1 \\
(B) 0.4 \\
(C) 1.0 \\
(D) Cannot be determined from given information
\end{multicols}\end{quote}

\item When two dice are rolled, the number of possible sample points is:
\begin{quote}\begin{multicols}{2}
\noindent (A) 6 \\
(B) 12 \\
(C) 36 \\
(D) 48
\end{multicols}\end{quote}

\item The probability of drawing a king of spades from a pack of 52 cards is:
\begin{quote}\begin{multicols}{2}
\noindent (A) $\dfrac{1}{4}$ \\
(B) $\dfrac{1}{13}$ \\
(C) $\dfrac{1}{26}$ \\
(D) $\dfrac{1}{52}$
\end{multicols}\end{quote}

\item If $P(A \text{ or } B)=P(A)$, then:
\begin{quote}\begin{multicols}{2}
\noindent (A) $P(A)+P(B)$ \\
(B) $P(A\cap B)$ is the joint probability of $A$ and $B$ \\
(C) $P(A\cap B)=0$ \\
(D) None of these
\end{multicols}\end{quote}

\item If a Venn diagram is drawn for events $A$ and $B$ which are mutually exclusive, which of the following would always be true of $A$ and $B$:
\begin{quote}\begin{multicols}{2}
\noindent (A) Their parts of the rectangle will overlap \\
(B) Their parts of the rectangle will be equal in area \\
(C) Their parts of the rectangle will not overlap \\
(D) None of these
\end{multicols}\end{quote}

\item When two coins are tossed simultaneously, $P(\text{one head})$ is:
\begin{quote}\begin{multicols}{2}
\noindent (A) $\dfrac{1}{4}$ \\
(B) $\dfrac{1}{2}$ \\
(C) $\dfrac{3}{4}$ \\
(D) 1.0
\end{multicols}\end{quote}

\item When three coins are tossed simultaneously, $P(3\text{ heads})$ is:
\begin{quote}\begin{multicols}{2}
\noindent (A) $\dfrac{3}{8}$ \\
(B) $\dfrac{1}{2}$ \\
(C) $\dfrac{1}{4}$ \\
(D) $\dfrac{1}{8}$
\end{multicols}\end{quote}

\item The probability of drawing two aces from a pack of 52 cards with replacement is:
\begin{quote}\begin{multicols}{2}
\noindent (A) $\dfrac{1}{4}$ \\
(B) $\dfrac{1}{16}$ \\
(C) $\dfrac{1}{169}$ \\
(D) $\dfrac{1}{256}$
\end{multicols}\end{quote}

\item For two mutually exclusive events $A$ and $B$, $P(A)=0.2$ and $P(B)=0.4$, then $P(A\cup B)$ is:
\begin{quote}\begin{multicols}{2}
\noindent (A) 0.8 \\
(B) 0.2 \\
(C) 0.6 \\
(D) 0.5
\end{multicols}\end{quote}

\item For two independent events $A$ and $B$, $P(A)=0.6$ and $P(B)=0.3$, then $P(A\cap B)$ is:
\begin{quote}\begin{multicols}{2}
\noindent (A) 0.9 \\
(B) 0.12 \\
(C) 0.5 \\
(D) 0.18
\end{multicols}\end{quote}

\item If $A$ and $B$ are mutually exclusive events, then $P(A\cup B)=P(A)+P(B)$. How does the calculation of $P(A\cup B)$ change if $A$ and $B$ are not mutually exclusive:
\begin{quote}\begin{multicols}{2}
\noindent (A) $P(A)+P(B)$ must be divided by $P(A\cap B)$ \\
(B) $[P(A)+P(B)]-P(A\cap B)$ \\
(C) $[P(A)+P(B)]+P(A\cap B)$ \\
(D) No change
\end{multicols}\end{quote}

\item Assume that a single fair die is rolled once. Which of the following is true:
\begin{quote}\begin{multicols}{2}
\noindent (A) The probability of rolling a number greater than `1' is $1-P\{\text{`1' is rolled}\}$ \\
(B) The probability of rolling a `3' is $1-P\{1,2,4,5\text{ or }6\text{ is rolled}\}$ \\
(C) The probability of rolling a `5 or 6' is higher than the probability of rolling a `3 or 4' \\
(D) Both a and b
\end{multicols}\end{quote}

\item If $P(B\mid A)=0.25$ and $P(A\cap B)=0.20$, then $P(A)$ is:
\begin{quote}\begin{multicols}{2}
\noindent (A) 0.05 \\
(B) 0.8 \\
(C) 0.95 \\
(D) 0.75
\end{multicols}\end{quote}

\item The probability of a red card out of 52 cards is:
\begin{quote}\begin{multicols}{2}
\noindent (A) $\dfrac{1}{4}$ \\
(B) $\dfrac{1}{2}$ \\
(C) $\dfrac{4}{52}$ \\
(D) Zero
\end{multicols}\end{quote}

\item When two coins are tossed simultaneously, the probability of at most one head is:
\begin{quote}\begin{multicols}{2}
\noindent (A) $\dfrac{1}{4}$ \\
(B) $\dfrac{1}{2}$ \\
(C) $\dfrac{3}{4}$ \\
(D) None of these
\end{multicols}\end{quote}

\item A letter is chosen at random from the word STATISTICS. The probability of getting a vowel is:
\begin{quote}\begin{multicols}{2}
\noindent (A) $\dfrac{1}{5}$ \\
(B) $\dfrac{3}{10}$ \\
(C) $\dfrac{1}{2}$ \\
(D) $\dfrac{2}{5}$
\end{multicols}\end{quote}

\item The probability of getting one red ball from a bag containing 4 red, 3 white and 3 black balls is:
\begin{quote}\begin{multicols}{2}
\noindent (A) $\dfrac{3}{10}$ \\
(B) $\dfrac{1}{5}$ \\
(C) $\dfrac{2}{5}$ \\
(D) $\dfrac{1}{2}$
\end{multicols}\end{quote}

\item $P(A\cap B)=P(A\mid B)P(B)$, then $A$ and $B$ are:
\begin{quote}\begin{multicols}{2}
\noindent (A) Independent \\
(B) Mutually exclusive \\
(C) Dependent \\
(D) None of these
\end{multicols}\end{quote}

\item The probability of getting two red balls with replacement from a bag containing 4 red, 3 white and 3 black balls is:
\begin{quote}\begin{multicols}{2}
\noindent (A) $\dfrac{4}{25}$ \\
(B) $\dfrac{1}{25}$ \\
(C) $\dfrac{9}{100}$ \\
(D) $\dfrac{2}{25}$
\end{multicols}\end{quote}

\item Ten numbered balls are placed in an urn. Numbers 1--4 are red and numbers 5--10 are blue. The probability that a ball drawn at random from the urn is blue is:
\begin{quote}\begin{multicols}{2}
\noindent (A) 0.1 \\
(B) 0.4 \\
(C) 0.6 \\
(D) 1.0
\end{multicols}\end{quote}

\item For two independent events $A$ and $B$, $P(A)=0.6$ and $P(B)=0.3$. Then $P(A\cup B)$ is:
\begin{quote}\begin{multicols}{2}
\noindent (A) 0.9 \\
(B) 0.18 \\
(C) 0.72 \\
(D) None of these
\end{multicols}\end{quote}

\item Which of the following approaches to probability would you apply to solve the problems given afterwards:
\begin{quote}\begin{multicols}{2}
\noindent (A) Classical approach \\
(B) Relative Frequency approach \\
(C) Subjective approach \\
(D) Axiomatic approach
\end{multicols}\end{quote}

\item The probability of getting an odd number when a balanced die is rolled is:
\begin{quote}\begin{multicols}{2}
\noindent (A) $\dfrac{1}{6}$ \\
(B) $\dfrac{1}{4}$ \\
(C) $\dfrac{1}{2}$ \\
(D) $\dfrac{1}{3}$
\end{multicols}\end{quote}

\item If $P(B\mid A)=P(B)$, then $A$ and $B$ are:
\begin{quote}\begin{multicols}{2}
\noindent (A) Mutually exclusive \\
(B) Dependent \\
(C) Independent \\
(D) None of these
\end{multicols}\end{quote}

\item If $P(A\cup B)=P(A)+P(B)$, then $A$ and $B$ are:
\begin{quote}\begin{multicols}{2}
\noindent (A) Mutually exclusive \\
(B) Independent \\
(C) Not mutually exclusive \\
(D) None of these
\end{multicols}\end{quote}

\item The probability of drawing any one spade card is:
\begin{quote}\begin{multicols}{2}
\noindent (A) $\dfrac{1}{32}$ \\
(B) $\dfrac{1}{13}$ \\
(C) $\dfrac{4}{13}$ \\
(D) $\dfrac{1}{4}$
\end{multicols}\end{quote}

\item In the tossing of two perfect coins, the probability that at least one head occurs is:
\begin{quote}\begin{multicols}{2}
\noindent (A) $\dfrac{1}{2}$ \\
(B) $\dfrac{1}{4}$ \\
(C) $\dfrac{3}{4}$ \\
(D) 1
\end{multicols}\end{quote}

\item The probability of an event cannot be:
\begin{quote}\begin{multicols}{2}
\noindent (A) $=0$ \\
(B) $>0$ \\
(C) $=1$ \\
(D) $<0$
\end{multicols}\end{quote}

\item If two coins are tossed, the probability of getting one head and one tail is:
\begin{quote}\begin{multicols}{2}
\noindent (A) $\dfrac{1}{4}$ \\
(B) $\dfrac{2}{4}$ \\
(C) $\dfrac{3}{4}$ \\
(D) $\dfrac{2}{3}$
\end{multicols}\end{quote}

\item A measure of the chance that an uncertain event will occur:
\begin{quote}\begin{multicols}{2}
\noindent (A) An experiment \\
(B) An event \\
(C) A probability \\
(D) A trial
\end{multicols}\end{quote}

\item When a die and a coin are rolled together, all possible outcomes are:
\begin{quote}\begin{multicols}{2}
\noindent (A) 6 \\
(B) 2 \\
(C) 36 \\
(D) 12
\end{multicols}\end{quote}

\item Two events $A$ and $B$ are called mutually exclusive if:
\begin{quote}\begin{multicols}{2}
\noindent (A) $A\cup B=\phi$ \\
(B) $A\cap B=\phi$ \\
(C) $A\cap B=S$ \\
(D) $A\cap B=1$
\end{multicols}\end{quote}

\item A set of numerical values assigned to a sample space is called:
\begin{quote}\begin{multicols}{2}
\noindent (A) Random sample \\
(B) Random variable \\
(C) Random numbers \\
(D) Random experiment
\end{multicols}\end{quote}

\item Events with equal probabilities are called:
\begin{quote}\begin{multicols}{2}
\noindent (A) Mutually exclusive events \\
(B) Exhaustive events \\
(C) Equally likely events \\
(D) Simple events
\end{multicols}\end{quote}

\item If $A$ and $B$ are mutually exclusive events, then $P(A\cap B)$ is:
\begin{quote}\begin{multicols}{2}
\noindent (A) Zero \\
(B) 1 \\
(C) 0.5 \\
(D) 0.8
\end{multicols}\end{quote}

\item The collection of all possible outcomes of a random experiment is called:
\begin{quote}\begin{multicols}{2}
\noindent (A) Sample point \\
(B) Sure event \\
(C) Sample space \\
(D) Simple event
\end{multicols}\end{quote}

\item A student solved 25 questions from the first 50 questions of a book to be solved. The probability that he will solve all the remaining questions is:
\begin{quote}\begin{multicols}{2}
\noindent (A) 0.25 \\
(B) 0.5 \\
(C) 1 \\
(D) 0
\end{multicols}\end{quote}

\item If $P(A)=\dfrac{1}{3}$, $P(B)=\dfrac{2}{3}$, and $A$ and $B$ are independent, then $P(A\cap B)$ is:
\begin{quote}\begin{multicols}{2}
\noindent (A) $\dfrac{2}{6}$ \\
(B) $\dfrac{2}{9}$ \\
(C) 0 \\
(D) $\dfrac{1}{9}$
\end{multicols}\end{quote}

\item The result of no interest in an experiment is called:
\begin{quote}\begin{multicols}{2}
\noindent (A) Constant \\
(B) Event \\
(C) Failure \\
(D) Success
\end{multicols}\end{quote}

\item Which is the impossible event when a die is rolled:
\begin{quote}\begin{multicols}{2}
\noindent (A) 5 or 6 \\
(B) 6 or 7 \\
(C) 2 or 3 \\
(D) 1
\end{multicols}\end{quote}

\item When two dice are rolled, the maximum total on the two faces of the dice will be:
\begin{quote}\begin{multicols}{2}
\noindent (A) 6 \\
(B) 36 \\
(C) 12 \\
(D) 2
\end{multicols}\end{quote}

\item The probability of an impossible event is:
\begin{quote}\begin{multicols}{2}
\noindent (A) Positive \\
(B) Zero \\
(C) Negative \\
(D) 1
\end{multicols}\end{quote}

\item Two coins are tossed; the probability of two tails is equal to:
\begin{quote}\begin{multicols}{2}
\noindent (A) $\dfrac{1}{4}$ \\
(B) $\dfrac{1}{2}$ \\
(C) $\dfrac{3}{4}$ \\
(D) 1
\end{multicols}\end{quote}

\item The range of probability is between:
\begin{quote}\begin{multicols}{2}
\noindent (A) 0 to 2 \\
(B) $-1$ to $+1$ \\
(C) 0 to $\infty$ \\
(D) 0 to 1
\end{multicols}\end{quote}

\item If $\Sigma(XY)=\Sigma(X)\Sigma(Y)$, then $X$ and $Y$ are:
\begin{quote}\begin{multicols}{2}
\noindent (A) Dependent \\
(B) Independent \\
(C) Biased \\
(D) Unbiased
\end{multicols}\end{quote}

\item Tossing two dice, the possible number of samples is:
\begin{quote}\begin{multicols}{2}
\noindent (A) 36 \\
(B) 12 \\
(C) 6 \\
(D) 2
\end{multicols}\end{quote}

\item The term ``event'' is used for:
\begin{quote}\begin{multicols}{2}
\noindent (A) Time \\
(B) Subset of the sample space \\
(C) Probability \\
(D) Total number of outcomes
\end{multicols}\end{quote}

\item Which of the following cannot be the probability of an event:
\begin{quote}\begin{multicols}{2}
\noindent (A) 0 \\
(B) 1 \\
(C) 0.32 \\
(D) 1.05
\end{multicols}\end{quote}

\item When a pair of dice is rolled, the sum of the uppermost dots varies from:
\begin{quote}\begin{multicols}{2}
\noindent (A) 0 to 10 \\
(B) 1 to 11 \\
(C) 2 to 19 \\
(D) 2 to 12
\end{multicols}\end{quote}

\item Every random experiment results in outcomes:
\begin{quote}\begin{multicols}{2}
\noindent (A) two or more \\
(B) three or more \\
(C) four or more \\
(D) five or more
\end{multicols}\end{quote}

\item When a pair of dice is thrown, the sum of upturned dots varies from:
\begin{quote}\begin{multicols}{2}
\noindent (A) 0 to 10 \\
(B) 2 to 12 \\
(C) 1 to 11 \\
(D) 3 to 13
\end{multicols}\end{quote}

\item Random variables are denoted by:
\begin{quote}\begin{multicols}{2}
\noindent (A) $L, M$ or $N$ \\
(B) $A, B$ or $C$ \\
(C) $X, Y$ or $Z$ \\
(D) $P, Q$ or $R$
\end{multicols}\end{quote}

\item The values of random variables are denoted by:
\begin{quote}\begin{multicols}{2}
\noindent (A) $p, q$ or $r$ \\
(B) $l, m$ or $n$ \\
(C) $a, b$ or $c$ \\
(D) $x, y$ or $z$
\end{multicols}\end{quote}

\item In a pair of dice, there are outcomes:
\begin{quote}\begin{multicols}{2}
\noindent (A) 32 \\
(B) 36 \\
(C) 38 \\
(D) 48
\end{multicols}\end{quote}

\item Each digit $\{0,1,2,\ldots,9\}$ has probability:
\begin{quote}\begin{multicols}{2}
\noindent (A) 0.4 \\
(B) 0.2 \\
(C) 0.1 \\
(D) 0.3
\end{multicols}\end{quote}

\item $\{0,1,2,\ldots,9\}$ are also called:
\begin{quote}\begin{multicols}{2}
\noindent (A) Random variables \\
(B) Random outcomes \\
(C) Random experiments \\
(D) Random digits
\end{multicols}\end{quote}

\item We can get random numbers from:
\begin{quote}\begin{multicols}{2}
\noindent (A) Games of chance \\
(B) Reading books \\
(C) Different people \\
(D) Reading tables
\end{multicols}\end{quote}

\item Tables of random numbers are widely:
\begin{quote}\begin{multicols}{2}
\noindent (A) Acceptable \\
(B) Available \\
(C) Readable \\
(D) Countable
\end{multicols}\end{quote}

\item In pseudo-random techniques, the digits will re-appear in the:
\begin{quote}\begin{multicols}{2}
\noindent (A) Cycle \\
(B) Manner \\
(C) Order \\
(D) Routine
\end{multicols}\end{quote}

\item Pseudo-random digits are not truly:
\begin{quote}\begin{multicols}{2}
\noindent (A) Constructed \\
(B) Possible \\
(C) Associated \\
(D) Random
\end{multicols}\end{quote}

\item Pseudo-random digits are completely:
\begin{quote}\begin{multicols}{2}
\noindent (A) Probabilistic \\
(B) Realistic \\
(C) Deterministic \\
(D) None of these
\end{multicols}\end{quote}

\item $ax_n+b \pmod{m}$ is equal to:
\begin{quote}\begin{multicols}{2}
\noindent (A) $x_{n+1}$ \\
(B) $x_{n-1}$ \\
(C) $x_n$ \\
(D) $x_{n+2}$
\end{multicols}\end{quote}

\item The seed $x_0$ (a starting number) is:
\begin{quote}\begin{multicols}{2}
\noindent (A) A whole number \\
(B) A natural number \\
(C) An integral number \\
(D) A rational number
\end{multicols}\end{quote}

\item The range of a pseudo-random digit is from:
\begin{quote}\begin{multicols}{2}
\noindent (A) $0$ to $m+1$ \\
(B) A natural number \\
(C) An integer \\
(D) A rational number
\end{multicols}\end{quote}

\item Random numbers are widely used in:
\begin{quote}\begin{multicols}{2}
\noindent (A) Analyzation \\
(B) Simulation \\
(C) Criticization \\
(D) Experimentation
\end{multicols}\end{quote}

\item Probability proportional to size is abbreviated as:
\begin{quote}\begin{multicols}{2}
\noindent (A) PPS \\
(B) SPP \\
(C) PPP \\
(D) PSP
\end{multicols}\end{quote}

\item The random numbers $0$ to $9$ are called:
\begin{quote}\begin{multicols}{2}
\noindent (A) $1$-digit random numbers \\
(B) $2$-digit random numbers \\
(C) $3$-digit random numbers \\
(D) $4$-digit random numbers
\end{multicols}\end{quote}

\item PPS is used when we have information about:
\begin{quote}\begin{multicols}{2}
\noindent (A) Random variables \\
(B) Probabilities \\
(C) Probability distribution \\
(D) Index numbers
\end{multicols}\end{quote}

\item $00$--$99$ are called \underline{\hspace{1.5cm}} random numbers:
\begin{quote}\begin{multicols}{2}
\noindent (A) $1$-digit \\
(B) $2$-digit \\
(C) $3$-digit \\
(D) $4$-digit
\end{multicols}\end{quote}

\item $000$--$999$ are called \underline{\hspace{1.5cm}} random numbers:
\begin{quote}\begin{multicols}{2}
\noindent (A) $1$-digit \\
(B) $2$-digit \\
(C) $3$-digit \\
(D) $4$-digit
\end{multicols}\end{quote}

\item If the set of values comprises isolated points on a real line, then the variable is called:
\begin{quote}\begin{multicols}{2}
\noindent (A) Discrete \\
(B) Continuous \\
(C) Given \\
(D) Random
\end{multicols}\end{quote}

\item A random variable may be discrete or:
\begin{quote}\begin{multicols}{2}
\noindent (A) Experimental \\
(B) Functional \\
(C) Given \\
(D) Continuous
\end{multicols}\end{quote}

\item A random variable may be discrete or:
\begin{quote}\begin{multicols}{2}
\noindent (A) Experimental \\
(B) Functional \\
(C) Given \\
(D) Continuous
\end{multicols}\end{quote}

\item A random variable is also called:
\begin{quote}\begin{multicols}{2}
\noindent (A) Discrete variable \\
(B) Continuous variable \\
(C) Chance variable \\
(D) Variable
\end{multicols}\end{quote}

\item If $X$ and $Y$ are random variables then $E(x+y)=E(x)+E(y)$ is equal to:
\begin{quote}\begin{multicols}{2}
\noindent (A) $E(x)+E(y)$ \\
(B) $E(x)-E(y)$ \\
(C) $X-E(y)$ \\
(D) $E(x)-Y$
\end{multicols}\end{quote}

\item A balanced die is rolled; the probability of getting an odd number is:
\begin{quote}\begin{multicols}{2}
\noindent (A) $1/2$ \\
(B) $1/4$ \\
(C) $1/6$ \\
(D) $1/3$
\end{multicols}\end{quote}

\item If two dice are rolled, the number of possible outcomes are:
\begin{quote}\begin{multicols}{2}
\noindent (A) $6$ \\
(B) $36$ \\
(C) $1$ \\
(D) Difficult to tell
\end{multicols}\end{quote}

\item A \underline{\hspace{1.5cm}} random variable can assume each and every value within an interval:
\begin{quote}\begin{multicols}{2}
\noindent (A) Continuous \\
(B) Qualitative \\
(C) Discrete \\
(D) None of these
\end{multicols}\end{quote}

\item A variable that can assume only some specific value within a given range is called:
\begin{quote}\begin{multicols}{2}
\noindent (A) Continuous r.v. \\
(B) Discrete r.v. \\
(C) Sample space \\
(D) Difficult to tell
\end{multicols}\end{quote}

\item Random numbers are a sequence of digits from the set:
\begin{quote}\begin{multicols}{2}
\noindent (A) $\{0,1,2,\ldots,9\}$ \\
(B) $\{1,2,3,\ldots,9\}$ \\
(C) $\{1,2,3,\ldots,10\}$ \\
(D) $\{2,4,6,\ldots,10\}$
\end{multicols}\end{quote}

\item A random variable which assumes a finite number of values is called:
\begin{quote}\begin{multicols}{2}
\noindent (A) Continuous random variable \\
(B) Deviated variable \\
(C) Discrete random variable \\
(D) All of these
\end{multicols}\end{quote}

\item A tabular arrangement of probabilities against each possible value of a discrete random variable is known as:
\begin{quote}\begin{multicols}{2}
\noindent (A) Distribution function \\
(B) Density function \\
(C) Probability density function \\
(D) None of these
\end{multicols}\end{quote}

\item A variable that is the result of a measuring process is called a:
\begin{quote}\begin{multicols}{2}
\noindent (A) Uniform variable \\
(B) Discrete variable \\
(C) Continuous variable \\
(D) All of these
\end{multicols}\end{quote}

\item Every random experiment has:
\begin{quote}\begin{multicols}{2}
\noindent (A) Only one outcome \\
(B) At least $2$ outcomes \\
(C) $0$ \\
(D) None of these
\end{multicols}\end{quote}

\item A discrete random variable takes only:
\begin{quote}\begin{multicols}{2}
\noindent (A) Infinite values \\
(B) Countable values \\
(C) Uncountable values \\
(D) None of these
\end{multicols}\end{quote}

\item The total area under the density curve is:
\begin{quote}\begin{multicols}{2}
\noindent (A) $0$ \\
(B) $1$ \\
(C) $2$ \\
(D) None of these
\end{multicols}\end{quote}

\item If $x$ is a continuous random variable, then the probability that $x$ takes a specific value is:
\begin{quote}\begin{multicols}{2}
\noindent (A) $1$ \\
(B) $0$ \\
(C) $2^2$ \\
(D) None of these
\end{multicols}\end{quote}

\item A random variable may not assume:
\begin{quote}\begin{multicols}{2}
\noindent (A) Positive values \\
(B) Negative values \\
(C) Zero \\
(D) None of these
\end{multicols}\end{quote}

\item A discrete probability function $f(x)$ is always:
\begin{quote}\begin{multicols}{2}
\noindent (A) Non-negative \\
(B) Negative \\
(C) $0$ \\
(D) None of these
\end{multicols}\end{quote}

\item A random variable may be discrete or:
\begin{quote}\begin{multicols}{2}
\noindent (A) Experimental \\
(B) Functional \\
(C) Given \\
(D) Continuous
\end{multicols}\end{quote}

\item A discrete probability distribution may be represented by:
\begin{quote}\begin{multicols}{2}
\noindent (A) A table \\
(B) A graph \\
(C) A mathematical equation \\
(D) All of these
\end{multicols}\end{quote}

\item A continuous probability distribution may be represented by:
\begin{quote}\begin{multicols}{2}
\noindent (A) A table \\
(B) A graph \\
(C) A mathematical equation \\
(D) Both
\end{multicols}\end{quote}

\item If the random variable $X$ denotes the number of heads when three distinct coins are tossed, then $X$ assumes the values:
\begin{quote}\begin{multicols}{2}
\noindent (A) $0,1,2,3$ \\
(B) $1,3,3,1$ \\
(C) $1,2,3$ \\
(D) None of these
\end{multicols}\end{quote}

\item If $X$ is a discrete random variable, then the function $f(x)$ is:
\begin{quote}\begin{multicols}{2}
\noindent (A) A probability function \\
(B) A probability density function \\
(C) A density function \\
(D) A distribution function
\end{multicols}\end{quote}

\item If $X$ is a continuous random variable, then the function $f(x)$ is:
\begin{quote}\begin{multicols}{2}
\noindent (A) A probability function \\
(B) A probability density function \\
(C) A density function \\
(D) Both
\end{multicols}\end{quote}

\item The distribution function $F(x)$ is equal to:
\begin{quote}\begin{multicols}{2}
\noindent (A) $P$ \\
(B) $P(X\leq x)$ \\
(C) $P(X\geq x)$ \\
(D) None of these
\end{multicols}\end{quote}

\item The expected value of a discrete random variable is:
\begin{quote}\begin{multicols}{2}
\noindent (A) Always an integer \\
(B) Always one of the values that the random variable can assume \\
(C) An interval of values \\
(D) None of these
\end{multicols}\end{quote}

\item If $X$ and $Y$ are random variables, then $E(X-Y)$ is equal to:
\begin{quote}\begin{multicols}{2}
\noindent (A) $E\{X\}+E\{Y\}$ \\
(B) $E\{X\}-E\{Y\}$ \\
(C) $X-E\{Y\}$ \\
(D) $E\{X\}-Y$
\end{multicols}\end{quote}

\item If $X$ and $Y$ are independent random variables, then $\mathrm{Var}(X-Y)$ is equal to:
\begin{quote}\begin{multicols}{2}
\noindent (A) $\mathrm{Var}\{X\}-\mathrm{Var}\{Y\}$ \\
(B) $\mathrm{Var}\{X\}+\mathrm{Var}\{Y\}$ \\
(C) $E\{X-Y\}$ \\
(D) $E\{X+Y\}$
\end{multicols}\end{quote}

\item If $X$ and $Y$ are independent random variables, then $\mathrm{S.D.}(X-Y)$ is equal to:
\begin{quote}\begin{multicols}{2}
\noindent (A) $\sqrt{\mathrm{Var}(X)+\mathrm{Var}(Y)}$ \\
(B) $\sqrt{\mathrm{Var}(X)-\mathrm{Var}(Y)}$ \\
(C) $\mathrm{S.D.}\{X\}+\mathrm{S.D.}\{Y\}$ \\
(D) $\mathrm{S.D.}\{X\}-\mathrm{S.D.}\{Y\}$
\end{multicols}\end{quote}

\item If $X$ and $Y$ are independent random variables, then $E(XY)$ is equal to:
\begin{quote}\begin{multicols}{2}
\noindent (A) $E\{XY\}$ \\
(B) $E\{X\}E\{Y\}$ \\
(C) $XE\{Y\}$ \\
(D) $YE\{X\}$
\end{multicols}\end{quote}

\item If $X$ is a random variable and $a$ and $b$ are constants, then $\mathrm{Var}(a+bX)$ is equal to:
\begin{quote}\begin{multicols}{2}
\noindent (A) $\mathrm{Var}\{a\}+\mathrm{Var}\{bX\}$ \\
(B) $\mathrm{Var}\{a\}+b\,\mathrm{Var}\{X\}$ \\
(C) $b^2\mathrm{Var}\{X\}$ \\
(D) $b\,\mathrm{Var}\{X\}$
\end{multicols}\end{quote}

\item If $X$ is a random variable and $a$ and $b$ are constants, then $\mathrm{S.D.}(a-bX)$ is equal to:
\begin{quote}\begin{multicols}{2}
\noindent (A) $\mathrm{S.D.}\{a\}-\mathrm{S.D.}\{bX\}$ \\
(B) $\mathrm{S.D.}\{a\}-b\,\mathrm{S.D.}\{X\}$ \\
(C) $b\sqrt{\mathrm{Var}(X)}$ \\
(D) None of these
\end{multicols}\end{quote}

\item If $X$ and $Y$ are random variables and $a$ and $b$ are constants, then $E(aX+bY)$ is equal to:
\begin{quote}\begin{multicols}{2}
\noindent (A) $E\{aX\}+E\{bX\}$ \\
(B) $aE\{X\}+bE\{Y\}$ \\
(C) $\{a+b\}\{E\{X\}+E\{Y\}\}$ \\
(D) None of these
\end{multicols}\end{quote}

\item If $E(X)=4$, then the arithmetic mean will be:
\begin{quote}\begin{multicols}{2}
\noindent (A) $4$ \\
(B) $8$ \\
(C) $0$ \\
(D) $1$
\end{multicols}\end{quote}

\item If $C$ is a non-random variable, then $E(C)$ is:
\begin{quote}\begin{multicols}{2}
\noindent (A) $c$ \\
(B) $0$ \\
(C) $1$ \\
(D) $x$
\end{multicols}\end{quote}

\item Given $\mathrm{Var}(X)=3$ and $\mathrm{Var}(Y)=5$. If $X$ and $Y$ are independent variables then $\mathrm{Var}(X-Y)=$:
\begin{quote}\begin{multicols}{2}
\noindent (A) $-2$ \\
(B) $8$ \\
(C) $24$ \\
(D) $2$
\end{multicols}\end{quote}

\item $E(X-\mu)^2 = $\underline{\hspace{1.5cm}}:
\begin{quote}\begin{multicols}{2}
\noindent (A) Zero \\
(B) S.D. \\
(C) Variance \\
(D) Mean deviation
\end{multicols}\end{quote}

\item An expected value of a random variable is equal to its:
\begin{quote}\begin{multicols}{2}
\noindent (A) Mean \\
(B) Variance \\
(C) S.D. \\
(D) Co-variance
\end{multicols}\end{quote}

\item If $y=5x+10$, then $\mathrm{S.D.}(y)$ is:
\begin{quote}\begin{multicols}{2}
\noindent (A) $5\,\mathrm{S.D.}\{x\}$ \\
(B) $25\,\mathrm{S.D.}\{x\}$ \\
(C) $10\,\mathrm{S.D.}\{x\}$ \\
(D) $5\,\mathrm{S.D.}\{x\}+\mathrm{S.D.}\{10\}$
\end{multicols}\end{quote}

\item If $X$ is a random variable, then $E(aX+b)$ is equal to:
\begin{quote}\begin{multicols}{2}
\noindent (A) $E\{X\}$ \\
(B) $aE\{X\}$ \\
(C) $aE\{X\}+b$ \\
(D) $E\{X\}+a$
\end{multicols}\end{quote}

\item If $Y_i=aX_i+b$, where $a$ and $b$ are any two constants, then $\mathrm{Var}(y)=$:
\begin{quote}\begin{multicols}{2}
\noindent (A) $a\,\mathrm{Var}\{X\}+b$ \\
(B) $\mathrm{Var}\{X\}+b$ \\
(C) $a\,\mathrm{Var}\{X\}$ \\
(D) $a^2\mathrm{Var}\{X\}$
\end{multicols}\end{quote}

\item For a random variable $X$, $\Sigma P(x)$ is equal to:
\begin{quote}\begin{multicols}{2}
\noindent (A) Zero \\
(B) $1$ \\
(C) Greater than $1$ \\
(D) Less than $1$
\end{multicols}\end{quote}

\item The probability that a continuous random variable $X$ takes on a specific value of $x$ is:
\begin{quote}\begin{multicols}{2}
\noindent (A) Greater than zero \\
(B) Less than zero \\
(C) Equal to zero \\
(D) $0$ to $1$
\end{multicols}\end{quote}

\item A probability density function is the probability function of a \underline{\hspace{1.5cm}} random variable:
\begin{quote}\begin{multicols}{2}
\noindent (A) Continuous \\
(B) Discrete \\
(C) Qualitative \\
(D) None
\end{multicols}\end{quote}

\item The probability density function $P(x)$ cannot exceed:
\begin{quote}\begin{multicols}{2}
\noindent (A) Zero \\
(B) One \\
(C) Mean \\
(D) Infinity
\end{multicols}\end{quote}

\item If $f(x)=\dfrac{1}{10}$, $X=10$, then $E(X)$ is:
\begin{quote}\begin{multicols}{2}
\noindent (A) Zero \\
(B) $6/8$ \\
(C) $1$ \\
(D) $-1$
\end{multicols}\end{quote}

\item Which of the following is not possible in a probability distribution:
\begin{quote}\begin{multicols}{2}
\noindent (A) $f(x)\geq 0$ \\
(B) $\Sigma f(x)=1$ \\
(C) $\Sigma xf(x)=2$ \\
(D) $f(x)\leq 1$
\end{multicols}\end{quote}

\item $\mathrm{Var}(x)$ is equal to:
\begin{quote}\begin{multicols}{2}
\noindent (A) $[E(x)]^2-E(x^2)$ \\
(B) $E(x^2)-[E(x)]^2$ \\
(C) $E(x^2)$ \\
(D) $E\{x\}$
\end{multicols}\end{quote}

\item In a discrete probability distribution, the sum of all the probabilities is always:
\begin{quote}\begin{multicols}{2}
\noindent (A) $0$ \\
(B) $1$ \\
(C) $-1$ \\
(D) $\infty$
\end{multicols}\end{quote}

\item For a given binomial distribution with $n$ fixed, if $p<0.5$, then:
\begin{quote}\begin{multicols}{2}
\noindent (A) The binomial distribution will be skewed to the left \\
(B) The binomial distribution will be skewed to the right \\
(C) The binomial distribution will be symmetric \\
(D) None of these
\end{multicols}\end{quote}

\item If the probability of success $p=0.4$ for a particular Bernoulli trial, the expression $\dfrac{7!}{3!\,4!}(0.4)^3(0.6)^4$ gives the probability of getting:
\begin{quote}\begin{multicols}{2}
\noindent (A) Exactly three successes in seven trials \\
(B) Exactly four successes in seven trials \\
(C) Three or more successes in seven trials \\
(D) Four or more successes in seven trials
\end{multicols}\end{quote}

\item The mean of a binomial distribution depends on:
\begin{quote}\begin{multicols}{2}
\noindent (A) Probability of success \\
(B) Probability of failure \\
(C) Number of trials \\
(D) Both (A) and (C)
\end{multicols}\end{quote}

\item The mean of a binomial distribution depends on:
\begin{quote}\begin{multicols}{2}
\noindent (A) Probability of success \\
(B) Probability of failure \\
(C) Number of trials \\
(D) Both
\end{multicols}\end{quote}

\item The standard deviation of a binomial distribution depends on:
\begin{quote}\begin{multicols}{2}
\noindent (A) Probability of success \\
(B) Probability of failure \\
(C) Number of trials \\
(D) Both
\end{multicols}\end{quote}

\item Which of the following can never be described by a binomial distribution:
\begin{quote}\begin{multicols}{2}
\noindent (A) The number of defective items produced by an assembly process \\
(B) The amount of water used by a single household \\
(C) The number of students in the class who can answer this question \\
(D) All of these can always be described by a binomial distribution
\end{multicols}\end{quote}

\item Which of the following is true for the binomial distribution $b(x;n,p)$:
\begin{quote}\begin{multicols}{2}
\noindent (A) Mean $>$ Variance \\
(B) Mean $<$ Variance \\
(C) Mean $=$ Variance \\
(D) Mean $=$ Standard deviation
\end{multicols}\end{quote}

\item The mean, median and mode of the binomial distribution $b(x;n,p)$ will be equal when:
\begin{quote}\begin{multicols}{2}
\noindent (A) $p=0.5$ \\
(B) $p<0.5$ \\
(C) $p>0.5$ \\
(D) None of these
\end{multicols}\end{quote}

\item The parameters of the hypergeometric distribution $h(x;N,n,k)$ are:
\begin{quote}\begin{multicols}{2}
\noindent (A) $x,N,k$ \\
(B) $N,n,k$ \\
(C) $x,n,k$ \\
(D) $x,N,n,k$
\end{multicols}\end{quote}

\item The number of possible outcomes in a Bernoulli trial is:
\begin{quote}\begin{multicols}{2}
\noindent (A) Three \\
(B) Four \\
(C) Two \\
(D) One
\end{multicols}\end{quote}

\item A binomial random variable is a(n):
\begin{quote}\begin{multicols}{2}
\noindent (A) Continuous random variable \\
(B) Discrete random variable \\
(C) Independent variable \\
(D) Dependent variable
\end{multicols}\end{quote}

\item A hypergeometric random variable is a(n):
\begin{quote}\begin{multicols}{2}
\noindent (A) Independent variable \\
(B) Continuous random variable \\
(C) Discrete random variable \\
(D) None of these
\end{multicols}\end{quote}

\item In which distribution does the probability of success remain constant from trial to trial:
\begin{quote}\begin{multicols}{2}
\noindent (A) Hypergeometric distribution \\
(B) Binomial distribution \\
(C) Sampling distribution \\
(D) Continuous distribution
\end{multicols}\end{quote}

\item In which distribution does the probability of success change from trial to trial:
\begin{quote}\begin{multicols}{2}
\noindent (A) Binomial distribution \\
(B) Sampling distribution \\
(C) Hypergeometric distribution \\
(D) Continuous distribution
\end{multicols}\end{quote}

\item In which distribution are the successive trials with replacement:
\begin{quote}\begin{multicols}{2}
\noindent (A) Hypergeometric distribution \\
(B) Binomial distribution \\
(C) Continuous distribution \\
(D) None of these
\end{multicols}\end{quote}

\item In which distribution are the successive trials without replacement:
\begin{quote}\begin{multicols}{2}
\noindent (A) Hypergeometric distribution \\
(B) Binomial distribution \\
(C) Continuous distribution \\
(D) None of these
\end{multicols}\end{quote}

\item Both binomial and hypergeometric distributions are:
\begin{quote}\begin{multicols}{2}
\noindent (A) Continuous probability distributions \\
(B) Discrete probability distributions \\
(C) Neither continuous nor discrete probability distributions \\
(D) Bivariate distributions
\end{multicols}\end{quote}

\item A fair coin is tossed four times. The probability of getting four heads is:
\begin{quote}\begin{multicols}{2}
\noindent (A) $\dfrac{1}{4}$ \\
(B) $\dfrac{1}{2}$ \\
(C) $\dfrac{1}{16}$ \\
(D) $1$
\end{multicols}\end{quote}

\item A fair die is rolled three times. The probability of getting three `aces' is:
\begin{quote}\begin{multicols}{2}
\noindent (A) $\dfrac{1}{3}$ \\
(B) $\dfrac{1}{6}$ \\
(C) $\dfrac{1}{216}$ \\
(D) $\dfrac{1}{27}$
\end{multicols}\end{quote}

\item A fair coin is tossed five times. The probability of getting zero heads is:
\begin{quote}\begin{multicols}{2}
\noindent (A) $0$ \\
(B) $\dfrac{1}{32}$ \\
(C) $\dfrac{1}{5}$ \\
(D) $\dfrac{1}{2}$
\end{multicols}\end{quote}

\item A fair coin is tossed five times. The probability of getting zero heads is:
\begin{quote}\begin{multicols}{2}
\noindent (A) $0$ \\
(B) $\dfrac{1}{32}$ \\
(C) $\dfrac{1}{5}$ \\
(D) $\dfrac{1}{2}$
\end{multicols}\end{quote}

\item In a hypergeometric distribution, the trials are:
\begin{quote}\begin{multicols}{2}
\noindent (A) Independent \\
(B) Dependent \\
(C) Independent and dependent \\
(D) None of these
\end{multicols}\end{quote}

\item In a binomial experiment, the successive trials are:
\begin{quote}\begin{multicols}{2}
\noindent (A) Dependent \\
(B) Independent \\
(C) Mutually exclusive \\
(D) Fixed
\end{multicols}\end{quote}

\item The binomial probability distribution is symmetrical when:
\begin{quote}\begin{multicols}{2}
\noindent (A) $P=q$ \\
(B) $P<q$ \\
(C) $P>q$ \\
(D) $np>npq$
\end{multicols}\end{quote}

\item The percentage of observations lying within the limits $\bar{X}\pm 3S$ in the normal distribution is:
\begin{quote}\begin{multicols}{2}
\noindent (A) $68.26\%$ \\
(B) $95.44\%$ \\
(C) $70.00\%$ \\
(D) $99.73\%$
\end{multicols}\end{quote}

\item In a hypergeometric distribution $N=6$, $n=2$, $K=3$, then mean $=$:
\begin{quote}\begin{multicols}{2}
\noindent (A) $2$ \\
(B) $3$ \\
(C) $1$ \\
(D) $4$
\end{multicols}\end{quote}

\item In a binomial distribution, the range of $X$ is:
\begin{quote}\begin{multicols}{2}
\noindent (A) Zero to $\infty$ \\
(B) Zero to $n$ \\
(C) $-1$ to $\infty$ \\
(D) Undefined
\end{multicols}\end{quote}

\item In a hypergeometric distribution, $n$ is:
\begin{quote}\begin{multicols}{2}
\noindent (A) Changed \\
(B) Zero \\
(C) Fixed \\
(D) Variable
\end{multicols}\end{quote}

\item In a binomial, $n=20$, $p=3/5$, then the variance of this distribution is:
\begin{quote}\begin{multicols}{2}
\noindent (A) $60$ \\
(B) $12$ \\
(C) $4.8$ \\
(D) $0$
\end{multicols}\end{quote}

\item If $n=10$ and $p=0.6$, then $P(X\geq 0)$ is:
\begin{quote}\begin{multicols}{2}
\noindent (A) $0.5$ \\
(B) $0.6$ \\
(C) $1.0$ \\
(D) $1.2$
\end{multicols}\end{quote}

\item A binomial distribution is positively skewed when:
\begin{quote}\begin{multicols}{2}
\noindent (A) $P=q$ \\
(B) $P>q$ \\
(C) $P<q$ \\
(D) $P=\dfrac{1}{2}$
\end{multicols}\end{quote}

\item If $P=q=\dfrac{1}{2}$, then the distribution is called:
\begin{quote}\begin{multicols}{2}
\noindent (A) Symmetrical \\
(B) Positively (skewed) \\
(C) Skewed \\
(D) Negatively (skewed)
\end{multicols}\end{quote}

\item The hypergeometric distribution has \underline{\hspace{1.5cm}} parameters:
\begin{quote}\begin{multicols}{2}
\noindent (A) $2$ \\
(B) $3$ \\
(C) $1$ \\
(D) $4$
\end{multicols}\end{quote}

\item The probability of failure is equal to:
\begin{quote}\begin{multicols}{2}
\noindent (A) $P$ \\
(B) $1-q$ \\
(C) $P-1$ \\
(D) $1-P$
\end{multicols}\end{quote}

\item The mean of the hypergeometric distribution is:
\begin{quote}\begin{multicols}{2}
\noindent (A) $\dfrac{Nn}{k}$ \\
(B) $\dfrac{Nk}{k}$ \\
(C) $\dfrac{nk}{N}$ \\
(D) $\dfrac{n+k}{N}$
\end{multicols}\end{quote}

\item The mean and variance of the binomial distribution are:
\begin{quote}\begin{multicols}{2}
\noindent (A) $np$ and $\sqrt{npq}$ \\
(B) $np$ and $npq$ \\
(C) $np$ and $nq$ \\
(D) $n$ and $p$
\end{multicols}\end{quote}

\item In a binomial experiment with three trials, the variable can take:
\begin{quote}\begin{multicols}{2}
\noindent (A) $2$ values \\
(B) $3$ values \\
(C) $4$ values \\
(D) $5$ values
\end{multicols}\end{quote}

\item The number of trials in a binomial distribution is:
\begin{quote}\begin{multicols}{2}
\noindent (A) Not fixed \\
(B) Fixed \\
(C) Large \\
(D) Small
\end{multicols}\end{quote}

\item In a hypergeometric distribution:
\begin{quote}\begin{multicols}{2}
\noindent (A) Mean $>$ Variance \\
(B) Mean $<$ Variance \\
(C) Mean $=$ Variance \\
(D) Mean $=$ zero
\end{multicols}\end{quote}

\item The binomial distribution is negatively skewed if:
\begin{quote}\begin{multicols}{2}
\noindent (A) $p<\dfrac{1}{2}$ \\
(B) $p=\dfrac{1}{2}$ \\
(C) $p>\dfrac{1}{2}$ \\
(D) $p=1$
\end{multicols}\end{quote}

\item If $N=40$, $n=5$, $k=4$, then the mean of the hypergeometric distribution is:
\begin{quote}\begin{multicols}{2}
\noindent (A) $1$ \\
(B) $1/2$ \\
(C) $1/4$ \\
(D) $1/3$
\end{multicols}\end{quote}

\item A binomial distribution is negatively skewed if:
\begin{quote}\begin{multicols}{2}
\noindent (A) $p<q$ \\
(B) $p>q$ \\
(C) $p=q$ \\
(D) $np=npq$
\end{multicols}\end{quote}

\end{mcqlist}

\newpage

\section{Short Answer Questions (1--667)}

\section*{Write short answers of the following questions.}

\begin{shortlist}

\item Define the word statistics and explain its different meanings?

\item Define the following term: Population and sample

\item Define the following term: Parameter and statistic

\item Distinguish between qualitative and quantitative variables.

\item Write the following using a summation sign with appropriate index:

\item Write the following using a summation sign with appropriate index:

\item Write the following using a summation sign with appropriate index:

\item Write the following using a summation sign with appropriate index:

\item Expand the following summation and product sign:

\item Expand the following summation and product signs:

\item Expand the following summation and product signs:

\item Expand the following summation and product signs:

\item Expand the following summation and product signs:

\item Expand the following summation and product signs:

\item Expand the following summation and production signs:

\item Classify the following as categorical, discrete or continuous variable: Sex of an insect

\item Classify the following as categorical, discrete or continuous variable: Weights of plants

\item Classify the following as categorical, discrete or continuous variable: Major crops of Pakistan

\item Classify the following as categorical, discrete or continuous variable: Level of satisfaction

\item Classify the following as categorical, discrete or continuous variable: Teaching standards

\item Classify the following as categorical, discrete or continuous variable: Temperature measured

\item What are the meanings of statistics? in Fahrenheit

\item Define data.

\item What is datum?

\item Define sampling?

\item What is ratio?

\item What is proportion?

\item How parameters are represented?

\item What is sampling variability?

\item Define experiment.

\item What is sample survey?

\item What is order statistic?

\item What is model?

\item What is random error?

\item What is sigma ?

\item Write Y+Y+Y+...+Y using a summation notation with suitable index. 5 6 7 10

\item Write using a summation notation with suitable index.

\item Expand

\item Expand

\item What is product ?

\item Write using product notation with suitable index.

\item Write using product notation with suitable index.

\item Expand

\item Expand

\item Evaluate n!.

\item Solve 7!.

\item Solve 5!.

\item Write down the importance of statistics in four disciplines.

\item When a variable is fixed?

\item When a variable is random?

\item What are the types of variable?

\item Define qantitative variable.

\item Define qualitative variable.

\item What are the types of quantitative variable?

\item Define continuous variable.

\item define discrete variable.

\item Give some examples of continuous variable.

\item Give some examples of discrete variable.

\item What are the sources of data?

\item What do you mean by statistical method?

\item What is meant by primary data?

\item Define inferential statistic.

\item Explain discrete and continuous data

\item Given x=5, 7, 10, 12 and y=8, 13, 18. Compute and i i

\item Define Statistics.

\item Name the sources of secondary data.

\item Give the uses of statistics. (Only two)

\item Expand the following summation. (a) (b)

\item Differentiate between quantitative and qualitative data.

\item Differentiate between population and sample.

\item Differentiate between a variable and a random variable.

\item Define statistics in singular sense.

\item Define the term error. Name the types of error.

\item Define secondary data with an example.

\item What is the difference between constant and variable?

\item Define statistic.

\item Define qualitative variable.

\item Distinguish between statistics and statistic.

\item What are the two types of quantitative variables?

\item What is description statistics?

\item Define the term population.

\item Enlist three main methods of collection of primary data.

\item Write the names of two branches of statistics.

\item Write two examples each for continuous and discrete variable.

\item Define variable.

\item Write two properties of G.M (Geometric Mean).

\item What is sample?

\item Define parameter.

\item What is statistic?

\item Define constant.

\item Define variable.

\item What is secondary data?

\item Define statistics in plural sense.

\item Expand the following

\item What are different methods of representation of statistical data?

\item What is different method of representation of statistical data

\item What is one-way classification?

\item What is two-way classification?

\item What is three-way classification?

\item Define many way classification.

\item What is tabulation?

\item Define title.

\item What are column captions?

\item Define box-head.

\item Define row captions.

\item What is stub?

\item What is body of table?

\item What is prefatory note?

\item What are footnotes?

\item What is source note?

\item Define frequency distribution.

\item How a frequency distribution is constructed?

\item What are open-end classes?

\item What are class limits?

\item What is class interval?

\item Define class frequency.

\item Define class mark.

\item Define class boundary.

\item What is cumulative frequency?

\item Define relative frequency.

\item Write down some common uses of graphs.

\item Write down some important points for drawing graphs.

\item What do you know about simple bar diagram?

\item What is multiple bar diagram?

\item What is component bar diagram?

\item What is pie-diagram?

\item Define histogram.

\item Give some advantages of histogram.

\item What is frequency polygon?

\item How frequency polygon constructed?

\item What is frequency curve?

\item What is scatter plot?

\item What is bivariate frequency distribution?

\item In pie-diagram, how can we obtain angle for each sector?

\item Write down the steps involved in the construction of simple bar diagram.

\item Write down the steps involved in the construction of multiple bar diagram.

\item Write down the steps involved in the construction of component bar diagram.

\item Write down the steps involoved in the construction of pie-diagram.

\item Write down the steps involved in the construction of histogram.

\item Write down the steps involved in the construction of frequencies polygon.

\item Write down the steps involved in the construction of ogive.

\item Give the frequencies of 5 classes 2, 6, 7, 3, 4, find the cumulative frequencies and relative frequencies.

\item What are the class boundries and size of class interval of (i) 3.5-4.6 (ii) -2-(+2)?

\item Out line the four major parts of s table.

\item Give the frequencies of 5 classes 3, 7, 8, 5 6, find cumulative frequencies, de-cumulative

\item The observations about the number of rooten potatoes from twenty equal sized samples frequencies and relative frequencies. taken from a store are available as follows: 1,2,3,4,3,0,1,2,3,1,1,0,2,1,0,2,3,0,0,1,3. Make frequency table. Also find relative frequencies.

\item How a frequency polygon is constructed?

\item Enlist the main parts of a table.

\item Differentiate between continous and discrete variable.

\item Define classification.

\item Define histogram.

\item Define tabulation.

\item What is simple bar chart?

\item What are the main parts of a table? (Any two only.

\item What is ogive?

\item Define pie chart.

\item Define relative frequency.

\item Explain the term “Open Ended” classes.

\item What is frequency distribution?

\item Differentiate between histogram and historigram.

\item Explain prefatory note and foot note in a table.

\item What is class mark?

\item Differentiate between ungrouped and grouped data.

\item What are the important graphs of frequency distribution?

\item What is size of class interval?

\item What is frequency polygon?

\item Differentiate between class mark and class width.

\item Define Venn diagram.

\item What are the class limits?

\item What graph can be made from quantitative data, name any three graphs?

\item Give three examples of tabular presentation.

\item Define classification.

\item Define histogram.

\item Define tabulation.

\item What is a simple bar chart?

\item What are the main parts of a table? (only any two).

\item What is ogive?

\item Define pie chart.

\item Define relative frequency.

\item Explain the term "Open Ended" classes.

\item What is frequency distribution.

\item Explain prefatory note ans foot note in a table.

\item What is a class mark?

\item Differentiate between ungrouped and grouped data.

\item What are the most graphs of frequency distribution?

\item What is size of class interval?

\item What is frequency polygon?

\item Differentiate between class mark and class width.

\item Define venn diagram.

\item What are the class limits?

\item What graph can be made from quantitative data, name any three graphs?

\item Give three examples of tabular presentation.

\item Find the geometric mean 50, 67, 39, 40, 36, 60, 54, 43.

\item Define Average.

\item What is measure of location?

\item Give some properties of good average.

\item Enlist the main measures of Central tendency.

\item Write down the formula for population mean.

\item Write down the formula for sample mean.

\item Calculate mean for 87,91,89,88,89,91,87,92,90,98.

\item Calculate mean for 10,15,20,25,30,35.

\item For 2,4,5,6,7,6 show that .

\item Prove that =0. (For upground data).

\item Prove that =0. (for grouped data).

\item Show that mean is effected by the change of scale.

\item Write down the coding formula for mean.

\item Write down the short formula for mean.

\item Give some merits of mean.

\item Give some demerits of mean.

\item What is weighted mean?

\item How we calculate weighted mean.

\item \textit{[Equation/content not recoverable --- source PDF contains a broken/unrendered equation at this item.]}

\item How we calculate geometric mean?

\item Calculate G.M for 0.5, 10.0, 2.7, 3.48 and 4.7.

\item Calculate the G.M for 45,30,335,40,44,332,42,37.

\item Write down the properties of G.M.

\item Give some merits of geometric mean.

\item Give some demerits of a geometric mean.

\item When geometric mean is appropriate?

\item When harmonic mean is useful?

\item How H.M is calculated?

\item Calculate H.M for 10,20,30,40,50

\item Calculate G.M for 3,8,11,0 and 10.

\item Calculate H.M for 3,8,11,0 and 10.

\item Write down the properties of harmonic mean?

\item How we denote median?

\item How we calculate median for odd data?

\item Find n=10.

\item How can we calculate median for even data?

\item Calculate the median for 87,88,89,89,90,91,91,92,98.

\item How can we calculate median for grouped data.

\item What are the merits of median?

\item Define quantities.

\item How can we calculate quantiles?

\item What are quartiles?

\item Define lower quartile.

\item How we calculate lower quartile?

\item How we define uppar quartile?

\item How we calculate uppar quartile?

\item Calculate Q1 from 84.60,88.60,94.90,94.50,95.05.

\item What are deciles?

\item How we donate deciles?

\item Where we calculate decile for grouped data?

\item How we calculate percentiles for grouped data?

\item How we find mode for ungrouped data?

\item Write down the merits of mode.

\item When a distribution is symmetrical?

\item What is the empircial relation between mean.

\item If mode=15, median =12, find them.

\item If mean =45, median =30 find mode

\item State some advantages of the mode.

\item The sum of deviation of 10 values from 15 is 20. Find arithmetic mean.

\item If mean =20, median=18.67, find mode.

\item Find G.M if x=1,3,9.

\item Given . Comment on the skewness

\item If mean =20, median=18, find mode for a set of data?

\item If a value of Q , d and P are equal to 72.32, then find the median of distribution

\item Find the mode of 3,3,7,8,10,11,10,12. 2 5 50

\item The sum of deviations of 15 values from 20 is 45. Find arithmetic mean.

\item Write three types of Average.

\item For a frequency distribution of a variable x, it is given that x=10+5u,

\item Find median of +1-3,3, -2 and 5. find the value of mean

\item A distribution consists of 3 components with respective size 42,40 and 65 along with

\item If mode =50, median=12, find mean. their means 2,2.5 and 5. Compute the combined mean.

\item Give various methods of a computing arithmetic mean for grouped data.

\item Why is difficult to calculate geometric mean for (a) 0,4,8 (b) 3,5,-2,1

\item \textit{[Equation/content not recoverable --- source PDF contains a broken/unrendered equation at this item.]}

\item Write down the empirical relation between mean, median and mode. If x1 =1 and x 2 =8 tjat A.M> G.M.

\item What is the median of the word "FREQUENCY"?

\item In a single-peaked symmetrical distribution, mean =5, median =5 then what will be the value of mode?

\item If mean =20, median=18.67, find mode.

\item Find G.M x=1,3,9.

\item Define G.M.

\item Write two qualities of good average.

\item Define mode with formula.

\item Some of deviations of 25 values from 20 is 45. find A.M.

\item If mode =15 and mean 10.5 find mediam.

\item What is the relationship between mean, median and mode in a positively skewed distribution?

\item Write down the formulas for different averages.

\item Define the median.

\item Write down empirical relationship b/w mean and mode

\item Give , fm =25 f1 20, f2 =20 and h=10. find mode.

\item Define harmonic mean.

\item In a symmetrical mean is 35. Find the value of mediann and mode.

\item Discuss advantages of median.

\item Find mode the letter "STATISTICS".

\item What is an array?

\item Average of 5 values is 70. Find the sum of values.

\item If mean is 20, Median =18. Find Mode.

\item Give two properties of arithmetic mean.

\item Observations. Find combined mean.

\item Differentiate between simple arithmetic and weighted arithmetic mean.

\item Calculate the geometric mean 1,1,1,27.

\item Define Arithmetic mean.

\item Calculate the geometric mean 1,1,1,27.

\item Define Arithmetic mean.

\item Compute the median for the data -2,5,0,-1,2,2.

\item In single paked symmetrical distribution mean =5, median =5, then what will be the value of mode?

\item A distribution consists of 3 components with respective size 45, 40 and 65, along with

\item Write down merits of mode (any two) their means 2,202 and 2.5 and cmputer the combined mean.

\item If x 2, x =8 show that A.M>G.M>H.M.

\item The arithmetic mean of 100 items is 5 and the arithmetic mean of 150 item is 4.5 Find 1= 2 combined mean.

\item What are two major properties of Median?

\item In symmetrical distribution mode 20. Find value of Mean and Median.

\item Give that median = 15 and mode = 20 by using the empricial relationship, find the value

\item Define the quartiles and percentages. of mean.

\item Find the mean and sum of the first n natural numbers.

\item What do you mean by statistical average?

\item Define the term measures of central tendency.

\item Calculate median from the given data. 88.03,94.50,95.05,84.60.

\item Define deciles and percentiles.

\item In a moderately a symmetrical distribution the value of median is 42 and the value of mode is 40. Find the mean.

\item Give X 1 4, X 2 =4, show that A.M = G.M = H.M.

\item Given each mean is based on eight values

\item Give the three dis-advantaged of H.M.

\item Give Various methods of computing A.M.

\item Find A.M when D=X-18 f D =150, n=25

\item If A.M = 80, H.M =60, Find G.M.

\item \textit{[Equation/content not recoverable --- source PDF contains a broken/unrendered equation at this item.]}

\item Given X1 =4, X2 =16. Show that G.M = Given 2.8 and n =5, find the sample mean.

\item The sum of deviation of 10 values from x=40 is 250, what is the value of arithmetic

\item Compute geometric mean of 5,25,125. mean?

\item If X1 2= and X2 =8, then shpoe that A.M>G.M>H.M.

\item If mean =75 and mode =70 find median.

\item State when it is impossible to calculate Harmonic Mean.

\item The Arithmetic mean of 20 values is 40.5 what is the total of values?

\item Calculate geometric mean of X=1,1,27.

\item Write down any two demerits of geometric mean.

\item Write down two merits and two de-merits of Harmonic Mean.

\item In a moderately skewed distribution, Mean =25 mode=31, Find value of Median.

\item For a frequency distribution of a variable X, iit is given that X=10+5

\item Write down two demerits of A.M. Find te value of mean

\item If G.M two values is 3. Find the product of two values.

\item Given data Find arithmetic mean of 'X'.

\item What are demerites of arithmetic mean?

\item In moderately skewed distribution, mode=15, median=12, Find its mean.

\item Why are Average called Measures of Central Tendency?

\item Write down the name of Average.

\item What are the types of dispersion?

\item Wha is Absolute dispersion?

\item What is relative dispersion?

\item How can we calculate relative dispersion?

\item What are the different types of absolute dispersion?

\item Define range.

\item How we calculate range from grouped data?

\item How we calculate range from grouped data?

\item What are the merits of range?

\item What are the demeirts of range?

\item How can calculate coefficient of range?

\item Calculate range for the following observations: 13,3,7,15,17,5,23,27.

\item What is quartile deviation?

\item How we calculate inter-quartile range?

\item How we calculate inter-quartile range?

\item If Q=88.03 and Q=94.90, find Q.D and I.Q.R 1 3

\item Calculate quartile deviation for Q=92.3750, Q=102.58333

\item How we calculated M.D for grouped data? 1 3

\item How can we calculating M.D for grouped data?

\item Write down the properties of mean deviation.

\item What are the merits of mean deviation?

\item What are the demirates of mean deviation?

\item How can calculate coefficient of mean deviation?

\item If M.D from Median =3.92 and =16.25, find coefficient of mean deviation.

\item What is variance?

\item How we calculate veriance for ungrouped data?

\item How we calculate variance for grouped data?

\item Define standard deviation.

\item How standard deviation is donated?

\item How can calculate standard deviation for ungrouped data?

\item Give some demerits of standard deviation.

\item Write down some properties of variance.

\item What are moments?

\item What are shappard's corrections?

\item What is coefficient of variation? What are its uses?

\item Define skewness.

\item Define absolute dispersion.

\item Given n=5. Find Co-efficient of variation.

\item Write down the types of dispersion.

\item Define range with formula.

\item What do you mean by co-efficient of variation, write the formula also.

\item Write the formula of Kari Pearson's coefficient of skewness.

\item Define the term "Moments".

\item If Q = 88.03 and Q 94.50. Compute coefficient of Q.D. 1 3

\item Define moment rations. In which units are they expressed?

\item The S.D of a symmetrical distribution is \%. Find its four moments about mean.

\item If mean =25, CV =64\% what will be the variance?

\item Find Bowley's coefficient of skewness. Given Q =53.67, Median= 65.0, Q =76.33. 1 3

\item If m=5, m=-12.8. Find b and comment

\item Define relative dispersion. 2 3 1

\item Compute the correct coefficient of variation when

\item Given =100, Mode =95 and S=10, find coefficient of =67.45, S=8.5275 and class interval is 3 skewness. 2

\item The mean of a distribution is 3 while it second moment about zero is 24, find its

\item Given =296 and mode=2.94. DIscuss the skewness standard deviation.

\item Given =1600, Find coefficient of mean deviation.

\item Define kurtosis

\item The first moment about x=15 is 2.5, the second moment about mean is 16. What is the

\item Give the formula for calculating combined standard deviation for three subgroups value of mean and standard deviation?

\item Give =55, S2 = 121. Find coefficient of variation

\item if Q1 = 88.06, Q3=94.50, compute coefficient of Q.D

\item A series comprises of hundred values each equal to "5". What will be the average and

\item Describe the properties of variance (any two). dispersion of the series?

\item What is symmetrical distribution?

\item Define combined variance.

\item Define standard deviation.

\item What can you say about skewness if mean =1393 and mode =145.8.

\item Write formula of Bowely's coefficient of skewness.

\item Find variance of given data 5,5,5,5,5,5 and 5.

\item If var(x) =8, find var (x+3).

\item If Q=136.62, Q =146 and Q =153.13 find bowley coefficient of skewness. 1 2 3

\item If var(x)=8, find var(2x).

\item Define mean deviation.

\item If Q =136.62, Q=153.13, find co-efficient of quartile deviation.

\item If S.D(x)=10 then find S.D (2x+10)=10? 1 3

\item Define Dispersion.

\item Given Q=125 Q=175. Find the coefficient of quartile deviation. 1 3

\item When the dist. Is lepto kurtic and platy kurtic.

\item When the dist. is lepto kurtic and platy kurtic.

\item Compute range of the data -1,-3, -4, -10, -20.

\item If some of squares of deviation from mean is 192 for 12 values, then find standard deviation of the values.

\item The coefficient of variation of a distribution is 45\% find S.D if its mean is 16.

\item The first two moments about a value x=1.5 of variable are 01 and 1.7. Calculate co- efficient of variation.

\item In the first moment about the number 2 is equal to 5. What is the mean?

\item The lower and upper quartile of a distribution are 142.36 and 167.73 respectively, while median is 153.50. Find co-efficient of skweness.

\item What would the shape and name of the frequency distribution if (a)

\item Define the semi inter quartile range. Mean=Median=Mode (b) Mean>Median>Mode.

\item If lower quartile is 20 and quartile deviation is 5, Find upper quartile.

\item Differentiate between positive and negative skewness.

\item How is variance defined? Also give its formula.

\item Define skewed distribution.

\item What can you say about skewness if median=76.29 and two quartiles are 71.38 and

\item Give names of any four methods for calculating dispersion. 81.25.

\item Write any two properties of standard deviation.

\item Find the range of -1, -3, 0, 2, 3,

\item If var(X) -25, find var (2X+4).

\item Define quartile deviation.

\item if n=10, =40 and S=2 then find coefficient of variation.

\item If Q1=10, Q 20 and Q=30, find the coefficient of skewness. 2= 3

\item What are two merits and one demirats of mean deviation?

\item If a distribuion has mean 14.3 and mode 41487. What can you say about the skewness?

\item If Q-20, Quartle Deviation =30, FInd Q

\item Given mean =50, Median =48, SD=6 Find coefficient of skewness. 1 3.

\item What do you understand by the term "Quantiles"?

\item Find quartile deviation and co efficient of quartile deviation of 7.4, 7.4, 7.4, 7.4, and 7.4.

\item Define negtively skewed data. Give one example from real life.

\item Find Bowley's coefficient of skewness if Q=84, Q=79 and median=81. 1 3

\item In a modertely skewed distribution, mean=100, mode=95 and S.D=10. Find coefficient of

\item Find K for the probality distribution. skewness. x 0 1 2 p(x) 3k 2k k

\item S.D of a distribution is 4,find 2nd moment about mean.

\item If var (x) =10 and y=5x+20, then find var(y).

\item Enlist various relative measures of dispersion.

\item Can mean, median and mode be same, if yes, state in what situation?

\item If first three moments about X=20 of a distribution are: 1, 4, 40, then find the value of "b1".

\item What is the range of Bowley's coefficient of skewness?

\item If Q=13.73, Q=38.29, compute Quartile deviation.

\item Calculate range of 13, 3, 4, 15, 17, 5, 13, 27. 1 3

\item What are different measures of absolute dispersion?

\item If var(x) =4, Find var (3x).

\item What is skewness?

\item When distribution is symmetrical?

\item When distribution is negatively skewed?

\item When distribution is normal?

\item When distribution is platykutic?

\item When distribution is lepokurtic?

\item For what values S.D and variance are equal?

\item Give the Bowley's formula of a coefficient or skewness.

\item What would be the shape of the distribution if Mean<Median<Mode?

\item Which absolute measure of dispersion is the best?

\item Define inter-quartile range and quartile deviation.

\item Write four relative measures of dispersion.

\item If coefficient of skewness is zero, then what will you say about mean, median and

\item Range of 5 observation is -6 and smallest observations is -2, find the largest mode? observation.

\item Given E(x)= 204 and C.V(x)=6, find S.D (x).

\item Define Price Relatives.

\item . Fin index number

\item Given most suitable average for index number, with reason.

\item Given W= 20, 30, 40, 50, 60, l = 100, 105, 110, 105, 120. Find weighted index.

\item What are quantity index numbers?

\item Define composite index number.

\item Define link relative.

\item Write the Formula for Fisher's index no

\item Given Find base year weighted index number.

\item What are types of Composite index number?

\item Write down some mein steps involved in construction of index number

\item What is variable.

\item What are random numbers?

\item How we generate random numbers?

\item What is pseudo random technique?

\item What is the application of random numbers?

\item Give some examples of discretes random variable.

\item Define continuous random variable.

\item Give some examples of continuous random variable.

\item What are the types of random variable?

\item Which methods are used in drawing random numbers?

\item Define random experiment.

\item Define continuous random variable?

\item What is mean by random variable?

\item Can a random variable assume negative values?

\item What is the application of random numbe?

\item What is discrete random variable?

\item Make random variable for tail when three coins are tossed.

\item Write down the assumptions of random numbers.

\item Differentiate between a variable and random variable.

\item Write down the properties of a continuous random vireable.

\item How random number can be generated?

\item Give two examples of random variable.

\item Describe the two different types of random variable.

\item What is the same sample space of a random experiment?

\item Properties of discrete random variable.

\item Make a random variable for number of heads when two coins are tossed.

\item Two dice are thrown make a random variable for sum of dots.

\item Define random numbers.

\item What are the expectations and standard deviation of a constant?

\item Given a random variable X with E (x) = 12 and S.D (X) =5. Find E(x2).

\item What does P.D.F stands for?

\item Write three properties of random experiment.

\item Define discrete random variable.

\item What is discrete random variable?

\item Give some examples of dicrete random variable?

\item Define continuous random variable.

\item Give some examples of continuous random variable.

\item How do we describe probability mass function?

\item How we calculated variance of variable $Y$?

\item If $E(X)=1.1$ and $E(X^2)=2.1$, find $\mathrm{Var}(X)$.

\item If $E(X)=1.1$, find $E(3X+5)$.

\item Given: 
\begin{center}
\begin{tabular}{c|ccc}
$X$ & 30 & 35 & 40 \\
\hline
$P(X)$ & $15/45$ & $24/45$ & $6/45$
\end{tabular}
\end{center}
Find $E(X)$.

\item If $E(X)=7$, $E(X^2)=54.83$, find $\mathrm{Var}(X)$.

\item If $E(X)=200$, $C.V(X)=7\%$, find $\mathrm{Var}(X)$.

\item If $E(Y)=2.1$, $E(Y^2)=9.8$, find $\mathrm{Var}(Y)$.

\item If $\mathrm{Var}(X)=2.07$, find $\mathrm{S.D}(X)$.

\item If $E(X)=9.94$, $\mathrm{S.D}(X)=1.44$, find $C.V(X)$.

\item Head:Tail $=3:1$, find $C.V(X)$.

\item If $E(X)=5$, find $E(-3X+2)$.

\item If $E(5X+10)=18$, $E(5X+10)^2=530$, find $\mathrm{Var}(5X+10)$.

\item Find $E(5-3Y)$, when $E(Y)=1.4$.

\item Check for $Y=1,2,3,4$. Is $f(y)=\dfrac{y}{10}$ a p.d.f?

\item If $\mathrm{Var}(X)=5$, $\mathrm{Var}(Y)=\mathrm{Var}(3X+2)$, then find variance and S.D. of $X$?

\item $x: 0\ 1\ 2$; $P(x): 4c,\ 3c,\ c$. Find the value of $c$.

\item If $E(X)=1.1$ and $E(X^2)=2.1$, find $\mathrm{Var}(X)$.

\item What is mathematical expectation?

\item Give the properties of probability distribution.

\item If $E(X)=2$ and $E(X^2)=20$. Find standard deviation of $X$.

\item If $E(X)=3$ and $E(Y)=2.5$, find $E(X+Y)$.

\item What is distribution function?

\item Define probability density function (p.d.f).

\item Write down properties of expectation (any two).

\item Given $E(X^2)=400$ and $\mathrm{Var}(X)=144$, then find mean of $X$ values (i.e. $E(X)$).

\item Write down properties of probability density function.

\item $S^2=3$ and $E(X)=3$. Find $E(X^2)$.

\item $\mathrm{Var}(X)=2$ and $\mathrm{Var}(Y)=5$. Find $\mathrm{Var}(3Y-X)$ if $X$ and $Y$ are independent.

\item What are the properties of discrete probability distribution?

\item If $E(X)=5$, find $E(-3X+2)$.

\item What is discrete probability distribution?

\item Given $E(X)=10$ and $E(X^2)=500$, find coefficient of variation.

\item What are properties of distribution function?

\item Given $E(X)=0$ and $E(Y)=1$, find $E(3X-2Y)$.

\item Given $f(x)=\dfrac{x}{6}$, $x=1,2,3$. Find $K$.

\item If $E(X^2)=400$, $S.D(X)=12$, find $E(X)$?

\item If $E(X)=0.63$, $\mathrm{Var}(X)=0.2331$, find $E(X^2)$?

\item If $E(X)=0.63$, $\mathrm{Var}(X)=0.2331$, find $E(X^2)$?

\item If $E(5X+10)=18$, $E(5X+10)^2=530$, find $\mathrm{Var}(5X+10)$.

\item Given $E(X)=2.5$, $E(X^2)=7.5$, find variance and coefficient of variation.

\item Given $E(X)=2.5$, $E(X^2)=7.5$, find variance and coefficient of variation.

\item Define continuous properties of expectation.

\item What is the mathematical expectation to win Rs.\ 50 if a head appears and lose Rs.\ 50 if a tail appears in a single toss of a fair coin?

\item Given $E(X)=0.55$, $\mathrm{Var}(X)=1.35$ and $Y=2X+1$. Find $E(Y)$ and $\mathrm{Var}(Y)$.

\item If $E(X)=3$, $E(Y)=2.5$, find $E(2X+Y)$.

\item What is an expected value of a random variable?

\item Given $X=-10,-20,30$ and $P(X)=2k,3k,5k$, find $K$.

\item Given $P(X)=C|X-2|$ and $X=-1,0,1,3$, find $E(X)$.

\item Given $X=0,2,3$ and $P(X)=\dfrac{X}{5}$, find $E(X)$.

\item If $E(X)=1.1$, find $E(3X+5)$.

\item Given $X=0,1,2$ and $P(X)=\dfrac{X}{3}$, is this a probability function?

\item A continuous random variable $X$ has a density function $f(x)=\dfrac{1}{4}$, $0\le x \le 4$.

\item If $E(X)=1.15$, then find $E(3X+5)$.

\item If $E(X)=1.1$, $E(X^2)=2.1$, find $\mathrm{Var}(X)$?

\item If $E(X)=1.1$, find $E(3X+5)$?

\item Find $E(X)$ when $\mathrm{Var}(X)=4$, $E(X^2)=20$.

\item Given $E(X^2)=400$, $S.D(X)=12$, then find $E(X)$.

\item Given $E(X^2)=400$, $S.D(X)=12$, find $E(X)$.

\item What is probability mass function?

\item Write down the properties of probability mass function.

\item Find $E(X)$, $E(X^2)$ from the following data.

\item A continuous random variable $X$ has a density function $f(x)=\dfrac{1}{2}$ for $a<x<4$. Find the value of $C$.

\item What does p.d.f. stand for?

\item Given $E(X)=0$ and $E(X)^2=1$. Find $E(2X^2-3X+5)$.

\item $E(X)=1.1$ and $E(X^2)=2.1$, find $\mathrm{Var}(X)$.

\item A die is rolled if $X=$ face value. Find $E(X)$.

\item What will be the $K$ in the given probability distribution?

\item Given $f(x)=kx$; $x=1,2,3$. Find $k$.

\item If $E(X)=10$, then find $E(5X-25)$.

\item Given $X=0,1,2$: $P(X)=4c,3c,c$, find the value of $c$.

\item Define trial.

\item What is Binomial experiment?

\item What are the properties of binomial distribution?

\item What are the parameters of the binomial distribution?

\item What are binomial coefficients?

\item When is binomial distribution skewed?

\item When is binomial distribution positively skewed?

\item When is binomial distribution negatively skewed?

\item If $n=4$, $p=0.5$, find $P(x=3)$.

\item If you toss a fair die 6 times, what is the probability of getting no even number?

\item Out of 800 families with 5 children each. How many would you expect to have 4 boys?

\item Give two properties of a Binomial experiment.

\item In binomial distribution $n=10$ and $p=0.5$, find its variance.

\item If $n=10$, $k=4$, $K=6$. Find variance of the hypergeometric distribution. OR find $E(X)$ and $S.D(X)$.

\item If $n=5$, $p=0.25$, find $P(X\ge1)$.

\item Define Bernoulli trials.

\item Criticize the statement that in a binomial distribution mean $=6$ and standard deviation is $2.5$.

\item Find mean and variance of following binomial distribution: $\left(\dfrac{1}{2}+\dfrac{1}{2}\right)^{2}$.

\item Define Hyper-geometric distribution.

\item In a binomial distribution $n=5$ and $p=0.5$. Find coefficient of variation.

\item In a binomial distribution, mean $=6$ and variance $=4$. Find its parameters.

\item What is the variance of Hypergeometric distribution?

\item Given $N=6$, $n=3$ and $K=4$, find $P(x=0)$.

\item What is a binomial experiment?

\item What are the parameters of hyper-geometric distribution?

\item If mean$=20$, $P=0.2$ in binomial distribution, then find $n$.

\item What is mean and variance of binomial distribution?

\item If $n=8$ and $p=0.4$, then find variance of binomial distribution.

\item If $N=10$, $n=3$, $k=4$, find variance of Hypergeometric distribution.

\item What is binomial probability distribution?

\item What is the range of binomial r.v. ``$X$''?

\item If a fair coin is tossed 5 times, find the probability of getting at the most 2 heads.

\item Define hyper-geometric experiment.

\item $N=10$, $n=2$ and $k=3$. Find $P(x=0)$.

\item If $n=8$, $p=0.3$ in binomial probability distribution, then find mean and variance.

\item Define the parameters of the binomial probability distribution.

\item Find the number of trials of a binomial distribution which has mean$=12$ and standard deviation$=2$.

\item Define binomial random variable.

\item If $N=10$, $N-K=4$, $n=5$. Find variance of Hypergeometric distribution.

\item If $E(X)=3$ and $\mathrm{Var}(X)=2$, then find parameters of binomial distribution of random variable $X$.

\item If $n=12$ and $p=0.6$, then find mean and variance of binomial distribution.

\item Is it possible to have binomial distribution with mean 10 and standard deviation 4?

\item What is hyper-geometric random variable?

\item If $N=10$, $N-K=4$, $n=5$, find $P(X=1)$.

\item If $n=10$ and $P=0.4$, find mean and variance of binomial distribution.

\item A random variable $X$ has a binomial distribution with $n=5$ and $p=0.2$; find $P(x=2)$.

\item If $E(X)=2.4$ and $p=0.3$, find the S.D. of binomial distribution.

\item Given $N=8$, $n=3$ and $K=4$. Find $P(x=2)$.

\item If $N=10$ and $q=0.4$, then find $P(x=2)$.

\item Is it possible to have a binomial distribution with mean $=5$ and S.D. $=3.7$?

\item If $X$ is binomially distributed with $n=10$ and $p=0.4$, then find variance of $Y=2X+3$.

\item For a binomial distribution with $n=10$ and $p=0.5$, find probability of successes.

\item Given $N=10$, $n=4$ and $k=3$. Find $P(x=1)$.

\item In a binomial distribution mean$=2.4$ and S.D.$=1.2$. Find the value of ``$n$''.

\item If $X\sim h(x;11,5,4)$. Find mean and variance of hyper-geometric distribution.

\item If ``$x$'' is a binomial random variable with $n=9$, $p=0.3$, then find $E(y)$ and $\mathrm{Var}(y)$.

\item In a binomial distribution with $n=5$, what is the value of ``$P$'' if $P(x=0)=P(X=1)$?

\item If $N=40$, $n=5$, $K=4$, then find values of mean and variance of Hyper-geometric distribution.

\item In a binomial distribution mean $=38$ and $q=0.83$, find $n$ and $p$.

\item Write down formula for mean and variance of hyper-geometric distribution.

\item What will be the mean and variance of binomial distribution if $n=6$ and $p=0.6$?

\item If $N=10$, $n=5$, $K=3$, find mean of hypergeometric distribution using formula of mean.

\item State probability function of Hypergeometric distribution. Also write its mean and variance.

\item Given $N=10$, $n=2$ and $K=2$. Find $P(X=0)$.

\item If $n=4$, $p=0.4$, find $P(x=3)$?

\item What is binomial distribution?

\item What is mean, variance and S.D. of $(q+p)^{n}$?

\item Given $P(x)=k\cdot{}^{4}C_{X}$ and $x=0,1,2,3,4$, find the values of $K$.

\item In hypergeometric, $N=40$, $n=5$, $K=4$. Find mean and variance.

\item Find the mean of hypergeometric distribution when $N=8$, $n=2$ and $K=8$.

\item If $n=10$ and $p=0.5$. Find mean, S.D. of binomial distribution.

\item In binomial distribution, if $n=10$ and $p=0.3$, then find variance of binomial distribution.

\item In binomial distribution, $n=7$, $p=0.7$. Find coefficient of variation.

\item Under which circumstance can we apply the binomial distribution and the hypergeometric distribution?

\item Give any two properties of hypergeometric experiment which are different from binomial experiment.

\item If in binomial distribution Mean $=3.2$, $\mathrm{Var}=1.152$. Find $n$ and $p$.

\item The sum of 10 items is 12 and sum of squares equal to 26.9, then find standard deviation.

\end{shortlist}

\newpage

\section{Long Questions (1--250)}

\section*{Write detailed answers of the following questions.}

\begin{longlist}

\item Explain in detail the main aspects of a statistical problem.

\item Define descriptive and inferential statistics and differentiate between them.

\item Distinguish between primary and secondary data and give different sources from which these are obtained.

\item Define the Histogram, the frequency polygon and frequency curve.

\item What do you understand by classification and tabulation? Discuss their importance in a statistical analysis.

\item Distinguish between one-way and two-way tables. Illustrate your answers with examples. Also explain the following: (i) Classification according to attributes (ii) Class limits (iii) Length of class interval (iv) Class frequency.

\item The following table gives the details of monthly budgets of two families. Represent these figures through a suitable diagram.
\begin{center}
\begin{tabular}{lcc}
Item & Family A & Family B \\
\hline
Food & Rs.\ 600 & Rs.\ 800 \\
Clothing & Rs.\ 100 & Rs.\ 100 \\
House Rent & Rs.\ 400 & Rs.\ 500 \\
Fuel and Lighting & Rs.\ 100 & Rs.\ 100 \\
Miscellaneous & Rs.\ 300 & Rs.\ 500 \\
\hline
Total & Rs.\ 1500 & Rs.\ 2000
\end{tabular}
\end{center}

\item Represent the following data through a pie diagram.
\begin{center}
\begin{tabular}{lc}
Items of Expenditure & Amount \\
\hline
Food & 4000 \\
Clothing & 1000 \\
House Rent & 2500 \\
Education & 1000 \\
Fuel and Light & 600 \\
Miscellaneous & 200
\end{tabular}
\end{center}

\item Define frequency Histogram. Draw a Histogram for the following frequency distribution giving the steps involved.
\begin{center}
\begin{tabular}{l|cccccccc}
Mid Values ($X$) & 32 & 37 & 42 & 47 & 52 & 57 & 62 & 67 \\
Frequency ($f$) & 3 & 17 & 28 & 47 & 54 & 31 & 14 & 4
\end{tabular}
\end{center}

\item [10.] (i)(a) Write down the important points for drawing graphs.\\
(b) In order to estimate the mean length of leaves from a certain tree, a sample of 100 leaves was chosen and their lengths measured in millimeters. A grouped frequency table was set up with results as follows:
\begin{center}
\begin{tabular}{l|ccccccccc}
Mid value & 2.2 & 2.7 & 3.2 & 3.7 & 4.2 & 4.7 & 5.2 & 5.7 & 6.2 \\
Frequency & 3 & 5 & 8 & 12 & 18 & 24 & 20 & 8 & 2
\end{tabular}
\end{center}
(ii)(a) Display the table in the form of a frequency polygon.\\
(b) What are the boundaries of the interval whose mid-point is 3.7 cm?

\item In a locality, total area is 500 acres where 250 acres are under sugarcane, 125 acres are under maize, 60 acres are under wheat and the remaining 65 acres are under other crops. Make a pie-diagram to represent the distribution of acreage under different crops.

\item A biologist was interested to know whether male spiders are longer than female spiders. He collected a random sample of female and male green lynx spiders given below. Advise him.
\begin{center}
\begin{tabular}{cccc|cccc}
\multicolumn{4}{c|}{Length of female (mm)} & \multicolumn{4}{c}{Length of male (mm)} \\
\hline
5.7 & 5.2 & 4.7 & 5.8 & 8.2 & 9.9 & 5.9 & 8.4 \\
6.1 & 6.3 & 7.0 & 5.7 & 9.6 & 7.0 & 6.6 & 7.8 \\
7.5 & 6.4 & 6.5 & 4.7 & 7.5 & 9.8 & 6.3 & 8.3 \\
6.2 & 5.4 & 6.2 & 5.8 & 8.0 & 9.1 & 6.3 & 8.4 \\
4.8 & 5.9 & 6.2 & 5.8 & 8.0 & 7.4 & 7.9 & 10.2 \\
5.6 & 5.5 & & & & & &
\end{tabular}
\end{center}

\item [13.] (i) What is meant by tabulation? Explain the main steps which are generally taken in tabulation.\\
(ii) What is frequency distribution? How is it constructed?

\item The following data gives the lifetime in minutes, recorded to the nearest tenth of a minute, of 50 sprayed insects.
\begin{center}
\begin{tabular}{cccccccccc}
1.2 & 2.2 & 0.7 & 1.7 & 1.9 & 1.4 & 1.8 & 1.8 & 2.0 & 4.3 \\
2.5 & 0.9 & 3.4 & 2.8 & 3.7 & 3.5 & 0.4 & 2.8 & 1.1 & 0.2 \\
3.9 & 6.3 & 2.5 & 2.1 & 1.3 & 2.1 & 0.3 & 0.4 & 2.4 & 2.1 \\
3.5 & 2.9 & 1.2 & 5.3 & 1.7 & 2.7 & 1.8 & 4.8 & 3.2 & 1.6 \\
2.6 & 1.8 & 2.3 & 1.3 & 3.1 & 1.5 & 2.6 & 5.9 & 2.0 & 2.3
\end{tabular}
\end{center}
Using 8 intervals with the lowest starting at 0.1:\\
(i) Form a frequency distribution and cumulative frequency distribution.\\
(ii) Also draw Histogram and frequency polygon for the frequency distribution so formed.

\item [15.] (i) What are the advantages of diagrammatic representation?\\
(ii) Explain the following: (a) A Bar diagram (b) Subdivided bar diagram (c) Multiple bar diagram.

\item The following data gives the record of a company's savings over the years. Draw a bar diagram to represent it.
\begin{center}
\begin{tabular}{l|cccccccc}
Year & 1950 & 1951 & 1952 & 1953 & 1954 & 1955 & 1956 & 1957 \\
Rs.\ (000) & 1010 & 2050 & 3458 & 1980 & 2300 & 1295 & 1520 & 1070
\end{tabular}
\end{center}

\item Draw a sub-divided bar diagram to represent the male and female population of four divisions of Punjab in 1961.
\begin{center}
\begin{tabular}{lccc}
Division & Male & Female & Both Sexes \\
\hline
Lahore & 35 & 30 & 65 \\
Multan & 35 & 31 & 66 \\
Sargodha & 32 & 28 & 60 \\
Rawalpindi & 21 & 19 & 40
\end{tabular}
\end{center}

\item The following information is available about the age (in years) and weights (in kilograms) of individuals as (age, weight) pairs. Make a scatter plot and bivariate frequency distribution.
\[
\begin{aligned}
&(20,60),(15,55),(14,45),(17,60),(16,48),(22,70),(16,63),(14,55),(18,57),(19,67),\\
&(21,67),(17,65),(13,60),(15,60),(17,49),(19,65),(23,73),(21,65),(22,70),(14,50),\\
&(16,60),(19,59),(15,62),(17,59),(24,75).
\end{aligned}
\]

\item [19.] (i) Define arithmetic mean, geometric mean and harmonic mean. Explain the situations when each of them is used most appropriately.\\
(ii) The relation between arithmetic mean (A.M.), geometric mean (G.M.) and harmonic mean (H.M.) is $A.M \ge G.M \ge H.M$.

\item A man traveling 100 kilometers has 5 stages, at equal intervals. The speed of the man in the various stages was observed to be $10,16,20,14,15$ kilometers per hour. Find the average speed at which the man travels.

\item Calculate mean, median and mode for the following table:
\begin{center}
\begin{tabular}{lcc}
Items & Family A & Family B \\
\hline
Food & Rs.\ 600 & Rs.\ 800 \\
Clothing & Rs.\ 100 & Rs.\ 100 \\
House rent & Rs.\ 400 & Rs.\ 500 \\
Fuel and lighting & Rs.\ 100 & Rs.\ 100 \\
Miscellaneous & Rs.\ 300 & Rs.\ 500 \\
\hline
Total & 1500 & 2000
\end{tabular}
\end{center}

\item Calculate mean, median and mode for the grouped data of the following table:
\begin{center}
\begin{tabular}{l|cccccccc}
Mid values ($x$) & 32 & 37 & 42 & 47 & 52 & 57 & 62 & 67 \\
Frequency ($f$) & 3 & 17 & 28 & 47 & 54 & 31 & 14 & 4
\end{tabular}
\end{center}

\item Calculate the following for the data of exercise 2.8(b): (i) calculate an estimate for the mean leaf length (ii) construct a cumulative frequency table and use it to estimate the sample median.
\begin{center}
\begin{tabular}{l|ccccccccc}
Mid value & 2.2 & 2.7 & 3.2 & 3.7 & 4.2 & 4.7 & 5.2 & 5.7 & 6.2 \\
Frequency & 3 & 5 & 8 & 12 & 18 & 24 & 20 & 8 & 2
\end{tabular}
\end{center}

\item What do you understand by weighted mean? In what circumstances is it preferred to ordinary mean and why?

\item Define the mode of a frequency distribution. How does it compare with other types of averages?

\item [26.] (i) Write down the empirical relation between mean, median and mode for a unimodal distribution of moderate asymmetry. Illustrate graphically the relative positions of the mean, median and mode for frequency curves which are skewed to the right and to the left.\\
(ii) For a certain frequency distribution, with the mean and median 45 and 36 respectively, find the mode approximately using the empirical relation between the three.

\item Bilal gets a rise of 10\% in salary at the end of his first year of service and further rises of 20\% and 25\% at the end of the second and third years respectively. The rise in each case is calculated on his salary at the beginning of the year. To what annual percentage increase is this equivalent?

\item Find the mean for the following distribution:
\begin{center}
\begin{tabular}{l|ccccccc}
Classes & 0--10 & 10--40 & 40--90 & 90--100 & 100--105 & 105--120 & 120--140 \\
$f$ & 40 & 110 & 150 & 200 & 120 & 30 & 20
\end{tabular}
\end{center}

\item The frequency distribution given below has been derived from the use of a working origin. If $D=X-18$, find arithmetic mean and geometric mean:
\begin{center}
\begin{tabular}{l|cccccccc}
$D$ & $-12$ & $-8$ & $-4$ & 0 & 4 & 8 & 12 & 16 \\
$F$ & 2 & 5 & 8 & 18 & 22 & 13 & 8 & 4
\end{tabular}
\end{center}

\item The reciprocals of 11 values of $x$ are given below: $0.500, 0.454, 0.0400, 0.333, 0.0285, 0.0232, 0.0213, 0.0200, 0.0182, 0.0151, 0.0143$. Calculate Harmonic Mean and Arithmetic Mean of the data.

\item Reciprocals of $x$ are given below: $0.0267, 0.0235, 0.0211, 0.0191, 0.0174, 0.0160, 0.0148$. Calculate Harmonic Mean of the data.

\item The cities A, B and C are equidistant from each other. Fatima travels from A to B at the speed of 30 miles per hour by car. From B to C at speed of 50 miles per hour. Determine her average speed for the entire trip.

\item Harmonic mean and Geometric mean of two numbers are 3.2 and 4 respectively. Find their arithmetic mean and both the numbers as well.

\item The arithmetic mean and geometric mean of three numbers are 34 and 18 respectively. Find all three numbers, when the geometric mean of the first two numbers is 9.

\item [35.] Find out: (i) The average rate of motion in the case of a person who rides the first mile at the rate of 10 miles per hour, the next mile at the rate of 8 miles per hour, the third at the rate of 6 miles per hour.\\
(ii) Increase in population which in the first decade has increased 20\%, in the next 25\%, and in the third 30\% (assumed).

\item The given table shows the distribution of the maximum loads (in short tons) supported by certain cables produced by a company. Determine Mean, Median and Mode.
\begin{center}
\begin{tabular}{lc}
Maximum Loads & No.\ of cables \\
\hline
9.3--9.7 & 2 \\
9.8--10.2 & 5 \\
10.3--10.7 & 12 \\
10.8--11.2 & 17 \\
11.3--11.7 & 14 \\
11.8--12.2 & 6 \\
12.3--12.7 & 3 \\
12.8--13.2 & 1
\end{tabular}
\end{center}

\item Compute mean, median, mode, 6th Decile and 74th percentile for the data given in the table:
\begin{center}
\begin{tabular}{lc}
Classes & Frequency \\
\hline
0.7312--0.7313 & 10 \\
0.7314--0.7315 & 15 \\
0.7316--0.7317 & 20 \\
0.7318--0.7319 & 25 \\
0.7320--0.7321 & 30 \\
0.7322--0.7323 & 8 \\
0.7324--0.7325 & 2
\end{tabular}
\end{center}

\item Find the value $Q_3$, $D_6$, $P_5$ and mode for the following data:
\begin{center}
\begin{tabular}{lc|lc}
Groups & Frequency & Groups & Frequency \\
\hline
0--4.9 & 3 & 25--29.9 & 13 \\
5--9.9 & 4 & 30--34.9 & 13 \\
10--14.9 & 9 & 35--39.9 & 5 \\
15--19.9 & 11 & 40--44.9 & 2 \\
20--24.9 & 15 & 45--49.9 & 2
\end{tabular}
\end{center}

\item If for any frequency distribution the mean is 45 and the median is 30, find mode approximately, using the formula connecting the three.

\item A bus traveling 200 miles has ten stages at equal intervals. The speed of the bus in the various stages was observed to be $10,15,20,75,20,30,40,50,30,40$ miles per hour. Find the average speed at which the bus has traveled.

\item The following data has been obtained from a frequency distribution of a continuous variable $x$ after making the substitution $u=\dfrac{x-A}{h}$.
\begin{center}
\begin{tabular}{l|cccccccc}
$U$ & $-4$ & $-3$ & $-2$ & $-1$ & 0 & 1 & 2 & 3 \\
$F$ & 2 & 5 & 8 & 18 & 22 & 13 & 8 & 4
\end{tabular}
\end{center}
Find Harmonic Mean.

\item Salman obtained the following marks in a certain examination. Find the weighted mean if weights $4,3,3,2,2$ respectively are allotted to the subjects.
\begin{center}
\begin{tabular}{lccccc}
& English & Urdu & Math & State & Physics \\
\hline
Marks & 73 & 82 & 80 & 57 & 62
\end{tabular}
\end{center}

\item Calculate weighted mean for the following items:
\begin{center}
\begin{tabular}{lcc}
Items & Expenditures & Weights \\
\hline
Food & 290 & 7.8 \\
Rent & 54 & 2 \\
Clothing & 98 & 1.5 \\
Fuel and Light & 75 & 1.0 \\
Other items & 75 & 0.5
\end{tabular}
\end{center}

\item For a certain distribution, if $\Sigma(X-\bar{X})=0$, what is the value of A.M. and why?

\item For a certain distribution, if $\Sigma(X-\bar{X})=0$, what is the value of A.M. and why?

\item Arithmetic Mean of 15 values is 20 and by adding more values, the mean remains 20. Find the new three values if ratio is $a:b:c :: 3:2:1$.

\item If $D=X-A$, find mean and G.M.
\begin{center}
\begin{tabular}{l|cccccccc}
$D$ & $-12$ & $-8$ & $-4$ & 0 & 4 & 8 & 12 & 16 \\
$F$ & 2 & 5 & 8 & 18 & 22 & 13 & 8 & 4
\end{tabular}
\end{center}

\item Find Median and Mode:
\begin{center}
\begin{tabular}{l|cccc}
Classes & 5--9 & 10--14 & 15--19 & 20--24 \\
$F$ & 7 & 18 & 13 & 4
\end{tabular}
\end{center}

\item Find geometric mean from the following frequency distribution:
\begin{center}
\begin{tabular}{l|ccccc}
$x$ & 2 & 3 & 4 & 5 & 6 \\
$f$ & 5 & 7 & 8 & 3 & 2
\end{tabular}
\end{center}

\item The average marks obtained by three sections of a statistics class are given below:
\begin{center}
\begin{tabular}{lcc}
Classes & No.\ of Students & Mean \\
\hline
A & 45 & 95 \\
B & 42 & 58 \\
C & 38 & 52
\end{tabular}
\end{center}
Find combined mean of the whole class.

\item Compute mode of the data:
\begin{center}
\begin{tabular}{l|cccccc}
Hourly Wages & 4--6 & 6--8 & 8--10 & 10--12 & 12--14 & 14--16 \\
No.\ of employees & 13 & 11 & 182 & 105 & 19 & 7
\end{tabular}
\end{center}

\item Find two numbers whose mean is 9 and G.M. is 7.2.

\item Find H.M. and G.M. from the following values: $2,4,6,8,10,12,14$.

\item Find mean and median from the following data:
\begin{center}
\begin{tabular}{l|ccccc}
Marks & 11--20 & 21--30 & 31--40 & 41--50 & 51--60 \\
$F$ & 2 & 4 & 6 & 8 & 6
\end{tabular}
\end{center}

\item Calculate the arithmetic mean using step deviation method from the following data:
\begin{center}
\begin{tabular}{l|ccccc}
Marks & 1--10 & 11--20 & 21--30 & 31--40 & 41--50 \\
$F$ & 5 & 30 & 40 & 25 & 10
\end{tabular}
\end{center}

\item Consider the following data:
\begin{center}
\begin{tabular}{l|ccccc}
Wages & 30--40 & 40--50 & 50--60 & 60--70 & 70--80 \\
$F$ & 2 & 5 & 6 & 5 & 2
\end{tabular}
\end{center}
Compute $Q_1$ and 7th decile.

\item The arithmetic mean and harmonic mean of two observations are 5 and 3.2 respectively. Find the geometric mean.

\item Compute the median from the following data:
\begin{center}
\begin{tabular}{l|ccccc}
Classes & 0--7 & 7--14 & 14--21 & 21--28 & 28--35 \\
$f$ & 2 & 11 & 15 & 9 & 4
\end{tabular}
\end{center}

\item Following are the numbers of mistakes made by typists in a set: $5,6,0,5,8,4,3,5,7,9,4,7,8,6,8,3,2,7,1,4,5,9,6,4,5,10,5,4,9,6$. Find A.M. also find mode.

\item Calculate geometric mean for the following data:
\begin{center}
\begin{tabular}{l|cccccc}
Marks & 10--19 & 20--29 & 30--39 & 40--49 & 50--59 & 60--69 \\
No.\ of students & 5 & 10 & 15 & 20 & 12 & 8
\end{tabular}
\end{center}

\item The weight of 100 students are recorded to the nearest kg:
\begin{center}
\begin{tabular}{l|ccccc}
Weight & 62--64 & 65--67 & 68--70 & 71--73 & 74--76 \\
$F$ & 19 & 25 & 34 & 17 & 5
\end{tabular}
\end{center}
Find arithmetic mean.

\item Calculate mode for the given data:
\begin{center}
\begin{tabular}{l|ccccc}
Group & 10--14 & 15--19 & 20--24 & 25--29 & 30--34 \\
$F$ & 4 & 8 & 12 & 9 & 3
\end{tabular}
\end{center}

\item Compute the median from the following data:
\begin{center}
\begin{tabular}{l|ccccc}
Group & 10--14 & 15--19 & 20--24 & 25--29 & 30--34 \\
$F$ & 1 & 4 & 10 & 5 & 2
\end{tabular}
\end{center}

\item The logarithms of 6 values are: $1.8129, 1.9037, 1.9777, 1.9294, 1.6990, 1.7782$. Find the arithmetic mean of variable $x$.

\item Calculate the Harmonic Mean from the following data:
\begin{center}
\begin{tabular}{l|cccccc}
Classes & 10--14 & 15--19 & 20--24 & 25--29 & 30--34 & 35--39 \\
$F$ & 1 & 4 & 5 & 11 & 9 & 2
\end{tabular}
\end{center}

\item Find geometric mean for the given data.
\begin{center}
\begin{tabular}{l|ccccc}
$X$ & 2 & 3 & 4 & 5 & 6 \\
$F$ & 5 & 7 & 8 & 3 & 2
\end{tabular}
\end{center}

\item Find the median from the given data:
\begin{center}
\begin{tabular}{lc}
Classes & Frequency \\
\hline
$0-10$ & 1 \\
$10-19$ & 2 \\
$10.19$ & 3 \\
$20-29$ & 11 \\
$30-39$ & 24 \\
$40-49$ & 32 \\
$50-59$ & 40
\end{tabular}
\end{center}

\item The mean weight of 150 students in a class is 60 kg. The mean weight of boys in the class is 70 kg and that of the girls is 55 kg. Find the number of boys and girls in the class.

\item The reciprocals of the variable ``$X$'' are $0.02857, 0.2222, 0.1534, 0.01176, 0.2222, 0.01333$. Show that A.M. $>$ Geometric mean $>$ Harmonic mean for the above data.

\item Calculate median for the following frequency distribution:
\begin{center}
\begin{tabular}{l|cccc}
Classes & 0.7312--0.7313 & 0.7314--0.7315 & 0.7316--0.7317 & 0.7318--0.7319 \\
$F$ & 10 & 15 & 20 & 8
\end{tabular}
\end{center}

\item Find Arithmetic mean and Harmonic mean of the following distribution.
\begin{center}
\begin{tabular}{l|ccccc}
Height & 60--62 & 63--65 & 66--68 & 69--71 & 72--74 \\
No.\ of Students & 5 & 18 & 42 & 27 & 8
\end{tabular}
\end{center}

\item Reciprocals of 8 values are $0.0667, 0.1111, 0.833, 0.556, 0.0357, 0.0500, 0.0278, 0.222$. Calculate median of $X$.

\item The mean weight of 45 boys is 120 lbs, that of 37 girls is 116.7 lbs and that of 5 teachers is 122.3 lbs. Calculate the mean weight of boys, girls and teachers combined.

\item A man gets a rise of 10\% in salary at the end of his first year of service and further rises of 20\% and 25\% at the end of the second and third years respectively. The rise in each case is calculated on his salary at the beginning of the year. To what annual percentage increase is this equivalent?

\item Find mean, median for the following data:
\begin{center}
\begin{tabular}{l|ccccc}
No.\ of rooms & 30 & 33 & 40 & 45 & 52 \\
No.\ of students & 3 & 10 & 8 & 5 & 3
\end{tabular}
\end{center}

\item Calculate H.M. and G.M. of the given distribution.
\begin{center}
\begin{tabular}{l|cccccc}
$x$ & 12 & 14 & 16 & 18 & 20 & 21 \\
$f$ & 1 & 4 & 6 & 10 & 7 & 2
\end{tabular}
\end{center}

\item Find the lower and upper quartiles from the data given below:
\begin{center}
\begin{tabular}{l|cccccc}
Time (Minutes) & 10--19 & 20--29 & 30--39 & 40--49 & 50--59 & 60--69 \\
No.\ of customers & 4 & 8 & 46 & 20 & 40 & 2
\end{tabular}
\end{center}

\item What do you understand by dispersion? What are the most useful methods of measuring dispersion? Indicate the advantages of these methods.

\item Define mean deviation and its coefficient. Discuss its advantages and uses.

\item [80.] (i) What is semi-interquartile range? (ii) Define range and discuss its uses.

\item Explain the difference between absolute dispersion and relative dispersion. Describe the properties of the standard deviation.

\item [82.] Define various measures of dispersion and give their formulae. The following table gives the marks of students:
\begin{center}
\begin{tabular}{l|ccccc}
Marks & 30--39 & 40--49 & 50--59 & 60--69 & 70--79 \\
$f$ & 8 & 87 & 190 & 86 & 20
\end{tabular}
\end{center}
Calculate: (a) Quartile Deviation (b) coefficient of skewness.

\item Find mean deviation from mean for the following data:
\begin{center}
\begin{tabular}{l|ccccc}
Marks & 30--39 & 40--49 & 50--59 & 60--69 & 70--79 \\
Frequency & 8 & 87 & 190 & 86 & 20
\end{tabular}
\end{center}

\item Find first moment about mean and second moment about mean of the following data:
\begin{center}
\begin{tabular}{l|ccccc}
$X$ & 12 & 14 & 16 & 18 & 20 \\
$F$ & 1 & 4 & 6 & 10 & 7
\end{tabular}
\end{center}

\item The data of differences between the ages of husbands and wives is given below. Calculate mean deviation from mean.
\begin{center}
\begin{tabular}{l|ccccc}
Difference & 0--5 & 5--10 & 10--15 & 15--20 & 20--25 \\
Freq. & 14 & 20 & 10 & 5 & 1
\end{tabular}
\end{center}

\item Given Mode $=2.944$ and the second moment about mean is $1.4884$. Calculate coefficient of skewness.

\item Find Quartile deviation for the following distribution:
\begin{center}
\begin{tabular}{l|ccccc}
Weights & 70--74 & 75--79 & 80--84 & 85--89 & 90--94 \\
$F$ & 1 & 7 & 15 & 5 & 2
\end{tabular}
\end{center}

\item The first three moments of a distribution about the value 2 of the variable are $1, 16$ and $-40$. Show that the mean is 3, the variance is 15 and $m_3$ is $-86$.

\item Compute coefficient of skewness for the data given below: $45,50,56,60,62,63,65,70$.

\item Calculate coefficient of variation from the following data:
\begin{center}
\begin{tabular}{l|ccccc}
$X$ & 0 & 1 & 2 & 3 & 4 \\
$f$ & 17 & 9 & 6 & 5 & 3
\end{tabular}
\end{center}

\item Calculate M.D. from mean for the following data:
\begin{center}
\begin{tabular}{l|cccccc}
$X$ & 110 & 120 & 130 & 140 & 150 & 160 \\
$f$ & 4 & 6 & 7 & 7 & 8 & 2
\end{tabular}
\end{center}

\item First three moments about $x=62$ are $1, 4$ and $10$. Find first three moments about mean.

\item Calculate the variance for the following data:
\begin{center}
\begin{tabular}{l|cccc}
Marks & 3--5 & 5--7 & 7--9 & 9--11 \\
Freq. & 8 & 6 & 7 & 4
\end{tabular}
\end{center}

\item First four moments about mean of a distribution are $0.3, 74, 45$ and $121$. Find out whether the distribution is leptokurtic or platykurtic.

\item Calculate coefficient of variation for each series and interpret the result -- which series is more consistent.

\item The lower and the upper quartiles of a distribution are $242.36$ and $267.74$ respectively, while the median is $253.60$. Find coefficient of skewness and interpret the result.

\item Calculate the mean deviation from the following data:
\begin{center}
\begin{tabular}{l|ccccc}
Classes & 10--19 & 20--29 & 30--39 & 40--49 & 50--59 \\
$F$ & 2 & 3 & 8 & 11 & 3
\end{tabular}
\end{center}

\item Calculate the first four moments about $X=15$ from the following data: $10,12,14,16,18$.

\item The mean of 200 items is 48 and their S.D is 3. Find the sum and sum of squares of all this data.

\item A manufacturer of two types of light bulbs A and B has the following results: Type A: $\bar{X}_A=1495$ hrs \& $S_A=280$ hrs; Type B: $\bar{X}_B=1895$ hrs \& $S_B=310$ hrs. Which bulb has greater: (i) Absolute dispersion (ii) Relative dispersion?

\item The first three moments of a distribution about the value $x=4$ are $1,4$ and $10$. Find the C.V. Is the distribution symmetrical?

\item The mean of 200 items is 48 and their S.D is 3. Find the sum and sum of squares of all these.

\item Calculate the quartile deviation for the following data.
\begin{center}
\begin{tabular}{l cccccc}
\toprule
Weight (kg) & 5 & 10 & 15 & 20 & 25 & 30 \\
No. of Roads & 5 & 10 & 15 & 20 & 20 & 30 \\
\bottomrule
\end{tabular}
\end{center}

\item Given Median $=2.944$, $\sum f(x-\bar{x})^3/n=120^\dagger$ and variance is $1.4884$, calculate Pearson's coefficient of skewness.

\item The following table gives the marks of students. Calculate the quartile deviation.
\begin{center}
\begin{tabular}{l ccccc}
\toprule
Marks & 30--39 & 40--49 & 50--59 & 60--69 & 70--79 \\
$f$ & 8 & 87 & 190 & 86 & 20 \\
\bottomrule
\end{tabular}
\end{center}

\item Find the four moments of a distribution about $x=2$: $1, 2.5, 5.5$ and $16$. Calculate the mean and coefficient of variation (C.V).

\item Find the mean deviation from the median for the following frequency distribution.
\begin{center}
\begin{tabular}{l cccc}
\toprule
Age (years) & 2--10 & 10--15 & 15--20 & 20--25 \\
Frequency & 10 & 20 & 30 & 15 \\
\bottomrule
\end{tabular}
\end{center}

\item The first three moments of a distribution about $y=4$ are $1$, $4$ and $-25$ respectively.$^\dagger$ Find the coefficient of variation (C.V). Is the distribution symmetrical, positively skewed, or negatively skewed?

\item Compute the quartile deviation for the following data:
\begin{center}
\begin{tabular}{l ccccc}
\toprule
Marks & 40--50 & 50--60 & 60--70 & 70--80 & 80--90 \\
No.\ of students & 30 & 50 & 60 & 40 & 20 \\
\bottomrule
\end{tabular}
\end{center}

\item The first four moments of a set of numbers about $3$ are $-2, 10, -25$ and $50$. Determine the corresponding moments about the mean.

\item What are index numbers? Give the uses of an index number.

\item Define an index number. Discuss the main points involved in the construction of an index number of prices.

\item Define an index number and describe the different types of index numbers.

\item Discuss the important problems involved in the construction of index numbers of prices.

\item Compare the following steps: (i) Simple and Composite index. (ii) Fixed and chain base method.

\item What is the weighted index number?

\item Find the index number of price from the following data taking the average price of all years as the base.
\begin{center}
\begin{tabular}{l cccccccc}
\toprule
Years & 1970 & 1971 & 1972 & 1973 & 1974 & 1975 & 1976 & 1977 \\
Prices & 15 & 19 & 21 & 30 & 37 & 38 & 40 & 48 \\
\bottomrule
\end{tabular}
\end{center}

\item Given the price of a commodity per maund for the period 1945 to 1960 as:
\begin{center}
\begin{tabular}{l r l r l r}
\toprule
Year & Price & Year & Price & Year & Price \\
\midrule
1945 & 12.7 & 1950 & 24.58 & 1955 & 15.65 \\
1946 & 18.97 & 1951 & 20.90 & 1956 & 16.15 \\
1947 & 19.70 & 1952 & 19.80 & 1957 & 20.20 \\
1948 & 13.50 & 1953 & 23.65 & 1958 & 25.25 \\
1949 & 15.65 & 1954 & 24.55 & 1959 & 32.40 \\
\bottomrule
\end{tabular}
\end{center}
Construct index numbers correct to 2 decimal places: (i) 1945 as base (ii) Average price of all years as base.

\item Find the index number using: (i) 1977 as base (ii) Average of the price as base:
\begin{center}
\begin{tabular}{l r l r l r}
\toprule
Years & Prices & Years & Prices & Years & Prices \\
\midrule
1977 & 22.5 & 1980 & 30 & 1983 & 37.5 \\
1978 & 25 & 1981 & 35 & 1984 & 47.5 \\
1979 & 27.5 & 1982 & 32.5 & 1985 & 45 \\
\bottomrule
\end{tabular}
\end{center}

\item Following the data below, find fixed index numbers taking (i) 1930 as base (ii) Average of first 3 years as base (iii) the year 1935 as base.
\begin{center}
\begin{tabular}{l r l r l r}
\toprule
Years & Prices & Years & Prices & Years & Prices \\
\midrule
1930 & 4 & 1933 & 7 & 1936 & 9 \\
1931 & 5 & 1934 & 8 & 1937 & 10 \\
1932 & 6 & 1935 & 10 & 1938 & 11 \\
\bottomrule
\end{tabular}
\end{center}

\item Construct index numbers of prices for the following data taking 1960 as base:
\begin{center}
\begin{tabular}{l r l r}
\toprule
Years & Prices & Years & Prices \\
\midrule
1960 & 50 & 1965 & 72 \\
1961 & 52 & 1966 & 73 \\
1962 & 55 & 1967 & 75 \\
1963 & 57 & 1968 & 71 \\
1964 & 62 & 1969 & 70 \\
\bottomrule
\end{tabular}
\end{center}

\item The price is Rs.\ per maund of coal sold during the years 1953--58 as given below. Compute the index number of price for the year 1953 as base.
\begin{center}
\begin{tabular}{l r l r l r}
\toprule
Years & Prices & Years & Prices & Years & Prices \\
\midrule
1953 & 14.95 & 1955 & 15.10 & 1957 & 16.28 \\
1954 & 14.95 & 1956 & 15.65 & 1958 & 16.28 \\
\bottomrule
\end{tabular}
\end{center}

\item From the data given below, compute the index numbers of prices, taking 1962 as base.
\begin{center}
\begin{tabular}{l cccc}
\toprule
Year & Firewood & Softcake & Kerosene oil & Matches \\
\midrule
1962 & 3.25 & 2.50 & 0.20 & 0.06 \\
1963 & 3.44 & 2.80 & 0.22 & 0.06 \\
1964 & 3.50 & 2.00 & 0.25 & 0.06 \\
1965 & 3.75 & 2.50 & 0.25 & 0.06 \\
\bottomrule
\end{tabular}
\end{center}

\item Compute the index number of prices from the following data taking 1981 as base and using median as average.
\begin{center}
\begin{tabular}{l ccc}
\toprule
Year & A & B & C \\
\midrule
1981 & 18 & 85 & 52 \\
1982 & 22 & 76 & 60 \\
1983 & 28 & 80 & 66 \\
1984 & 31 & 95 & 80 \\
\bottomrule
\end{tabular}
\end{center}

\item Find chain index numbers (using geometric mean to average the relatives) for the following data of prices, taking 1970 as the base year.
\begin{center}
\begin{tabular}{l ccccc}
\toprule
Commodities & 1970 & 1971 & 1972 & 1973 & 1974 \\
\midrule
A & 40 & 43 & 45 & 42 & 50 \\
B & 160 & 162 & 165 & 161 & 168 \\
C & 240 & 29 & 52 & 23 & 27 \\
D & 240 & 245 & 247 & 250 & 255 \\
\bottomrule
\end{tabular}
\end{center}

\item The following table gives the average wholesale prices in rupees per unit of gold, wheat and cotton during the years 1912--1917. Construct an index number with 1912 as base using: (i) A.M. (ii) G.M.
\begin{center}
\begin{tabular}{l cccccc}
\toprule
Commodities & 1912 & 1913 & 1914 & 1915 & 1916 & 1917 \\
\midrule
Gold & 25.3 & 30.8 & 33.4 & 35.5 & 35.4 & 36.0 \\
Wheat & 17.3 & 14.5 & 4.9 & 5.7 & 17.7 & 11.6 \\
Cotton & 7.8 & 5.4 & 6.7 & 5.6 & 7.2 & 10.2 \\
\bottomrule
\end{tabular}
\end{center}

\item Construct chain indices for the following years, taking 1940 as base.
\begin{center}
\begin{tabular}{l ccccc}
\toprule
Item & 1940 & 1941 & 1942 & 1943 & 1944 \\
\midrule
Wheat & 2.80 & 3.40 & 3.60 & 4.00 & 4.20 \\
Rice & 2.95 & 3.60 & 2.90 & 2.75 & 2.75 \\
Maize & 3.10 & 3.50 & 3.40 & 4.50 & 3.70 \\
\bottomrule
\end{tabular}
\end{center}

\item Construct index numbers for 1963 assuming 1953 as the base period by (i) Laspeyre's formula (ii) Paasche's formula.
\begin{center}
\begin{tabular}{l cccc}
\toprule
& \multicolumn{2}{c}{1953} & \multicolumn{2}{c}{1963} \\
Commodities & Price & Quantity & Price & Quantity \\
\midrule
A & 2 & 50 & 10 & 40 \\
B & 3 & 10 & 8 & 5 \\
C & 4 & 5 & 4 & 5 \\
\bottomrule
\end{tabular}
\end{center}

\item Compute the weighted index numbers for 1964 from the following data with 1960 as base:
\begin{center}
\begin{tabular}{l cc cc}
\toprule
& \multicolumn{2}{c}{Price} & \multicolumn{2}{c}{Quantity} \\
Commodities & 1960 & 1964 & 1960 & 1964 \\
\midrule
Milk & 3.95 & 4.25 & 97.75 & 104.36 \\
Cheese & 34.80 & 38.90 & 78 & 83 \\
Butter & 61.56 & 59.70 & 118 & 116 \\
\bottomrule
\end{tabular}
\end{center}

\item Compute Fisher's index number for the following data.
\begin{center}
\begin{tabular}{l cccc}
\toprule
& \multicolumn{2}{c}{Base Year} & \multicolumn{2}{c}{Current Year} \\
Commodities & Price & Quantity & Price & Quantity \\
\midrule
A & 7 & 70 & 5 & 49 \\
B & 5 & 27 & 7 & 28 \\
C & 10 & 35 & 9 & 29 \\
D & 9 & 50 & 4 & 42 \\
E & 3 & 16 & 10 & 25 \\
\bottomrule
\end{tabular}
\end{center}

\item Construct the following with the help of the data given below: Fisher's ideal index taking 1970 as base.
\begin{center}
\begin{tabular}{l cc cc}
\toprule
& \multicolumn{2}{c}{Total Production (tons)} & \multicolumn{2}{c}{Harvest Price (Rs)} \\
Commodities & 1970 & 1971 & 1970 & 1971 \\
\midrule
Rice & 71 & 26 & 3.80 & 3.50 \\
Barley & 107 & 83 & 2.90 & 1.90 \\
Maize & 72 & 48 & 2.90 & 1.80 \\
\bottomrule
\end{tabular}
\end{center}

\item Calculate Fisher's ideal index from the following data with 1965 as base year.
\begin{center}
\begin{tabular}{l cccc}
\toprule
& \multicolumn{2}{c}{1965} & \multicolumn{2}{c}{1970} \\
Commodities & Price & Quantity & Price & Quantity \\
\midrule
A & 4.6 & 102 & 9.50 & 96 \\
B & 3.7 & 15 & 7.36 & 28 \\
C & 10.2 & 17 & 8.42 & 21 \\
D & 8.9 & 19 & 9.87 & 13 \\
\bottomrule
\end{tabular}
\end{center}

\item Define weighted and unweighted index numbers and explain why weighted index numbers are preferred over unweighted index numbers.

\item Calculate Laspeyre's, Paasche's and Fisher's ideal index for the following data with 1992 as base.
\begin{center}
\begin{tabular}{l cc cc}
\toprule
& \multicolumn{2}{c}{Average Price (Rs)} & \multicolumn{2}{c}{Quantity (Units)} \\
Item & 1992 & 1993 & 1992 & 1993 \\
\midrule
Wheat flour & 4.38 & 4.57 & 20kg & 16kg \\
Rice & 14.15 & 15.58 & 10kg & 12kg \\
Moong pulse & 18.67 & 17.28 & 1kg & 1kg \\
Gram pulse & 10.41 & 16.36 & 1kg & 1kg \\
\bottomrule
\end{tabular}
\end{center}

\item The following figures give the average annual price in U.K. for beef and mutton.
\begin{center}
\begin{tabular}{l cc l cc l cc}
\toprule
Years & Beef & Mutton & Year & Beef & Mutton & Year & Beef & Mutton \\
\midrule
1935 & 54 & 75 & 1938 & 62 & 62 & 1941 & 72 & 85 \\
1936 & 54 & 73 & 1939 & 61 & 68 & 1942 & 76 & 90 \\
1937 & 61 & 73 & 1940 & 72 & 85 & 1943 & 79 & 90 \\
\bottomrule
\end{tabular}
\end{center}
Construct an index number of meat prices giving weights 2 and 1 for beef and mutton respectively. Take the year 1935 as the base year.

\item The following table shows the average price in rupees for wheat, rice and barley.
\begin{center}
\begin{tabular}{l ccc}
\toprule
Year & Wheat & Rice & Barley \\
\midrule
1980 & 175.5 & 480.4 & 82.4 \\
1981 & 180.3 & 509.7 & 90.6 \\
\bottomrule
\end{tabular}
\end{center}
Taking 1980 as the base year, construct the price index number by the weighted average of relatives method for the year 1981 using weights 20 for wheat, 12 for rice and 4 for barley.

\item An inquiry into the budgets of middle-class families in a city of England gave the following information. What change in the cost of living do the figures of 1929 show as compared to 1928?
\begin{center}
\begin{tabular}{l cccc}
\toprule
Expenses on & Food 35\% & Rent 15\% & Clothing 20\% & Fuel 10\% \\
\midrule
Price (1928) & 150 & 30 & 75 & 25 \\
Price (1929) & 145 & 30 & 65 & 23 \\
\bottomrule
\end{tabular}
\end{center}

\item (i) Define random variable and give an example to explain it. (ii) What are random numbers and how are these generated? Also give an example to explain their application. (iii) Classify each of the following random variables as either discrete or continuous: (a) the number of pages in a book, (b) the number of questions asked in an oral examination, (c) the lifetime of a light bulb, (d) the amount of rainfall at a particular location during different months of 1996.

\item Generate the first 6 random digits using the pseudo-random number generator with $m=100$, $a=21$, $b=7$, $x_0=10$.

\item Count the number of heads and tails when a single coin is tossed 10 times, without actually throwing a coin.

\item Two coins are tossed. Let $Y$ denote the number of heads; the possible number of heads and their corresponding probabilities are given in the adjoining table.
\begin{center}
\begin{tabular}{cc}
\toprule
$Y$ & $p(Y)$ \\
\midrule
0 & 0.25 \\
1 & 0.50 \\
2 & 0.25 \\
\bottomrule
\end{tabular}
\end{center}
Simulate the sequence of number of heads for 20 tosses without actually throwing a coin.

\item Let the digits $0,1,2,3,4$ represent heads and $5,6,7,8,9$ represent tails. Use random numbers to simulate 20 flips of a coin.

\item Two students in a class are asked to take one of the two courses Mathematics (M) or Biology (B). Define the random variable $Y$ as the number of students taking Biology. Write down the possible outcomes and the values assigned by the random variable $Y$.

\item Two coins are tossed and let the random variable $Y$ denote the number of heads. Write down the possible outcomes and the values assigned to the random variable.

\item There are three children in a family. Let the random variable denote the number of boys in the family. Write down the possible outcomes and the values assigned by the random variable, assuming equal chances for boys and girls.

\item Four balls are drawn from a bag containing 5 white and 3 black balls. If $X$ denotes the number of white balls drawn, write down the possible values of the random variable $X$.

\item Differentiate between discrete and continuous random variables and give an example of each.

\item What are random numbers? How can they be generated? Explain the applications of random numbers.

\item Three balls are drawn from a bag containing 5 white and 3 black balls. If $X$ denotes the number of white balls drawn from the bag, find the probability distribution of $X$.

\item There are seven candidates for three positions of typist. Four of the candidates know Urdu typing while the other three do not. If three candidates are selected at random, find the probability distribution of the number of persons knowing Urdu typing among those selected.

\item (i) What is meant by a probability distribution? Distinguish between discrete and continuous random variables by giving examples. (ii) A coin is tossed 4 times. If $X$ denotes the number of tails, what is the probability distribution of $X$? Draw a probability histogram. (iii) A bag contains 4 red and 6 black balls. For a sample of 4 balls, find the probability distribution of $X$.

\item (i) Define a probability function and give an example to explain it. (ii) Two fair dice are thrown and $Y$ denotes the product of the two scores; obtain the probability distribution of $Y$.

\item (i) What properties should a mathematical function possess to be a probability function/probability density function? (ii) Check whether the following functions satisfy the conditions of a probability function: a. $f(y)=\dfrac{1}{4}$ for $y=1,2,3,4,5$, b. $f(y)=\dfrac{1}{y}$ for $y=1,2,3,4$, c. $f(y)=\dfrac{y}{15}$ for $y=0,1,2,3,4,5$, d. $f(y)=\dfrac{(5-y)^2}{6}$ for $y=0,1,2,3$.

\item Determine the value of $c$ so that the function can serve as a probability function of a random variable: a. $cy$ for $y=1,2,3,4,5$; b. $(1-c)c^y$ for $y=0,1,2$.

\item (i) Given the following probability distribution:
\begin{center}
\begin{tabular}{l ccccc}
\toprule
$x$ & 0 & 1 & 2 & 3 & 4 \\
$P(x)$ & $1/126$ & $20/126$ & $60/126$ & $40/126$ & $5/126$ \\
\bottomrule
\end{tabular}
\end{center}
Verify that $E(2X+3)=2E(X)+3$. (ii) Let $X$ be a random variable with probability distribution:
\begin{center}
\begin{tabular}{l ccccc}
\toprule
$x$ & $-1$ & 0 & 1 & 2 & 3 \\
$P(x)$ & 0.125 & 0.500 & 0.200 & 0.050 & 0.125 \\
\bottomrule
\end{tabular}
\end{center}
Find (a) $E(X)$ and $\operatorname{Var}(X)$; (b) the probability distribution of the random variable $Y=2X+1$, and using it determine $E(Y)$ and $\operatorname{Var}(Y)$.

\item \emph{(Repeat of the previous question, as given in the original.)} (i) Given the following probability distribution:
\begin{center}
\begin{tabular}{l ccccc}
\toprule
$x$ & 0 & 1 & 2 & 3 & 4 \\
$P(x)$ & $1/126$ & $20/126$ & $60/126$ & $40/126$ & $5/126$ \\
\bottomrule
\end{tabular}
\end{center}
Verify that $E(2X+3)=2E(X)+3$. (ii) Let $X$ be a random variable with probability distribution:
\begin{center}
\begin{tabular}{l ccccc}
\toprule
$x$ & $-1$ & 0 & 1 & 2 & 3 \\
$P(x)$ & 0.125 & 0.500 & 0.200 & 0.050 & 0.125 \\
\bottomrule
\end{tabular}
\end{center}
Find (a) $E(X)$ and $\operatorname{Var}(X)$; (b) the probability distribution of the random variable $Y=2X+1$, and using it determine $E(Y)$ and $\operatorname{Var}(Y)$.

\item The following table gives the probability distribution of the random variable $Y$, the number of courses taught by a teacher during the spring semester in the University of Agriculture Faisalabad.
\begin{center}
\begin{tabular}{l ccccc}
\toprule
$y$ & 1 & 2 & 3 & 4 & 5 \\
$P(y)$ & 0.1 & 0.2 & 0.3 & 0.3 & 0.1 \\
\bottomrule
\end{tabular}
\end{center}
Find the probability that: (i) the teacher taught 2 courses, (ii) the teacher taught less than 4 courses, (iii) the teacher taught between 2 and 5 courses, (iv) the teacher taught at least 3 courses, (v) the teacher taught at most 4 courses.

\item A box contains five slips of paper marked $1,2,3,4$ and $5$. Two slips are selected without replacement; list the possible values for each of the following random variables: (i) the sum of the two numbers on the slips, (ii) the difference between the first and second number.

\item Four randomly selected students from a class are asked their opinion about the teaching system as satisfactory (S) or not satisfactory (NS). Let $Y$ denote the number of students saying satisfactory. Write down the possible outcomes and the possible values.

\item A point is randomly selected on the surface of a lake that has a maximum depth of 30 feet. Let $y$ denote the depth of the lake at the randomly selected point. What are the possible values for $Y$? Is $Y$ a discrete variable or a continuous variable?

\item Calculate the mean and variance of the following probability distribution.
\begin{center}
\begin{tabular}{l cccccc}
\toprule
$y$ & 0 & 1 & 2 & 3 & 4 & 5 \\
$P(y)$ & 0.1 & 0.2 & 0.3 & 0.2 & 0.1 & 0.1 \\
\bottomrule
\end{tabular}
\end{center}

\item (i) A continuous random variable $X$ having values only between 0 and 4 has a density function given by $f(x)=\dfrac{1}{2}-ax$, where $a$ is a constant. Find (a) $a$, (b) $P(1<X<2)$. (ii) A continuous random variable $X$ has probability density function $f(x)=cx$ for $0<X<2$. Find (a) $c$, (b) probability that $1<X<1.5$, (c) probability that $X<1.5$. (iii) If $f(x)$ has probability density $kx^2$, $0<x<1$, determine its kind and find the probability that $1/3<X<1/2$.

\item What is a random variable? Distinguish between discrete and continuous random variables, giving examples.

\item Find the probability distribution of the number of boys in families with three children, assuming equal probabilities for boys and girls.

\item \emph{(Repeat of the previous question, as given in the original.)} Find the probability distribution of the number of boys in families with three children, assuming equal probabilities for boys and girls.

\item Define continuous random variable and its probability distribution. (ii) Find the constant $k$ so that the function $f(x)$ defined as $f(x)=kx^\dagger$ for the appropriate range, and $f(x)=0$ elsewhere, may be a density function.

\item (i) What do you mean by expected value? What are the properties of expectation? (ii) Given the following discrete probability distribution:
\begin{center}
\begin{tabular}{l cccccc}
\toprule
$x$ & 0 & 1 & 2 & 3 & 4 & 5 \\
$P(x)$ & $6/36$ & $10/36$ & $8/36$ & $6/36$ & $4/36$ & $2/36$ \\
\bottomrule
\end{tabular}
\end{center}
Compute its mean, variance, standard deviation and coefficient of variation.

\item Let $X$ be a random variable with probability distribution as follows:
\begin{center}
\begin{tabular}{l ccccc}
\toprule
$x$ & 1 & 2 & 3 & 4 & 5 \\
$f(x)$ & 0.125 & 0.45 & 0.25 & 0.05 & 0.125 \\
\bottomrule
\end{tabular}
\end{center}
Find the mean and variance.

\item (i) A continuous random variable $X$ has a density function $f(x)=\dfrac{x+1}{8}$ for $x=2$ to $x=4$. Find (a) $P(X<3.5)$, (b) $P(2.4<X<3.5)$, (c) $P(X=1.5)$. (ii) A continuous random variable $X$ has density function $f(x)=2x$, $0\leq x\leq 1$. Find (a) $P(X=1/2)$, (b) $P(X>1/4)$, (c) $P(1/4<X<1/2)$.

\item In a summer season, a dealer of desert room coolers can earn Rs.\ 800 per day if it is cloudy. Find his mathematical expectation if the probability of the day being hot is 0.40 and of being cloudy is 0.35.

\item A committee of size 5 is to be selected from 3 female and 5 male members. Find the expected number of female members on the committee.

\item What are the properties of probability functions and probability density functions?

\item (i) Explain the concept of a distribution function. (ii) 10 vegetable cans, all of the same size, have lost their labels. It is known that 5 contain tomatoes and 5 contain corn. If 5 cans are selected at random, find the probability distribution and the distribution function for the number of tomato cans in the sample.

\item (i) Define the expected value for a random variable. (ii) The number of automobile accidents in a city are $1,2,3,4$ with corresponding probabilities $1/8, 2/8, 2/8$ and $3/8$. What is the expected number of daily accidents?

\item A and B throw a die for a prize of Rs.\ 11, to be won by the player who first throws a 6. If A has the first throw, what are their respective expectations?

\item A bag contains 2 white and 2 black balls. Three men A, B and C draw a ball, without replacement. The person who draws the white ball first receives Rs.\ 12. What are their respective expectations?

\item Three balls are drawn from a bag containing 5 white and 3 black balls. Find the probability distribution of the number of white balls drawn, and find the mean and variance of the distribution when it is repeated 4 times.

\item Approximately 10\% of the glass bottles coming from a production line have serious defects. If two bottles are selected at random, find the expected number of bottles having serious defects.

\item A random variable $X$ takes the values $-3,-2,2,3$ and $4$ with probabilities $P(X)$ equal to $1/5,1/10,1/10,1/5$ and $2/5$ respectively. Compute $E(X)$ and show that $E(5X+10)=5E(X)+10$. Also compute $\operatorname{Var}(X)$ and $\operatorname{Var}(5X+10)$. Find the ratio of the two variances.

\item If $f(x)=\dfrac{1}{11}$ for $x=2,3,4,\ldots,12$, then find the mean and variance of the random variable $X$.

\item \emph{(Repeat of the previous question, as given in the original.)} If $f(x)=\dfrac{1}{11}$ for $x=2,3,4,\ldots,12$, then find the mean and variance of the random variable $X$.

\item For the following probability distribution, find: (i) $E(X)$, (ii) $E(X^2)$, (iii) $E[X-E(X)]^2$.
\begin{center}
\begin{tabular}{l ccc}
\toprule
$x$ & $-10$ & $-20$ & 30 \\
$P(X)$ & $1/5$ & $3/10$ & $1/2$ \\
\bottomrule
\end{tabular}
\end{center}

\item (i) Define the expectation of a random variable. (ii) The probability distribution of a discrete random variable $X$ is given by $f(x)=\dbinom{3}{x}\left(\dfrac{1}{4}\right)^x\left(\dfrac{3}{4}\right)^{3-x}$, $x=0,1,2,3$. Find $E(X)$ and $E(X^2)$.

\item Find the mean and variance of the following probability distribution.
\begin{center}
\begin{tabular}{l cccccc}
\toprule
$x$ & 0 & 1 & 2 & 3 & 4 & 5 \\
$P(x)$ & 0.1 & 0.2 & 0.3 & 0.2 & 0.1 & 0.1 \\
\bottomrule
\end{tabular}
\end{center}

\item If $F(x)=1-e^{-x}$ for $x>0^\dagger$, show that it is a p.d.f (via its density) and also find $P(X<4)$.

\item Show that $f(x)=\dfrac{3+2x}{2}$, $2\leq x\leq 4$ is a proper density function$^\dagger$ and find the probability $P(2\leq x\leq 3)$.

\item A pair of dice is rolled. Let $X$ represent the sum of dots on both dice. Find (i) the range of $X$, (ii) the probability distribution of $X$.

\item A random variable $X$ has the probability distribution $f(x)=A\binom{3}{x}^\dagger$, $x=0,1,2,3$. (i) Determine the value of $A$. (ii) Determine $F(x)$.

\item A continuous random variable $X$ can assume values between $x=1$ and $x=3$ and has a density function given by $f(x)=\dfrac{x}{4}^\dagger$. (i) Show that the area under the curve is equal to one. (ii) Find $P(x<2)$.

\item A man draws 2 balls from a bag containing 2 white and 5 black balls. Make a probability distribution of the number of white balls.

\item A bag contains 5 white and 7 black balls. If 3 balls are drawn from the bag, what is the probability that (i) two are white and one black, (ii) all are the same colour?

\item Find $E(X)$ and $\operatorname{Var}(X)$ from the following probability distribution:
\begin{center}
\begin{tabular}{l ccccc}
\toprule
$x$ & $-2$ & $-1$ & 0 & 1 & 2 \\
$P(x)$ & $1/10$ & $2/10$ & $4/10$ & $2/10$ & $1/10$ \\
\bottomrule
\end{tabular}
\end{center}

\item Find the variance of the probability distribution:
\begin{center}
\begin{tabular}{l cccc}
\toprule
$X$ & 1 & 2 & 3 & 4 \\
$P(x)$ & 0.2 & 0.3 & 0.3 & 0.2 \\
\bottomrule
\end{tabular}
\end{center}

\item A bag contains 4 white and 2 black marbles. If two marbles are drawn at random, find the expected number of black marbles.

\item If $f(x)$ has probability density $kx^2$, $0<x<1$, determine ``$k$'' and find the probability that $\tfrac{1}{2}<X<1$.

\item Given the following probability distribution.
\begin{center}
\begin{tabular}{l ccccc}
\toprule
$X$ & 0 & 1 & 2 & 3 & 4 \\
$P(X)$ & $\dfrac{1}{126}$ & $\dfrac{20}{126}$ & $\dfrac{60}{126}$ & $\dfrac{40}{126}$ & $\dfrac{5}{126}$ \\
\bottomrule
\end{tabular}
\end{center}
Verify that $E(2x+3)=2E(x)+3$.

\item A random variable ``$x$'' has the following probability distribution.
\begin{center}
\begin{tabular}{l ccc}
\toprule
$x$ & $-2$ & 3 & 1 \\
$f(x)$ & $\dfrac{1}{3}$ & $\dfrac{1}{2}$ & $\dfrac{1}{6}$ \\
\bottomrule
\end{tabular}
\end{center}
Find $E(x)$ and $E(x^2)$.

\item A random variable $X$ has the following probability distribution.
\begin{center}
\begin{tabular}{l cccccc}
\toprule
$X$ & $-2$ & $-1$ & 0 & 1 & 2 & 3 \\
$P(x)$ & 0.1 & $k$ & 0.2 & $2k$ & 0.3 & $3k$ \\
\bottomrule
\end{tabular}
\end{center}
Find (i) $k$, (ii) $P(X<2)$.

\item Define the binomial probability distribution.

\item What is a binomial experiment? Give its properties.

\item An event has probability $p=0.3^\dagger$. Find the complete binomial distribution for $n=5$ trials.

\item Find the probability, when tossing a fair coin four times, that there will appear (i) 4 heads, (ii) 1 tail 3 heads, (iii) at least 2 heads, (iv) at most 2 heads.

\item If 20\% of the bolts produced by a machine are defective, determine the probability that out of 4 bolts chosen at random: (i) zero, (ii) 2 bolts are defective.

\item Given the probability of passing an examination is 0.75, what is the probability of (i) passing at least two examinations if you take six? (ii) failing at least 2 examinations if you take four?

\item The experience of a house agent indicates that he can provide suitable accommodation for 75\% of the clients who come to him independently. Calculate the probability that (i) less than 4 clients will get satisfactory accommodation, (ii) at least 5 clients will get satisfactory accommodation.

\item The incidence of an occupational disease in an industry is such that the workers have a 20\% chance of suffering from it. What is the probability that out of 6 workmen: (i) not more than 2 will catch the disease? (ii) 4 or more will catch the disease?

\item If 8 coins are tossed, what is the probability that (i) exactly 5 are heads, (ii) 1 to 7 are heads?

\item If $X$ is binomially distributed with $n=10$, $p=0.4$, then find the mean and variance of $Y=5X+3^\dagger$.

\item If a true die is thrown once, determine the probability of getting $0,1,2,3,4,5$ sixes. Find the mean and variance of the probability distribution so obtained.

\item Five dice are tossed 96 times. Find the expected frequencies when throwing of 4, 5 or 6 is regarded as success.

\item Four dice are thrown and the number of sixes in each throw is recorded. This is repeated 180 times. Write down the theoretical frequencies for 0, 1, 2, 3 and 4 sixes. Calculate the mean number of sixes in a single throw.

\item The mean and variance of a binomial distribution are 6 and 2.4 respectively. Find $p$ and $n$, the two parameters of the binomial distribution.

\item If the probability of a defective bolt is 0.1, find the mean and standard deviation of the distribution of defective bolts in a total of 500.

\item If in a binomial distribution the mean is 3 and the standard deviation is 1.5, find its parameters.

\item Is it possible to have a binomial distribution with mean $=5$ and S.D. $=3.7$?

\item In a binomial distribution with $n=5$, what is the value of the other parameter of the binomial, if $P(X=0)=P(X=1)$? Find the mean of the distribution.

\item Find the binomial distribution whose mean is 12 and standard deviation is 3.

\item Find the mean and variance of the binomial $(q+p)^3$.

\item Find the mean and standard deviation of the binomial distribution $(q+p)^n$.

\item Find the probability that in 5 tosses of a fair die ``3 or 4'' appears (i) at no time, (ii) once, (iii) twice, (iv) at most twice.

\item Given that $X$ is a hyper-geometric random variable with $N=8$, $n=3$, $k=5$, find $P(X\leq 3)$.

\item In a binomial distribution, Mean $=2.4$ and S.D $=1.2$. Find (i) $n$ and $p$, (ii) $P(X=2)$.

\item Four balls are drawn from a bag containing 4 white and 6 black balls. Find the probability distribution of the number of white balls.

\item Five balls are drawn from a box containing 7 red and 4 blue balls. If $X$ denotes the number of red balls drawn from the box, obtain the probability distribution of $X$. Also compute the mean and variance.

\item A fair die is rolled 5 times and the number of sixes is counted. Find the probability of (i) no six, (ii) at least one six.

\item In a hyper-geometric distribution $n=4$, $N=10$ and $k=3$, then find (i) $P(X\leq 3)$, (ii) $P(X=4)$.

\item Let $X$ be a binomial r.v with $n=5$ and $p=0.2$; find the complete binomial distribution.

\item A committee of size 4 is to be selected at random from 3 women and 4 men. Find the probability that (i) no women, (ii) at least 2 women will be selected on the committee.

\item Given that $X$ is a hyper-geometric r.v with $N=10$, $n=4$, $k=5$, compute (i) $P(X\leq 3)$, (ii) $P(X=4)$.

\item A random variable $X$ is believed to follow a binomial distribution $b(x;5,p)$. If $P(X=0)=0.03125^\dagger$, find $P(X=3)$.

\item Find the mean and standard deviation of the binomial distribution if $n=2$.

\item If $N=8$, $K=4$ and $n=3$, find the complete hyper-geometric distribution and show that the mean is 1.5.

\item A committee of size 3 is selected from 4 men and 2 women. Find the probability distribution, via the hyper-geometric experiment, for the number of men in the committee.

\item If $X$ is a binomial distribution with $n=10$, $p=0.35$, find the probability of (i) exactly 6 successes, (ii) at most 2 successes.

\item If $X$ is binomially distributed with $n=10$, $p=0.41$, then find the mean and variance of $Y=3X-2^\dagger$.

\item Three balls are drawn without replacement from a box containing 3 white and 7 red balls. If $X$ denotes the number of red balls, obtain the probability distribution of $X$.

\item Workers have a 20\% chance of suffering from an occupational disease. What is the probability that out of 6 workers: (i) exactly 2 will suffer from the disease, (ii) at least 2 will suffer from the disease?

\item Find the mean and S.D of $(q+p)^n$.

\item A bag contains 6 green and 4 red balls. Three are selected at random without replacement. Compute the probability distribution of the number of green balls.

\item \emph{(Repeat of the previous question, as given in the original.)} A bag contains 6 green and 4 red balls. Three are selected at random without replacement. Compute the probability distribution of the number of green balls.

\item Three balls are drawn from a bag containing 5 white and 3 black balls. If $X$ denotes the number of white balls, find the probability distribution of $X$.

\item The mean and variance of a binomial distribution are 2 and $1.6^\dagger$. Find (i) the parameters, (ii) $P(X=2)$.

\item If $N=4$, $n=10$ and $K=5$, find (i) $P(X=3)$, (ii) $P(X\leq 1)$.

\item Find the parameters of the binomial distribution.

\item If $N=10$, $K=6$, $n=5$, then find the mean and variance of the hyper-geometric probability distribution.

\item Assuming that each baby has probability 0.4 of being male, find the probability that a family of 4 children will have (i) at least one boy, (ii) at most one girl.

\item If $n=5$, $N=12$ and $K=3$, then find (i) $P(X\leq 1)$, (ii) $P(X>1)$.

\item In a hyper-geometric distribution $n=4$, $N=10$ and $K=3$, then find (i) $P(X\leq 2)$, (ii) $P(X=4)$.

\item If a coin is tossed 5 times and the probability that it shows 3 heads is twice the probability that it shows one tail, then find $p$ and $q$.

\item From a lot of 10 missiles, 4 are selected at random and fired. If the lot contains 3 defective missiles that will not fire, what is the probability that (i) all 4 will fire, (ii) at least 2 will not fire?

\end{longlist}

\end{document}